% Electromagnetic Induction — Physics Upper Sixth — Cours signé OnBuch+
% Official MINESEC syllabus. Compiler avec : tectonic PHY08-electromagnetic-induction.tex
\newcommand{\DOCMATIERE}{Physics}
\newcommand{\DOCNIVEAU}{Upper Sixth}
\newcommand{\DOCMODULE}{Upper Sixth — Physics}
\newcommand{\DOCLECON}{8}
\newcommand{\DOCTITRE}{Electromagnetic Induction}
% OnBuch+ — préambule « cours signés » ENRICHI (compilé avec tectonic / XeLaTeX)
% Système visuel « Pop » d'OnBuch+ : fond crème (jamais blanc), bordures encre
% épaisses, ombres « dures » décalées, titres Archivo Black en capitales.
% Version MATHÉMATIQUES : graphes (pgfplots/tikz), boîtes pédagogiques
% (définition, propriété, méthode, attention, le savais-tu, exercice résolu,
% exercice, TP, bilan), page de garde et signature OnBuch+ ancrée en bas.
%
% Macros attendues dans main.tex AVANT \input{preamble} :
%   \DOCMATIERE  \DOCNIVEAU  \DOCMODULE  \DOCLECON (numéro)  \DOCTITRE
\documentclass[11pt]{article}

\usepackage{fontspec}
\usepackage[english]{babel}
\usepackage[a4paper,margin=2cm,top=3.4cm,bottom=2.8cm,headheight=1.4cm,footskip=1.3cm]{geometry}
\usepackage{amsmath,amssymb,amsfonts}
\usepackage{mathtools} % \xmapsto, \coloneqq... souvent utilisés par les modèles
\usepackage{cancel} % simplifications barrées (bilans de forces)
% \abs{z} (module, valeur absolue) et \norm{u} (norme), utilisés par les modèles
\providecommand{\abs}{}\let\abs\relax\DeclarePairedDelimiter{\abs}{\lvert}{\rvert}
\providecommand{\norm}{}\let\norm\relax\DeclarePairedDelimiter{\norm}{\lVert}{\rVert}
\usepackage[table]{xcolor} % [table] : fournit aussi \rowcolors (alternance de lignes)
\usepackage{pagecolor}
\usepackage{tikz}
\usetikzlibrary{arrows.meta,positioning,calc,shapes.geometric,shapes.misc,shapes.arrows,decorations.pathmorphing,decorations.pathreplacing,decorations.markings,patterns,shadows,mindmap,trees,fit,backgrounds,matrix,intersections,angles,quotes}
\usepackage{pgfplots}
\usepackage[version=4]{mhchem} % équations nucléaires \ce{^{235}_{92}U}
\usepackage[european,siunitx]{circuitikz} % schémas électriques
\pgfplotsset{compat=1.18}
\usepgfplotslibrary{fillbetween,groupplots} % aires entre courbes (calcul intégral)
\usepackage{enumitem}
\usepackage{fancyhdr}
% [most] charge toutes les bibliothèques tcolorbox (listings, minted, raster,
% documentation...) et gonfle inutilement la mémoire TeX sur un document long
% (~300 boîtes) : on ne charge que ce qui est réellement utilisé.
\usepackage[skins,breakable]{tcolorbox}
\usepackage{graphicx}
\usepackage{colortbl}
\usepackage{booktabs}
\usepackage{tabularx}
\usepackage{diagbox} % case d'en-tête diagonale des tableaux à double entrée
% Colonnes extensibles centrée / gauche / droite (C, L, R), souvent utilisées par les modèles
\newcolumntype{C}{>{\centering\arraybackslash}X}
\newcolumntype{L}{>{\raggedright\arraybackslash}X}
\newcolumntype{R}{>{\raggedleft\arraybackslash}X}
\usepackage{array}
\usepackage{multicol}
\usepackage{siunitx}
% parse-numbers=false : accepte aussi \SI{2,5 \times 10^{-3}}{...}
\sisetup{locale=UK,output-decimal-marker={.},group-minimum-digits=4,parse-numbers=false}
\DeclareSIUnit\ohm{\ensuremath{\symup{\Omega}}} % U+2126 absent de la police
\DeclareSIUnit\division{div} % réglages d'oscilloscope (ms/div, V/div)
\DeclareSIUnit\tr{tr} % tours (vitesse de rotation, ex. \SI{1500}{\tr\per\minute})
\usepackage{qrcode}
% \degree : commande fréquemment utilisée par les modèles pour un angle ou une
% température en dehors de \SI{}{} (non fournie par les paquets chargés).
\providecommand{\degree}{^{\circ}}
% \qed (fin de démonstration), utilisé par les modèles sans amsthm
\providecommand{\qed}{\hfill$\square$}
% \barre : double barre des valeurs interdites dans un tableau de variations (tkz-tab)
\providecommand{\barre}{\ensuremath{\|}}
% environnement proof (amsthm non chargé) utilisé par les modèles
\ifdefined\proof\else\newenvironment{proof}[1][Proof]{\par\noindent\textit{#1.}\ }{\qed\par}\fi
\usepackage{lastpage}
\usepackage{hyperref}
\hypersetup{hidelinks}

% ============================================================
% Palette « Pop » OnBuch+ — mêmes hex que la classe P (plus_ui.dart)
% ============================================================
\definecolor{poporange}{HTML}{F59321}
\definecolor{popdark}{HTML}{E07A0C}
\definecolor{popcream}{HTML}{FFF6E8}
\definecolor{popink}{HTML}{1C1714}
\definecolor{poppurple}{HTML}{7A5AE0}
\definecolor{popgreen}{HTML}{2E9E6A}
\definecolor{poppink}{HTML}{E0457B}
\definecolor{popblue}{HTML}{2F7DE1}
\definecolor{popgold}{HTML}{C98A12}
\definecolor{popmuted}{HTML}{8A7F76}
\definecolor{popline}{HTML}{EADFD2}
\definecolor{popgray}{HTML}{8A7F76}
\colorlet{popgrey}{popgray}
% Alias tolérants (le modèle nomme parfois les couleurs différemment)
\colorlet{popwhite}{white}
\colorlet{popblack}{popink}
\colorlet{popred}{poppink}
\colorlet{popyellow}{popgold}
% Teintes claires pour fonds de tableaux / zones
\colorlet{popblueL}{popblue!10!white}
\colorlet{poporangeL}{poporange!14!white}
\colorlet{popgreenL}{popgreen!12!white}
\colorlet{poppurpleL}{poppurple!12!white}
\colorlet{poppinkL}{poppink!12!white}
\colorlet{popgoldL}{popgold!14!white}

\pagecolor{popcream}

% ============================================================
% Typographie
% ============================================================
\defaultfontfeatures{Ligatures=TeX}
\setmainfont{PlusJakartaSans-Medium.ttf}[Path=../../../fonts/,BoldFont=PlusJakartaSans-Bold.ttf]
% Maths en sans-serif (Fira Math) avec chiffres et lettres droites de Plus
% Jakarta, pour que les formules se fondent dans le texte.
\usepackage[math-style=ISO,bold-style=ISO]{unicode-math}
\setmathfont{FiraMath-Regular.otf}[Path=../../../fonts/]
\setmathfont{PlusJakartaSans-Medium.ttf}[Path=../../../fonts/,range={up/num,up/{Latin,latin},\mathup}]
% (icomma retiré : décimales avec point en anglais)
% Caractères mathématiques Unicode absents des polices de texte (Plus Jakarta,
% Archivo Black) : rendus par leur équivalent mathématique, en texte comme en
% formule — sinon ils s'affichent en carré vide (ex. « Arithmétique dans ℤ »).
\usepackage{newunicodechar}
\newunicodechar{ℝ}{\ensuremath{\mathbb{R}}}
\newunicodechar{ℤ}{\ensuremath{\mathbb{Z}}}
\newunicodechar{ℂ}{\ensuremath{\mathbb{C}}}
\newunicodechar{ℕ}{\ensuremath{\mathbb{N}}}
\newunicodechar{ℚ}{\ensuremath{\mathbb{Q}}}
\newunicodechar{𝔻}{\ensuremath{\mathbb{D}}}
\newunicodechar{α}{\ensuremath{\alpha}}
\newunicodechar{β}{\ensuremath{\beta}}
\newunicodechar{γ}{\ensuremath{\gamma}}
\newunicodechar{δ}{\ensuremath{\delta}}
\newunicodechar{ε}{\ensuremath{\varepsilon}}
\newunicodechar{θ}{\ensuremath{\theta}}
\newunicodechar{λ}{\ensuremath{\lambda}}
\newunicodechar{μ}{\ensuremath{\mu}}
\newunicodechar{σ}{\ensuremath{\sigma}}
\newunicodechar{τ}{\ensuremath{\tau}}
\newunicodechar{φ}{\ensuremath{\varphi}}
\newunicodechar{ψ}{\ensuremath{\psi}}
\newunicodechar{ω}{\ensuremath{\omega}}
\newunicodechar{Σ}{\ensuremath{\Sigma}}
% majuscules produites par \MakeUppercase dans les titres (α-aminés -> Α)
\newunicodechar{Α}{\ensuremath{\alpha}}
\newunicodechar{Β}{\ensuremath{\beta}}
\newunicodechar{Γ}{\ensuremath{\Gamma}}
\newunicodechar{Δ}{\ensuremath{\Delta}}
\newunicodechar{Ε}{\ensuremath{\varepsilon}}
\newunicodechar{Θ}{\ensuremath{\Theta}}
\newunicodechar{Λ}{\ensuremath{\Lambda}}
\newunicodechar{Μ}{\ensuremath{\mu}}
\newunicodechar{Τ}{\ensuremath{\tau}}
\newunicodechar{Φ}{\ensuremath{\Phi}}
\newunicodechar{Ψ}{\ensuremath{\Psi}}
\newunicodechar{Ω}{\ensuremath{\Omega}}
\newunicodechar{↦}{\ensuremath{\mapsto}}
\newunicodechar{∈}{\ensuremath{\in}}
\newunicodechar{∉}{\ensuremath{\notin}}
\newunicodechar{∧}{\ensuremath{\wedge}}
\newunicodechar{⇔}{\ensuremath{\Leftrightarrow}}
\newunicodechar{⟺}{\ensuremath{\Longleftrightarrow}}
\newunicodechar{⇒}{\ensuremath{\Rightarrow}}
\newunicodechar{⟹}{\ensuremath{\Longrightarrow}}
\newunicodechar{∘}{\ensuremath{\circ}}
\newunicodechar{∪}{\ensuremath{\cup}}
\newunicodechar{∩}{\ensuremath{\cap}}
\newunicodechar{≡}{\ensuremath{\equiv}}
\newunicodechar{⊂}{\ensuremath{\subset}}
\newunicodechar{⊥}{\ensuremath{\perp}}
\newunicodechar{∃}{\ensuremath{\exists}}
\newunicodechar{∀}{\ensuremath{\forall}}
\newunicodechar{∛}{\ensuremath{\sqrt[3]{\,}}}
\newunicodechar{∜}{\ensuremath{\sqrt[4]{\,}}}
\newunicodechar{⃗}{\ensuremath{{}^{\to}}}
\newunicodechar{ⁿ}{\ifmmode^{n}\else\textsuperscript{\ensuremath{n}}\fi}
\newunicodechar{ˣ}{\ifmmode^{x}\else\textsuperscript{\ensuremath{x}}\fi}
\newunicodechar{ᵇ}{\ifmmode^{b}\else\textsuperscript{\ensuremath{b}}\fi}
\newunicodechar{ʳ}{\ifmmode^{r}\else\textsuperscript{\ensuremath{r}}\fi}
\newunicodechar{ᵘ}{\ifmmode^{u}\else\textsuperscript{\ensuremath{u}}\fi}
\newunicodechar{ʸ}{\ifmmode^{y}\else\textsuperscript{\ensuremath{y}}\fi}
\newunicodechar{ᵃ}{\ifmmode^{a}\else\textsuperscript{\ensuremath{a}}\fi}
\newunicodechar{ᵉ}{\ifmmode^{e}\else\textsuperscript{\ensuremath{e}}\fi}
\newunicodechar{ᵏ}{\ifmmode^{k}\else\textsuperscript{\ensuremath{k}}\fi}
\newunicodechar{ᵖ}{\ifmmode^{p}\else\textsuperscript{\ensuremath{p}}\fi}
\newunicodechar{ᴺ}{\ifmmode^{N}\else\textsuperscript{\ensuremath{N}}\fi}
\newunicodechar{ᶜ}{\ifmmode^{c}\else\textsuperscript{\ensuremath{c}}\fi}
\newunicodechar{ʰ}{\ifmmode^{h}\else\textsuperscript{\ensuremath{h}}\fi}
\newunicodechar{ᵠ}{\ifmmode^{q}\else\textsuperscript{\ensuremath{q}}\fi}
\newunicodechar{ₙ}{\ifmmode_{n}\else\textsubscript{\ensuremath{n}}\fi}
\newunicodechar{ₐ}{\ifmmode_{a}\else\textsubscript{\ensuremath{a}}\fi}
\newunicodechar{ᵢ}{\ifmmode_{i}\else\textsubscript{\ensuremath{i}}\fi}
\newunicodechar{ᵣ}{\ifmmode_{r}\else\textsubscript{\ensuremath{r}}\fi}
\newunicodechar{ₖ}{\ifmmode_{k}\else\textsubscript{\ensuremath{k}}\fi}
\newunicodechar{ᵦ}{\ifmmode_{\beta}\else\textsubscript{\ensuremath{\beta}}\fi}
\newunicodechar{ₓ}{\ifmmode_{x}\else\textsubscript{\ensuremath{x}}\fi}
\newfontfamily\popdisplay{ArchivoBlack-Regular.ttf}[Path=../../../fonts/]
\newfontfamily\poplogo{SpaceGrotesk-Bold.ttf}[Path=../../../fonts/]
\newfontfamily\popsemi{PlusJakartaSans-SemiBold.ttf}[Path=../../../fonts/]
\newfontfamily\popxbold{PlusJakartaSans-ExtraBold.ttf}[Path=../../../fonts/]
\newfontfamily\popxboldls{PlusJakartaSans-ExtraBold.ttf}[Path=../../../fonts/,LetterSpace=3.0]

\setlist[enumerate]{leftmargin=1.7em,itemsep=5pt,topsep=4pt}
\setlist[itemize]{leftmargin=1.7em,itemsep=4pt,topsep=4pt,label={\color{poporange}\textbullet}}
\setlength{\parindent}{0pt}
\setlength{\parskip}{5pt}
\renewcommand{\arraystretch}{1.25}

% pgfplots : style Pop par défaut
\pgfplotsset{
  every axis/.append style={axis line style={popink,line width=1pt},tick style={popink},
    grid style={popline,line width=0.5pt},label style={font=\small},tick label style={font=\footnotesize}},
  popcurve/.style={poporange,line width=1.8pt,no markers,smooth},
}
% Même style disponible dans un \draw TikZ simple (hors pgfplots)
\tikzset{popcurve/.style={poporange,line width=1.8pt}}
\tikzset{popgrid/.style={popline,very thin}}
\colorlet{popcurve}{poporange} % le nom du style est parfois utilisé comme couleur

% ============================================================
% Pill / squiggle
% ============================================================
\newcommand{\poppill}[3]{%
  \tikz[baseline=-0.55ex]\node[draw=popink,line width=1.2pt,rounded corners=2.6pt,%
    fill=#2,inner xsep=6.5pt,inner ysep=3pt]{%
    \popxboldls\fontsize{7}{7}\selectfont\color{#3}\MakeUppercase{#1}};%
}
\newcommand{\popsquiggle}[1]{%
  \tikz[baseline=-0.4ex]{
    \draw[#1,line width=2.6pt,line cap=round]
      plot [smooth] coordinates {(0,0.09) (0.34,0.01) (0.68,0.21) (1.02,0.01) (1.36,0.10)};
  }%
}
% Mot-clé surligné (marqueur orange sous le texte)
\newcommand{\cle}[1]{{\popxbold\color{popdark}#1}}

% ============================================================
% En-tête / pied
% ============================================================
\pagestyle{fancy}
\fancyhf{}
\renewcommand{\headrulewidth}{0pt}
\renewcommand{\footrulewidth}{0pt}
\lhead{\raisebox{-0.15ex}{\poplogo\fontsize{13}{13}\selectfont\color{popink}OnBuch\color{poporange}+}}
\rhead{\poppill{\DOCMATIERE\ \textbullet\ \DOCNIVEAU\ \textbullet\ Chapter \DOCLECON}{white}{popink}}
\lfoot{\popxbold\fontsize{7.6}{7.6}\selectfont\color{popmuted}\MakeUppercase{Signed course OnBuch+}\ \textbullet\ {\popsemi\DOCTITRE}}
\rfoot{\tikz[baseline=-0.6ex]\node[rounded corners=8.5pt,fill=popink,minimum height=17pt,minimum width=17pt,inner xsep=5pt,inner sep=0]{\popxbold\fontsize{8}{8}\selectfont\color{popcream}\thepage\,/\,\pageref*{LastPage}};}

% ============================================================
% Titres
% ============================================================
\newcommand{\chaptitre}[2]{%
  \begin{tcolorbox}[enhanced,colback=poporange,colframe=popink,boxrule=2.6pt,arc=11pt,
      drop shadow=popink,left=15pt,right=15pt,top=12pt,bottom=11pt,before skip=3pt,after skip=18pt]
    \poppill{#2}{popink}{popcream}\par\vspace{9pt}
    {\raggedright\hyphenpenalty=10000\exhyphenpenalty=10000\popdisplay\fontsize{22}{24}\selectfont\color{popcream}\MakeUppercase{#1}\par}
  \end{tcolorbox}
}
% Section numérotée : pastille orange avec le numéro + titre Archivo
\newcounter{popsec}
\newcommand{\coursec}[1]{%
  \stepcounter{popsec}\setcounter{popsub}{0}%
  \par\vspace{16pt}\noindent
  \begin{minipage}{\linewidth}
  \tikz[baseline=-0.7ex]\node[draw=popink,line width=1.6pt,fill=poporange,rounded corners=4pt,minimum size=22pt,inner sep=0]{\popdisplay\fontsize{12}{12}\selectfont\color{popcream}\thepopsec};%
  \hspace{8pt}{\popdisplay\fontsize{15}{17}\selectfont\color{popink}\MakeUppercase{#1}}\par
  \vspace{4pt}\noindent{\color{poporange}\rule{36pt}{3.6pt}}
  \end{minipage}\par\nopagebreak\vspace{8pt}
}
\newcounter{popsub}
\newcommand{\courssub}[1]{%
  \stepcounter{popsub}%
  \par\vspace{9pt}\noindent{\popxbold\fontsize{11.5}{13}\selectfont\color{popink}{\color{poporange}\thepopsec.\thepopsub}\hspace{6pt}#1}\par\nopagebreak\vspace{4pt}
}

% ============================================================
% Boîtes pédagogiques — même recette : fond blanc, bordure épaisse colorée,
% ombre dure encre, badge plein. Titre optionnel [#1].
% ============================================================
\tcbset{popbase/.style={breakable,enhanced,colback=white,boxrule=2.3pt,arc=9pt,
  shadow={3pt}{-3pt}{0pt}{popink},left=14pt,right=14pt,top=11pt,bottom=11pt,before skip=11pt,after skip=15pt,
  fontupper=\popsemi\fontsize{10.6}{14.5}\selectfont\color{popink},
  fontlower=\popsemi\fontsize{10.6}{14.5}\selectfont\color{popink}}}
% Coche dessinée (le glyphe ✓ n'existe pas dans les polices du document)
\newcommand{\popcheck}[1]{\tikz[baseline=-0.5ex]\draw[#1,line width=1.6pt,line cap=round,line join=round](-0.13,0.0)--(-0.04,-0.09)--(0.13,0.1);}
\providecommand{\coche}{\popcheck{popgreen}}
\providecommand{\resultat}[1]{\par\smallskip\noindent\fcolorbox{popgreen}{popgreenL}{\parbox{\dimexpr\linewidth-2\fboxsep-2\fboxrule\relax}{\textbf{Result.} #1}}\par\smallskip}
\newcommand{\popboxhead}[3]{\poppill{#1}{#2}{white}\ifx\relax#3\relax\else\hspace{6pt}{\popxbold\fontsize{10.6}{12}\selectfont\color{popink}#3}\fi\par\vspace{7pt}}

\newtcolorbox{aretenir}[1][]{popbase,colframe=popgreen,before upper={\popboxhead{Key points}{popgreen}{#1}}}
\newtcolorbox{exemplebox}[1][]{popbase,colframe=poporange,before upper={\popboxhead{Exemple}{poporange}{#1}}}
\newtcolorbox{definition}[1][]{popbase,colframe=popblue,before upper={\popboxhead{Definition}{popblue}{#1}}}
\newtcolorbox{propriete}[1][]{popbase,colframe=poppurple,before upper={\popboxhead{Property}{poppurple}{#1}}}
\newtcolorbox{methode}[1][]{popbase,colframe=popgold,before upper={\popboxhead{Method}{popgold}{#1}}}
\newtcolorbox{attention}[1][]{popbase,colframe=poppink,before upper={\popboxhead{Warning}{poppink}{#1}}}
\newtcolorbox{experience}[1][]{popbase,colframe=popblue,colback=popblueL,before upper={\popboxhead{Experiment}{popblue}{#1}}}
% « Le savais-tu ? » : carte violette pleine (contexte, culture, Cameroun)
\newtcolorbox{savaistu}[1][]{popbase,colback=poppurple,colframe=popink,
  fontupper=\popsemi\fontsize{10.4}{14.2}\selectfont\color{white},
  before upper={\poppill{Did you know?}{white}{poppurple}\ifx\relax#1\relax\else\hspace{6pt}{\popxbold\color{white}#1}\fi\par\vspace{7pt}}}
% Exercice résolu : énoncé en haut, \tcblower puis la solution sur fond crème
\newcounter{popexo}\newcounter{popexr}
\newtcolorbox{exoresolu}[1][]{popbase,colframe=poporange,colbacklower=popcream,
  segmentation style={popink,line width=1.2pt,dashed},
  before upper={\stepcounter{popexr}\popboxhead{Worked example \thepopexr}{poporange}{#1}},
  before lower={\poppill{Solution}{popink}{popcream}\par\vspace{6pt}}}
% Exercice (non résolu en place) — la correction est dans la partie Corrigés
\newtcolorbox{exercice}[1][]{popbase,colframe=popink,
  before upper={\stepcounter{popexo}\popboxhead{Exercise \thepopexo}{popink}{#1}}}
% Corrigé
\newtcolorbox{corrige}[1][]{popbase,colframe=popgreen,colback=popgreenL,
  before upper={\popboxhead{Solution}{popgreen}{#1}}}
% Objectifs / prérequis / bilan
\newtcolorbox{objectifs}{popbase,colframe=popink,colback=poporangeL,
  before upper={\popboxhead{Chapter objectives}{popink}{}}}
\newtcolorbox{prerequis}{popbase,colframe=popink,colback=popblueL,
  before upper={\popboxhead{Prerequisites}{popblue}{}}}
\newtcolorbox{bilan}{popbase,colframe=popink,colback=white,boxrule=2.8pt,
  before upper={\popboxhead{Fiche bilan}{poporange}{L'essentiel à connaître pour l'examen}}}
% Figure encadrée avec légende
\newtcolorbox{popfigure}[1][]{enhanced,breakable=false,colback=white,colframe=popline,boxrule=1.4pt,arc=8pt,
  left=10pt,right=10pt,top=10pt,bottom=8pt,before skip=10pt,after skip=12pt,halign=center,
  lower separated=false,halign lower=center,
  fontlower=\popsemi\fontsize{9}{11.5}\selectfont\color{popmuted},
  #1}
\newcommand{\legende}[1]{\par\vspace{6pt}{\popsemi\fontsize{9}{11.5}\selectfont\color{popmuted}#1}}

% ============================================================
% Signature OnBuch+
% ============================================================
\newcommand{\poplogoinline}[1]{{\poplogo\fontsize{#1}{#1}\selectfont\color{popink}OnBuch\color{poporange}+}}
\newcommand{\signatureonbuch}{%
  \begin{tcolorbox}[enhanced,colback=popink,colframe=popink,boxrule=2.2pt,arc=10pt,
      drop shadow=poporange,left=14pt,right=14pt,top=10pt,bottom=10pt,before skip=0pt,after skip=0pt]
    \tikz[baseline=-0.7ex]\node[circle,fill=popgreen,draw=popcream,line width=1.3pt,minimum size=22pt,inner sep=0]{\popcheck{white}};%
    \hspace{9pt}\begin{minipage}[c]{0.62\linewidth}
      {\popxbold\fontsize{12}{13}\selectfont\color{popcream}Signed\ \textbullet\ {\poplogo OnBuch\color{poporange}+}}\par\vspace{2pt}
      {\raggedright\popsemi\fontsize{8}{10.5}\selectfont\color{popline}\MakeUppercase{OnBuch+ teaching team}\\\MakeUppercase{Follows the official MINESEC syllabus}\par}
    \end{minipage}\hfill
    {\poplogo\fontsize{22}{22}\selectfont\color{popcream}OnBuch\color{poporange}+}
  \end{tcolorbox}%
}

% ============================================================
% Page de garde
% #1 titre · #2 matière · #3 niveau · #4 module · #5 n° leçon · #6 durée ·
% #7 description / objectifs (texte ou itemize)
% ============================================================
\newcommand{\pagegarde}[7]{%
  \thispagestyle{empty}
  \newgeometry{a4paper,margin=1.7cm,top=1.6cm,bottom=1.4cm}%
  \begin{tikzpicture}[remember picture,overlay]
    \fill[poporange!18] ([xshift=-3.2cm,yshift=-3.6cm]current page.north east) circle (4.6cm);
    \fill[poppurple!14] ([xshift=2.4cm,yshift=7.6cm]current page.south west) circle (3.8cm);
  \end{tikzpicture}%
  {\poplogo\fontsize{30}{30}\selectfont\color{popink}OnBuch\color{poporange}+}\hfill
  \raisebox{0.6ex}{\poppill{Signed course \textbullet\ Premium}{popink}{popcream}}\par
  \vspace{1pt}\popsquiggle{poppurple}\par
  \vspace{8pt}
  \begin{tcolorbox}[enhanced,colback=poporange,colframe=popink,boxrule=2.8pt,arc=13pt,
      drop shadow=popink,left=18pt,right=18pt,top=12pt,bottom=13pt,after skip=10pt]
    \poppill{#2 \textbullet\ #3}{popink}{popcream}\hspace{5pt}\poppill{#4}{white}{popink}\par\vspace{8pt}
    {\popdisplay\fontsize{14}{14}\selectfont\color{popink}CHAPTER #5}\par\vspace{4pt}
    {\raggedright\hyphenpenalty=10000\exhyphenpenalty=10000\popdisplay\fontsize{25}{27}\selectfont\color{popcream}\MakeUppercase{#1}\par}
    \vspace{8pt}
    {\popxbold\fontsize{9.5}{9.5}\selectfont\color{popink}\MakeUppercase{Suggested time: #6}}
  \end{tcolorbox}
  \begin{tcolorbox}[enhanced,colback=white,colframe=popink,boxrule=2.2pt,arc=10pt,
      drop shadow=popink,left=16pt,right=16pt,top=10pt,bottom=10pt,after skip=8pt]
    \poppill{In this chapter}{poppurple}{white}\par\vspace{5pt}
    \popsemi\fontsize{9.3}{12.6}\selectfont\color{popink}#7
  \end{tcolorbox}
  \noindent\begin{minipage}[t]{0.487\linewidth}
    \begin{tcolorbox}[enhanced,colback=popgreen,colframe=popink,boxrule=2.2pt,arc=10pt,
        drop shadow=popink,left=11pt,right=11pt,top=10pt,bottom=10pt,height=3.1cm,valign=center]
      \tcbox[colback=white,colframe=popink,boxrule=1.3pt,arc=5pt,boxsep=0pt,nobeforeafter,
        left=3pt,right=3pt,top=3pt,bottom=3pt,valign=center]{\qrcode[height=1.9cm]{https://play.google.com/store/apps/details?id=com.luvvix.onbuch&hl=en}}%
      \hspace{8pt}\begin{minipage}[c]{0.5\linewidth}
        {\popdisplay\fontsize{11}{12.5}\selectfont\color{white}\MakeUppercase{Download OnBuch}}\par\vspace{4pt}
        {\raggedright\popsemi\fontsize{8.4}{11}\selectfont\color{white}Free \textbullet\ Google Play\\AI tutor, past papers, results\par}
      \end{minipage}
    \end{tcolorbox}
  \end{minipage}\hfill
  \begin{minipage}[t]{0.487\linewidth}
    \begin{tcolorbox}[enhanced,colback=poppurple,colframe=popink,boxrule=2.2pt,arc=10pt,
        drop shadow=popink,left=11pt,right=11pt,top=10pt,bottom=10pt,height=3.1cm,valign=center]
      \tikz[baseline=-0.7ex]\node[circle,fill=white,draw=popink,line width=1.3pt,minimum size=1.6cm,inner sep=0]{\popdisplay\fontsize{24}{24}\selectfont\color{poppurple}+};%
      \hspace{8pt}\begin{minipage}[c]{0.6\linewidth}
        {\popdisplay\fontsize{11}{12.5}\selectfont\color{white}\MakeUppercase{Go OnBuch+}}\par\vspace{4pt}
        {\raggedright\popsemi\fontsize{8.4}{11}\selectfont\color{white}All signed courses, unlimited AI tutor, PDF sheets — OnBuch+ tab in the app\par}
      \end{minipage}
    \end{tcolorbox}
  \end{minipage}
  \vspace{8pt}
  \popcontenu
  \vfill
  \signatureonbuch
  \restoregeometry
  \newpage
}

% Rangée « Ce cours contient » (page de garde)
\newcommand{\poptile}[3]{% #1 couleur #2 gros texte #3 légende
  \begin{tcolorbox}[enhanced,colback=#1,colframe=popink,boxrule=2pt,arc=9pt,shadow={2.5pt}{-2.5pt}{0pt}{popink},
      left=6pt,right=6pt,top=9pt,bottom=9pt,halign=center,valign=center,height=1.75cm,width=\linewidth,nobeforeafter]
    {\popdisplay\fontsize{12.5}{13}\selectfont\color{white}\MakeUppercase{#2}}\par\vspace{3pt}
    {\centering\popsemi\fontsize{7.8}{9.5}\selectfont\color{white}#3\par}
  \end{tcolorbox}}
\newcommand{\popcontenu}{%
  \par\noindent{\popxboldls\fontsize{7.5}{7.5}\selectfont\color{popmuted}THIS SIGNED COURSE CONTAINS}\par\vspace{7pt}
  \noindent\begin{minipage}[t]{0.235\linewidth}\poptile{popblue}{Course}{complete, illustrated, step by step}\end{minipage}\hfill
  \begin{minipage}[t]{0.235\linewidth}\poptile{poporange}{Methods}{and worked examples}\end{minipage}\hfill
  \begin{minipage}[t]{0.235\linewidth}\poptile{poppink}{Exercises}{exam style + solutions}\end{minipage}\hfill
  \begin{minipage}[t]{0.235\linewidth}\poptile{popgreen}{Summary sheet}{the essentials to remember}\end{minipage}\par}

% Sommaire
\newtcolorbox{sommaire}{enhanced,colback=white,colframe=popink,boxrule=2.2pt,arc=10pt,
  drop shadow=popink,left=18pt,right=18pt,top=13pt,bottom=13pt,before skip=6pt,after skip=20pt,
  before upper={\poppill{Contents}{poppurple}{white}\par\vspace{9pt}\popsemi\fontsize{10.6}{14.5}\selectfont\color{popink}}}

% Signature de fin de cours — poussée tout en bas de la dernière page
\newcommand{\findecours}{%
  \par\vfill\nopagebreak
  \begin{center}\popsquiggle{poporange}\end{center}\vspace{4pt}
  \signatureonbuch
}
\AtBeginDocument{\def\llbracket{\mathopen{[\mkern-3.5mu[}}\def\rrbracket{\mathclose{]\mkern-3.5mu]}}} % intervalles d'entiers (glyphe ⟦ absent de la police)
% Glyphes absents de Fira Math : redéfinis à partir de glyphes disponibles
\AtBeginDocument{%
  \def\setminus{\mathbin{\backslash}}\let\smallsetminus\setminus
  \def\vdots{\mathord{\vbox{\baselineskip3.2pt\lineskiplimit0pt\kern4pt\hbox{.}\hbox{.}\hbox{.}}}}%
  \def\ddots{\mathinner{\mkern1mu\raise6.5pt\hbox{.}\mkern2mu\raise3.6pt\hbox{.}\mkern2mu\raise0.7pt\hbox{.}\mkern1mu}}%
  \def\triangle{\mathord{\scalebox{0.9}{$\symup{\Delta}$}}}%
  \def\nabla{\mathord{\raisebox{\depth}{\scalebox{1}[-1]{$\symup{\Delta}$}}}}%
  \def\checkmark{\popcheck{popdark}}%
  \def\star{\mathbin{\ast}}\def\bigstar{\mathord{\ast}}%
  \def\diamond{\mathbin{\scalebox{0.6}{$\Diamond$}}}%
  \def\bigcup{\mathop{\mathchoice{\raisebox{-0.3em}{\scalebox{1.7}{$\cup$}}}{\scalebox{1.25}{$\cup$}}{\cup}{\cup}}\displaylimits}%
  \def\bigcap{\mathop{\mathchoice{\raisebox{-0.3em}{\scalebox{1.7}{$\cap$}}}{\scalebox{1.25}{$\cap$}}{\cap}{\cap}}\displaylimits}%
  \def\lesseqqgtr{\mathrel{\lessgtr}}\def\bigoplus{\mathop{\oplus}\displaylimits}%
}
\providecommand{\ssi}{if and only if} % abréviation fréquente
\AtBeginDocument{\providecommand{\wideparen}{\overparen}} % arc de cercle

%% FIGFIX-BEGIN v5 — correctifs globaux des figures (docs/onbuch-generation/tools/figfix)
\usepackage{adjustbox}\usepackage{etoolbox}\tcbuselibrary{raster}
% toute figure TikZ/pgfplots plus large que la ligne est réduite à la largeur disponible
\BeforeBeginEnvironment{tikzpicture}{\begin{adjustbox}{max width=\linewidth}}
\AfterEndEnvironment{tikzpicture}{\end{adjustbox}}
% légendes pgfplots lisibles par-dessus les courbes
\pgfplotsset{every axis/.append style={legend style={fill=white,fill opacity=0.92,text opacity=1,draw=black!15,inner sep=3pt}}}
% étiquettes d'arêtes (syntaxe edge["..."]) avec fond blanc
\tikzset{every edge quotes/.append style={fill=white,inner sep=1.2pt}}
%% FIGFIX-END
\begin{document}

\pagegarde{Electromagnetic Induction}{Physics}{Upper Sixth}{Upper Sixth — Physics}{8}{2 periods}{This chapter introduces electromagnetic induction: how a changing magnetic flux produces an induced e.m.f. (Faraday's and Lenz's laws), and how this principle powers AC and DC generators. It then develops self- and mutual inductance, the growth and decay of current in an L-R circuit, and the energy stored in an inductor.
\vspace{4pt}
\begin{multicols}{2}\small
\begin{itemize}
\item By the end of this chapter I can define magnetic flux and calculate the flux through a coil of N turns in a uniform field.
\item By the end of this chapter I can state and apply Faraday's law e = -N dΦ/dt to compute an induced e.m.f.
\item By the end of this chapter I can use Lenz's law (and Fleming's right-hand rule) to predict the direction of an induced current.
\item By the end of this chapter I can derive and apply e = Blv for a conductor moving through a uniform magnetic field.
\item By the end of this chapter I can explain the working principles of simple AC and DC generators and sketch their output e.m.f. against time.
\item By the end of this chapter I can define self-inductance L and mutual inductance M through NΦ = LI and NΦ₂ = MI₁, and compute them in simple cases.
\end{itemize}\end{multicols}}

\begin{sommaire}
\begin{multicols}{2}
\begin{enumerate}
\item Magnetic Flux and Faraday's Law of Induction
\item Lenz's Law and the Motional E.M.F. (e = Blv)
\item Simple AC and DC Generators
\item Self- and Mutual Inductance; the L-R Circuit and Stored Energy
\item Integration activity
\item Methods and worked examples
\item Exercises
\item Solutions to the exercises
\item Summary sheet
\end{enumerate}
\end{multicols}
\end{sommaire}

\begin{objectifs}
\begin{itemize}
\item By the end of this chapter I can define magnetic flux and calculate the flux through a coil of N turns in a uniform field.
\item By the end of this chapter I can state and apply Faraday's law e = -N dΦ/dt to compute an induced e.m.f.
\item By the end of this chapter I can use Lenz's law (and Fleming's right-hand rule) to predict the direction of an induced current.
\item By the end of this chapter I can derive and apply e = Blv for a conductor moving through a uniform magnetic field.
\item By the end of this chapter I can explain the working principles of simple AC and DC generators and sketch their output e.m.f. against time.
\item By the end of this chapter I can define self-inductance L and mutual inductance M through NΦ = LI and NΦ₂ = MI₁, and compute them in simple cases.
\item By the end of this chapter I can write and solve the differential equation of an L-R circuit for the growth and decay of current.
\item By the end of this chapter I can calculate the energy stored in an inductor, E = ½LI².
\item By the end of this chapter I can interpret time constants and exponential curves in induction problems, as required in GCE structured questions.
\end{itemize}
\end{objectifs}

\begin{prerequis}
\begin{itemize}
\item Magnetic field B, force on a current-carrying conductor F = BIl, and Fleming's left-hand rule (Lower Sixth, magnetic fields chapter)
\item Magnetic flux Φ = BA cos θ and flux linkage NΦ
\item Electric current, e.m.f., resistance and Ohm's law; Kirchhoff's laws
\item Exponential functions and simple first-order differential equations (Further/Pure Mathematics)
\item Alternating current basics: sinusoidal variation, peak and r.m.s. values (Lower Sixth)
\end{itemize}
\end{prerequis}

\newpage

\coursec{Magnetic Flux and Faraday's Law of Induction}

During the long rainy-season blackouts in Douala and Bamenda, many families switch on a small petrol generator --- yet few realise that this machine works by spinning a coil in a magnetic field, exactly like the alternator of a motorcycle. Even more surprising: you can light a bulb by simply waving a magnet near a coil of wire, with no battery at all. How does a changing magnetic field `push' charges around a circuit? Why does the induced current always fight the change that creates it? And why does a big coil in a DC circuit make the current rise slowly instead of instantly? This chapter answers these questions and explains the physics behind generators, transformers and the ENEO grid itself.

We begin with the central discovery, made by Michael Faraday in 1831: \emph{a changing magnetic flux through a circuit generates an electromotive force (e.m.f.) in that circuit}. This phenomenon is called \cle{electromagnetic induction}. Before we can state the law quantitatively, we need a precise way of measuring ``how much magnetic field passes through a circuit''. That quantity is the magnetic flux.

\courssub{Magnetic Flux and Flux Linkage}

\begin{definition}[Magnetic flux]
The \cle{magnetic flux} $\Phi$ through a plane surface of area $A$ placed in a uniform magnetic field of magnitude $B$ is
\[
\Phi = BA\cos\theta
\]
where $\theta$ is the angle between the field vector $\vec{B}$ and the \emph{normal} (perpendicular) to the surface. The SI unit of flux is the \cle{weber} (Wb): $1\ \mathrm{Wb} = 1\ \mathrm{T\,m^{2}}$.
\end{definition}

Think of magnetic field lines as rain falling vertically and the coil as a hoop you hold in the rain. The amount of rain caught by the hoop depends on three things: how hard it rains ($B$), how big the hoop is ($A$), and how you tilt it ($\theta$). Hold the hoop horizontally (normal vertical, $\theta = 0$) and you catch the maximum, $\Phi = BA$; hold it edge-on ($\theta = 90^{\circ}$) and you catch nothing, $\Phi = 0$. Flux is exactly this ``amount of field threading the loop''.

\begin{definition}[Flux linkage]
For a coil of $N$ turns, each turn of area $A$ threading the same flux $\Phi$, the \cle{flux linkage} is
\[
N\Phi = NBA\cos\theta
\]
Its unit is the weber-turn (Wb-turns), dimensionally identical to the weber.
\end{definition}

\begin{attention}[Angle with the normal, not with the plane]
The angle $\theta$ in $\Phi = BA\cos\theta$ is measured between $\vec{B}$ and the \textbf{normal to the plane of the coil}, not between $\vec{B}$ and the plane itself. If a question gives the angle between the field and the \emph{plane} of the coil, you must use $\Phi = BA\sin\theta$ instead. This is one of the most frequent slips at the GCE.
\end{attention}

\begin{exemplebox}[Computing a flux linkage]
A circular coil of $200$ turns and radius $5.0\ \mathrm{cm}$ is placed with its plane perpendicular to a uniform field $B = 0.30\ \mathrm{T}$. Then $\theta = 0$ and
\[
N\Phi = NBA = 200 \times 0.30 \times \pi (0.050)^{2} = 0.47\ \mathrm{Wb\text{-}turns}.
\]
\end{exemplebox}

\courssub{Experimental Evidence of Electromagnetic Induction}

Faraday's conclusions rest on simple experiments that you can perform in the school laboratory with a bar magnet, a coil of insulated copper wire and a centre-zero galvanometer (a sensitive current detector whose needle can swing left or right).

\begin{experience}[Inducing a current with a magnet and a coil]
Connect the two ends of a coil of many turns to a centre-zero galvanometer. There is \textbf{no battery} anywhere in the circuit.
\begin{enumerate}
\item \textbf{Magnet pushed into the coil:} while the north pole of the magnet is moving towards and into the coil, the needle kicks to one side (say, to the right). The faster the magnet moves, the larger the kick.
\item \textbf{Magnet held at rest inside the coil:} the needle sits exactly at zero, even though the field through the coil is now at its strongest.
\item \textbf{Magnet pulled out:} the needle kicks to the \emph{opposite} side (to the left).
\item \textbf{Coil moved instead of magnet:} the same deflections occur --- only the \emph{relative} motion matters.
\item \textbf{Neighbouring circuit switched on/off:} replace the magnet by a second coil carrying a current from a battery. Nothing happens while the current is steady, but at the instant the switch is \emph{closed} or \emph{opened}, the galvanometer kicks --- in opposite directions for closing and opening.
\end{enumerate}
\end{experience}

\begin{popfigure}
\begin{tikzpicture}[scale=1.0, >=stealth]
% Galvanometer
\draw[thick, popblue] (6.5,-2.6) circle (0.85);
\node[popblue] at (6.5,-2.6) {\textbf{G}};
\draw[thick, popdark] (6.5,-2.6) -- (7.05,-2.05);
\draw[popdark] (6.5,-1.72) arc[start angle=90, end angle=30, radius=0.88];
\draw[popdark] (6.5,-1.72) arc[start angle=90, end angle=150, radius=0.88];
\node[popdark, font=\small] at (7.6,-2.0) {$+$};
\node[popdark, font=\small] at (5.4,-2.0) {$-$};
% Coil (solenoid) drawn as loops
\foreach \x in {2.6,3.0,3.4,3.8,4.2}{
  \draw[thick, poporange] (\x,0.55) ellipse (0.14 and 0.55);
}
\draw[thick, poporange] (2.6,1.1) -- (4.2,1.1);
\draw[thick, poporange] (2.6,-0.55) -- (4.2,-0.55);
% Wires to galvanometer
\draw[thick, popink] (4.2,1.1) -- (5.6,1.1) -- (5.6,-1.9) -- (5.85,-2.2);
\draw[thick, popink] (4.2,-0.55) -- (7.4,-0.55) -- (7.4,-1.9) -- (7.15,-2.2);
% Bar magnet
\draw[thick, fill=popblueL] (-1.4,-0.35) rectangle (0.0,0.35);
\draw[thick, fill=poppink!60] (0.0,-0.35) rectangle (1.4,0.35);
\node at (-0.7,0) {\textbf{S}};
\node at (0.7,0) {\textbf{N}};
% Motion arrow
\draw[->, very thick, popgreen] (1.7,0) -- (2.5,0);
\node[popgreen, font=\small] at (2.1,0.35) {$\vec{v}$};
\node[font=\small, align=center] at (3.4,1.5) {coil of $N$ turns};
\end{tikzpicture}
\legende{A bar magnet pushed into a coil connected to a centre-zero galvanometer. The needle deflects only \emph{while the magnet is moving}; reversing the motion reverses the deflection.}
\end{popfigure}

The crucial observation is this:

\begin{aretenir}[What the experiments show]
\begin{itemize}
\item An e.m.f. (and hence a current) is induced \textbf{only while the flux linkage through the circuit is changing}.
\item The \emph{magnitude} of the induced e.m.f. increases with the \emph{speed} of the change.
\item The \emph{direction} of the induced current depends on the \emph{sense} of the change (flux increasing or decreasing).
\item A steady flux, however large, induces nothing.
\end{itemize}
\end{aretenir}

\courssub{Faraday's Law of Electromagnetic Induction}

Faraday turned these observations into a quantitative law: the induced e.m.f. equals the rate of change of flux linkage.

\begin{propriete}[Faraday's law of electromagnetic induction]
When the flux linkage through a circuit of $N$ turns changes, an e.m.f. is induced given by
\[
e = -N\dfrac{d\Phi}{dt} = -\dfrac{d(N\Phi)}{dt}
\]
The \emph{magnitude} of the induced e.m.f. equals the rate of change of flux linkage; the minus sign expresses the \emph{direction} of the e.m.f. (Lenz's law, studied in the next section). SI check: weber per second = volt, since $1\ \mathrm{V} = 1\ \mathrm{Wb\,s^{-1}}$.
\end{propriete}

The law contains everything the experiments showed. If $N\Phi$ is constant, $d(N\Phi)/dt = 0$ and $e = 0$: no change, no e.m.f. If the change is twice as fast, the e.m.f. is twice as big. If the flux decreases instead of increasing, the sign of the derivative flips, and so does the direction of the e.m.f.

\begin{attention}[A change of flux, not a large flux]
The most common misconception is to believe that a strong field or a large flux induces a large e.m.f. It does not. A coil sitting motionless in the strongest steady magnet in the laboratory has \emph{zero} induced e.m.f., because $d\Phi/dt = 0$. Only the \textbf{rate of change} of flux linkage matters. Always ask yourself first: \emph{what is changing --- $B$, $A$, or $\theta$?}
\end{attention}

Since $\Phi = BA\cos\theta$, there are exactly three independent ways to change the flux through a circuit:
\begin{itemize}
\item \textbf{Change $B$:} move a magnet, switch an electromagnet on or off, or vary the current in a neighbouring coil.
\item \textbf{Change $A$:} deform the loop, slide one rail of a rectangular circuit, or withdraw the coil from the field region (the \emph{effective} area threading the field changes).
\item \textbf{Change $\theta$:} rotate the coil in the field --- this is the principle of the alternator and of every generator in the ENEO grid.
\end{itemize}

\begin{methode}[Three-step recipe for induced-e.m.f. problems]
\begin{enumerate}
\item \textbf{Write} the flux linkage $N\Phi$ as a function of the quantity that is changing ($B(t)$, $A(t)$ or $\theta(t)$).
\item \textbf{Differentiate} with respect to time to get $e = -d(N\Phi)/dt$.
\item \textbf{Take the magnitude} $|e|$ for the size of the e.m.f., and determine the \emph{direction separately} (Lenz's law, next section). Never try to read the direction off the algebra alone.
\end{enumerate}
\end{methode}

\courssub{Induced Current and Induced Charge}

If the circuit is closed with total resistance $R$, the induced e.m.f. drives an induced current
\[
I = \dfrac{|e|}{R} = \dfrac{N}{R}\left|\dfrac{d\Phi}{dt}\right|.
\]
The charge that flows during a change lasting from $t_1$ to $t_2$ is obtained by integrating:
\[
Q = \int_{t_1}^{t_2} I\,dt = \frac{N}{R}\int_{t_1}^{t_2}\left|\frac{d\Phi}{dt}\right|dt = \frac{N\,|\Delta\Phi|}{R}.
\]

\begin{propriete}[Induced charge]
The total charge circulated while the flux linkage changes by $\Delta(N\Phi)$ is
\[
Q = \dfrac{N\,|\Delta\Phi|}{R}
\]
where $R$ is the total resistance of the circuit. Remarkably, $Q$ depends only on the \emph{total change} of flux linkage, \textbf{not on how fast} the change occurs: a slow withdrawal and a quick withdrawal of the coil pass the same total charge (the quick one gives a larger current for a shorter time). This result is frequently examined at the GCE and is the working principle of the \emph{fluxmeter} and of the ballistic galvanometer used to measure magnetic fields.
\end{propriete}

\courssub{Graphical View: from $N\Phi(t)$ to $e(t)$}

Because $e = -d(N\Phi)/dt$, the induced e.m.f. at each instant is \emph{minus the slope} of the flux-linkage--time graph. Where $N\Phi$ is flat, $e = 0$; where $N\Phi$ falls steeply, $e$ is large and positive.

\begin{popfigure}
\begin{tikzpicture}
\begin{axis}[
  width=11cm, height=4.6cm,
  axis lines=middle,
  xlabel={$t$ (s)}, ylabel={$N\Phi$ (Wb-turns)},
  xmin=-0.3, xmax=6.3, ymin=-0.1, ymax=1.3,
  xtick={1,2,3,4,5}, ytick={1},
  domain=0:6, samples=200, thick]
\addplot[popblue, very thick] coordinates {(0,1) (1,1)};
\addplot[popblue, very thick, domain=1:3] {1 - (x-1)/2};
\addplot[popblue, very thick] coordinates {(3,0) (4,0)};
\addplot[popblue, very thick, domain=4:5] {0.6*(x-4)};
\addplot[popblue, very thick] coordinates {(5,0.6) (6,0.6)};
\end{axis}
\end{tikzpicture}

\vspace{2mm}
\begin{tikzpicture}
\begin{axis}[
  width=11cm, height=4.6cm,
  axis lines=middle,
  xlabel={$t$ (s)}, ylabel={$e$ (V)},
  xmin=-0.3, xmax=6.3, ymin=-1.1, ymax=1.1,
  xtick={1,2,3,4,5}, ytick={-0.6,0.5},
  thick]
\addplot[poporange, very thick] coordinates {(0,0) (1,0)};
\addplot[poporange, very thick] coordinates {(1,0.5) (3,0.5)};
\addplot[poporange, very thick] coordinates {(3,0) (4,0)};
\addplot[poporange, very thick] coordinates {(4,-0.6) (5,-0.6)};
\addplot[poporange, very thick] coordinates {(5,0) (6,0)};
\addplot[poporange, thick, dashed] coordinates {(1,0) (1,0.5)};
\addplot[poporange, thick, dashed] coordinates {(3,0.5) (3,0)};
\addplot[poporange, thick, dashed] coordinates {(4,0) (4,-0.6)};
\addplot[poporange, thick, dashed] coordinates {(5,-0.6) (5,0)};
\end{axis}
\end{tikzpicture}
\legende{Top: flux linkage $N\Phi(t)$. Bottom: induced e.m.f. $e = -d(N\Phi)/dt$, which is minus the slope of the upper graph. A steady flux (flat portions) gives zero e.m.f.; a decreasing flux gives a positive e.m.f. and an increasing flux a negative one.}
\end{popfigure}

\courssub{Worked Examples}

\begin{exoresolu}[Coil withdrawn from a magnetic field]
A rectangular coil of $N = 50$ turns and area $A = 2.0\times10^{-3}\ \mathrm{m^{2}}$ lies perpendicular to a uniform field $B = 0.40\ \mathrm{T}$. It is pulled completely out of the field in $0.25\ \mathrm{s}$. The coil has resistance $R = 10\ \Omega$. Find (a) the average induced e.m.f., (b) the average induced current, (c) the total charge that flows.
\tcblower
\textbf{Solution.} The changing quantity is the effective area in the field; the flux goes from its full value to zero.

(a) Initial flux linkage: $N\Phi_1 = NBA = 50 \times 0.40 \times 2.0\times10^{-3} = 4.0\times10^{-2}\ \mathrm{Wb\text{-}turns}$. Final flux linkage: $N\Phi_2 = 0$. By Faraday's law,
\[
|e| = \left|\frac{\Delta(N\Phi)}{\Delta t}\right| = \frac{4.0\times10^{-2}}{0.25} = 0.16\ \mathrm{V}.
\]

(b) Induced current:
\[
I = \frac{|e|}{R} = \frac{0.16}{10} = 1.6\times10^{-2}\ \mathrm{A} = 16\ \mathrm{mA}.
\]

(c) Induced charge (note that the time cancels):
\[
Q = \frac{N|\Delta\Phi|}{R} = \frac{4.0\times10^{-2}}{10} = 4.0\times10^{-3}\ \mathrm{C} = 4.0\ \mathrm{mC}.
\]
Check: $Q = I\,\Delta t = 0.016 \times 0.25 = 4.0\times10^{-3}\ \mathrm{C}$. \checkmark
\end{exoresolu}

\begin{exoresolu}[Field switched off inside a solenoid]
A small search coil of $500$ turns and cross-sectional area $1.0\ \mathrm{cm^{2}}$ sits inside a long solenoid, coaxial with it, where the field is $B = 0.080\ \mathrm{T}$. The solenoid current is switched off and the field falls to zero uniformly in $40\ \mathrm{ms}$. Calculate the average e.m.f. induced in the search coil.
\tcblower
\textbf{Solution.} Here it is $B$ that changes; $N$, $A$ and $\theta = 0$ are fixed.
\[
|e| = NA\left|\frac{\Delta B}{\Delta t}\right| = 500 \times 1.0\times10^{-4} \times \frac{0.080}{0.040} = 0.10\ \mathrm{V}.
\]
This is exactly how the field inside a solenoid is measured in the laboratory: the search coil is connected to a fluxmeter, and the induced charge $Q = N\Delta\Phi/R$ gives $B$.
\end{exoresolu}

\begin{exoresolu}[Rotating coil --- a taste of the alternator]
A flat coil of $N = 100$ turns and area $A = 5.0\times10^{-3}\ \mathrm{m^{2}}$ rotates at constant angular speed $\omega = 20\ \mathrm{rad\,s^{-1}}$ about an axis in its plane, perpendicular to a uniform field $B = 0.25\ \mathrm{T}$. (a) Write the flux linkage as a function of time. (b) Find the maximum induced e.m.f.
\tcblower
\textbf{Solution.} (a) If $\theta = 0$ when the plane of the coil is perpendicular to $\vec{B}$, then $\theta = \omega t$ and
\[
N\Phi(t) = NBA\cos\omega t = 100 \times 0.25 \times 5.0\times10^{-3}\cos(20t) = 0.125\cos(20t)\ \ \mathrm{Wb\text{-}turns}.
\]
(b) Differentiating (step 2 of the method):
\[
e = -\frac{d(N\Phi)}{dt} = 0.125 \times 20 \sin(20t) = 2.5\sin(20t)\ \ \mathrm{V}.
\]
The maximum e.m.f. is $e_{\max} = NBA\omega = 2.5\ \mathrm{V}$. The e.m.f. is \emph{alternating}: this is precisely the principle of the AC generator, developed in Section 3. Note that the e.m.f. is greatest when the flux is \emph{zero} (coil edge-on to the field), because that is where the flux changes fastest.
\end{exoresolu}

\begin{attention}[Pitfalls to avoid at the GCE]
\begin{itemize}
\item Do not write $e = N\Phi/t$ with the \emph{total} flux: the law involves the \emph{change} $\Delta(N\Phi)$.
\item Do not forget the factor $N$: examiners routinely give the number of turns and many candidates drop it.
\item In $Q = N\Delta\Phi/R$, $R$ is the \emph{total} resistance of the whole circuit, including the coil itself.
\item Keep units in SI: areas in $\mathrm{m^{2}}$, times in seconds; convert $\mathrm{cm^{2}}$ ($1\ \mathrm{cm^{2}} = 10^{-4}\ \mathrm{m^{2}}$) and milliseconds before substituting.
\end{itemize}
\end{attention}

\begin{exercice}[]
A coil of $300$ turns, each of area $4.0\times10^{-4}\ \mathrm{m^{2}}$, is placed with its plane perpendicular to a field that increases uniformly from $0.10\ \mathrm{T}$ to $0.50\ \mathrm{T}$ in $0.10\ \mathrm{s}$. The circuit resistance is $6.0\ \Omega$. Calculate (a) the induced e.m.f., (b) the induced current, (c) the charge that circulates. (d) What would the charge be if the same change took $1.0\ \mathrm{s}$ instead?
\end{exercice}

\begin{corrige}[Exercise 1]
(a) $|e| = NA\,\Delta B/\Delta t = 300 \times 4.0\times10^{-4} \times 0.40/0.10 = 0.48\ \mathrm{V}$.
(b) $I = 0.48/6.0 = 0.080\ \mathrm{A}$.
(c) $Q = NA\Delta B/R = 300 \times 4.0\times10^{-4} \times 0.40/6.0 = 8.0\times10^{-3}\ \mathrm{C}$.
(d) The charge is unchanged, $Q = 8.0\ \mathrm{mC}$, because $Q = N\Delta\Phi/R$ does not depend on the duration; only the current ($8.0\ \mathrm{mA}$) and e.m.f. ($48\ \mathrm{mV}$) would be ten times smaller.
\end{corrige}

\begin{aretenir}[Key points of this section]
\begin{itemize}
\item Magnetic flux: $\Phi = BA\cos\theta$ (weber); flux linkage $N\Phi$ for $N$ turns.
\item An e.m.f. is induced \emph{only while the flux linkage is changing}.
\item Faraday's law: $e = -N\,d\Phi/dt$; magnitude = rate of change of flux linkage.
\item Three ways to change the flux: change $B$, change $A$, change $\theta$.
\item Induced current $I = |e|/R$; induced charge $Q = N|\Delta\Phi|/R$, independent of the speed of the change.
\item The minus sign in Faraday's law encodes the \emph{direction} of the e.m.f. --- this is Lenz's law, the subject of the next section.
\end{itemize}
\end{aretenir}

\coursec{Lenz's Law and the Motional E.M.F. (e = Blv)}

In the previous section we discovered Faraday's law: whenever the magnetic flux $\Phi$ linking a circuit \emph{changes}, an e.m.f. is induced, of magnitude $|e| = \left|\dfrac{d\Phi}{dt}\right|$. But a magnitude alone is not enough. To predict the behaviour of a circuit we must also know the \emph{direction} of the induced e.m.f. and of the induced current. This is the question answered by the Russian physicist Heinrich Lenz, and it will lead us to one of the most useful results of this chapter: the motional e.m.f. $e = Blv$.

\courssub{Lenz's Law}

\textbf{A simple observation.}\par

Take a bar magnet and push its North pole towards a coil connected to a sensitive galvanometer. The pointer deflects: a current is induced. Now pull the magnet \emph{away}: the pointer deflects the \emph{other way}. Push the South pole in instead: the deflection reverses again. The induced current always seems to ``fight'' whatever you are doing. This is not a coincidence --- it is a fundamental law of nature.

\begin{propriete}[Lenz's Law (law)]
The direction of the induced e.m.f. (and hence of the induced current) is always such that it \cle{opposes the change} that produces it.
\end{propriete}

Let us make this concrete with the magnet and the coil:

\begin{itemize}
\item A \emph{North pole approaching} the coil increases the flux through it. The induced current flows so that the \emph{near face of the coil becomes a North pole}. Since like poles repel, the coil pushes back against the incoming magnet --- it opposes the approach.
\item A \emph{North pole receding} decreases the flux. The induced current reverses, so the near face becomes a \emph{South pole}, attracting the magnet back --- it opposes the retreat.
\end{itemize}

\begin{popfigure}
\begin{tikzpicture}[scale=1.0]
% coil
\foreach \x in {0,0.25,0.5,0.75,1.0} {
  \draw[thick,popblue] (3.2+\x,0) ellipse (0.09 and 0.8);
}
\draw[thick,popblue] (3.2,0.8) -- (1.6,0.8) -- (1.6,-0.8) -- (3.2,-0.8);
\node[popblue] at (3.7,1.15) {coil};
% magnet
\draw[fill=poporange!70] (-2.2,-0.35) rectangle (-1.1,0.35);
\draw[fill=popink!80] (-1.1,-0.35) rectangle (0.0,0.35);
\node[white] at (-1.65,0) {\textbf{S}};
\node[white] at (-0.55,0) {\textbf{N}};
% velocity arrow
\draw[-{Stealth},very thick,popdark] (0.15,0) -- (1.3,0);
\node[popdark] at (0.7,0.35) {$\vec{v}$};
% induced pole label
\node[popblue] at (3.05,-1.25) {induced \textbf{N} pole};
% current arrows on coil
\draw[-{Stealth},thick,popgreen] (3.2,0.8) arc (90:270:0.09 and 0.8);
\node[popgreen] at (4.6,-1.25) {induced current $I$};
% opposition comment
\node[align=center,popmuted] at (0.7,-1.6) {coil repels the\\approaching magnet};
\end{tikzpicture}
\legende{A North pole approaching a coil: the induced current makes the near face a North pole, opposing the approach.}
\end{popfigure}

\textbf{The pole rule.} A quick way to remember the result:

\begin{aretenir}[Pole rule for a magnet and a coil]
An \emph{approaching} pole induces the \emph{same} pole on the near face of the coil (repulsion); a \emph{receding} pole induces the \emph{opposite} pole (attraction). Approaching North $\Rightarrow$ near face becomes North.
\end{aretenir}

\textbf{The minus sign in Faraday's law.}\par

We can now write Faraday's law in its complete form:
\[
e = -\frac{d\Phi}{dt} = -N\frac{d\Phi_B}{dt}
\]
where $\Phi = N\Phi_B$ is the total flux linkage of a coil of $N$ turns, each turn carrying flux $\Phi_B$. The \cle{minus sign} is the mathematical expression of Lenz's law. If we choose a positive sense of circulation around the coil and define the flux $\Phi$ as positive when it passes in the corresponding direction (right-hand grip rule), then the induced e.m.f. $e$ is \emph{negative} when the flux is \emph{increasing} ($d\Phi/dt > 0$): the induced e.m.f. drives current in the negative sense, creating a flux that \emph{cancels the increase}. The minus sign is not a decoration --- it is Lenz's law written in algebra.

\begin{attention}[Common mistake]
Do not say ``the induced current opposes the magnetic field''. It opposes the \emph{change of flux}, not the field itself. If the flux is decreasing, the induced current actually \emph{reinforces} the existing field to slow the decrease.
\end{attention}

\courssub{Lenz's Law and the Conservation of Energy}

Why must nature behave this way? Suppose, for a moment, that the induced current \emph{helped} the change instead of opposing it. Then an approaching North pole would induce a South pole on the near face, attracting the magnet, accelerating it further, inducing more current, attracting it more strongly\ldots{} We would obtain an ever-growing current and an ever-faster magnet \emph{for free}. This is impossible: it violates the \cle{conservation of energy}.

\begin{propriete}[Energy conservation in induction (property)]
Lenz's law is a direct consequence of the conservation of energy. To change the flux through a circuit, an external agent must do \cle{mechanical work} against the electromagnetic opposition; this work is converted into \cle{electrical energy} in the circuit (and ultimately into heat in the resistance):
\[
\text{mechanical work done} = \text{electrical energy produced} \; (+\ \text{heat}).
\]
\end{propriete}

When you push the magnet into the coil, you \emph{feel} the repulsion: you must push harder, and the work you supply is exactly what lights the lamp connected to the coil. Every dynamo, every alternator at the Song Loulou or Memve'ele hydroelectric stations, obeys this principle: the turbines supply mechanical work, and Lenz's law guarantees that this work --- and no less --- is needed to generate the electrical energy.

\courssub{Fleming's Right-Hand Rule}

For a straight conductor moving through a magnetic field, there is a practical tool for finding the direction of the induced current: \cle{Fleming's right-hand (dynamo) rule}.

\begin{methode}[Using Fleming's right-hand rule]
Hold the thumb, first finger and second finger of the \textbf{right} hand mutually at right angles:
\begin{enumerate}
\item \textbf{F}irst finger $\rightarrow$ \textbf{F}ield (direction of $\vec{B}$, North to South);
\item thu\textbf{M}b $\rightarrow$ \textbf{M}otion (direction of the velocity $\vec{v}$ of the conductor);
\item se\textbf{C}ond finger $\rightarrow$ \textbf{C}urrent (direction of the induced conventional current).
\end{enumerate}
\end{methode}

\begin{attention}[Right hand or left hand?]
A classic GCE trap! Use the \textbf{right} hand for \emph{induction} (generators, dynamos: motion $\rightarrow$ current). Use the \textbf{left} hand for the \emph{motor effect} (force on a current-carrying conductor: current + field $\rightarrow$ force). Memory aid: \textbf{RI}ght = \textbf{I}nduction; \textbf{L}eft = motor ($\vec{F} = I\vec{l}\wedge\vec{B}$).
\end{attention}

\courssub{The Motional E.M.F.: Derivation of e = Blv}

Consider a straight conducting rod $PQ$ of length $l$ sliding with constant velocity $\vec{v}$ on two parallel conducting rails, in a uniform magnetic field $\vec{B}$ directed into the page. The rails are joined at one end, so the rod and rails form a closed rectangular circuit.

\begin{popfigure}
\begin{tikzpicture}[scale=1.1]
% field into page
\foreach \x in {0.4,1.4,2.4,3.4,4.4} {
  \foreach \y in {0.5,1.5} {
    \node[popmuted] at (\x,\y) {$\times$};
  }
}
\node[popmuted] at (4.9,1.9) {$\vec{B}$ (into page)};
% rails
\draw[thick] (0,0) -- (5,0);
\draw[thick] (0,2) -- (5,2);
\draw[thick] (0,0) -- (0,2);
% rod
\draw[line width=2.5pt,popblue] (2.6,0) -- (2.6,2);
\node[popblue] at (2.6,2.3) {rod, length $l$};
% velocity
\draw[-{Stealth},very thick,popdark] (2.6,1.0) -- (3.7,1.0);
\node[popdark] at (3.15,1.3) {$\vec{v}$};
% current arrows
\draw[-{Stealth},thick,popgreen] (2.6,1.6) -- (2.6,0.6);
\node[popgreen] at (2.25,1.1) {$I$};
\draw[-{Stealth},thick,popgreen] (1.6,0) -- (0.7,0);
\draw[-{Stealth},thick,popgreen] (0.7,2) -- (1.6,2);
% opposing force
\draw[-{Stealth},very thick,poporange] (2.6,0.4) -- (1.6,0.4);
\node[poporange] at (2.1,0.05) {$\vec{F}$};
\end{tikzpicture}
\legende{Rod sliding on U-shaped rails in a field into the page. The induced current $I$ flows round the loop; the magnetic force $\vec{F}$ on the rod opposes its motion (Lenz's law).}
\end{popfigure}

\textbf{Derivation from Faraday's law.}\par

Let $x$ be the distance of the rod from the closed end of the rails. The area of the circuit is $A = lx$, so the flux linking the circuit is
\[
\Phi = BA = Blx .
\]
As the rod moves, $x$ changes at the rate $\dfrac{dx}{dt} = v$. Applying Faraday's law:
\[
|e| = \frac{d\Phi}{dt} = Bl\,\frac{dx}{dt} = Blv .
\]

\begin{propriete}[Motional e.m.f. (formula)]
A straight conductor of length $l$ moving with velocity $v$ perpendicular to a uniform magnetic field $B$ (conductor, field and velocity mutually perpendicular) has an induced e.m.f.
\[
\boxed{e = Blv}
\]
with $B$ in tesla (T), $l$ in metres (m), $v$ in $\mathrm{m\,s^{-1}}$, $e$ in volts (V). If the velocity makes an angle $\theta$ with the field, only the component of $\vec{v}$ perpendicular to $\vec{B}$ counts:
\[
e = Blv\sin\theta .
\]
\end{propriete}

\textbf{Microscopic origin.} The free electrons in the rod are carried along at velocity $\vec{v}$ through the field, so each electron experiences a magnetic force $evB$ along the rod. Electrons accumulate at one end, leaving the other end positive: the rod becomes a battery of e.m.f. $Blv$. Fleming's right-hand rule gives the direction of the conventional current this ``battery'' drives.

\courssub{Force on the Induced Current and the Power Balance}

Once the induced current $I$ flows in the rod, the rod is a current-carrying conductor sitting in a magnetic field: it experiences the Laplace force
\[
F = BIl .
\]
By Fleming's \emph{left}-hand rule, this force points \emph{opposite} to $\vec{v}$ --- exactly as Lenz's law demands. To keep the rod moving at constant velocity, an external agent must pull with an equal force $F = BIl$.

Now check the energy balance. The mechanical power supplied by the agent is
\[
P_{\text{mech}} = Fv = BIl \times v = (Blv)\,I = eI = P_{\text{elec}} .
\]
The mechanical power is converted \emph{entirely} into electrical power $eI$ (which then appears as heat $RI^2$ in the circuit resistance). The equations $e = Blv$ and $F = BIl$ are two faces of the same energy conversion.

\begin{exoresolu}[Rod sliding on conducting rails]
A rod of length $l = 0.50\ \mathrm{m}$ slides at constant speed $v = 4.0\ \mathrm{m\,s^{-1}}$ on frictionless horizontal rails, in a uniform vertical field $B = 0.20\ \mathrm{T}$. The total resistance of the circuit is $R = 0.10\ \Omega$. Find (a) the induced e.m.f., (b) the induced current, (c) the force needed to keep the rod moving, (d) the mechanical power and the electrical power dissipated.
\tcblower
\textbf{Solution.}
\begin{enumerate}
\item[(a)] Motional e.m.f.: $e = Blv = 0.20 \times 0.50 \times 4.0 = \SI{0.40}{V}$.
\item[(b)] Induced current: $I = \dfrac{e}{R} = \dfrac{0.40}{0.10} = \SI{4.0}{A}$.
\item[(c)] Opposing force: $F = BIl = 0.20 \times 4.0 \times 0.50 = \SI{0.40}{N}$. The agent must pull with $0.40\ \mathrm{N}$.
\item[(d)] Mechanical power: $P = Fv = 0.40 \times 4.0 = \SI{1.6}{W}$. Electrical power: $P = eI = 0.40 \times 4.0 = \SI{1.6}{W}$ (equivalently $RI^2 = 0.10 \times 16 = 1.6\ \mathrm{W}$). The two are equal: energy is conserved.
\end{enumerate}
\end{exoresolu}

\begin{exoresolu}[Aeroplane wing cutting the Earth's field]
An aeroplane with wingspan $l = 30\ \mathrm{m}$ flies horizontally at $v = 250\ \mathrm{m\,s^{-1}}$ over a region where the vertical component of the Earth's magnetic field is $B = 4.0\times 10^{-5}\ \mathrm{T}$. Calculate the e.m.f. induced between the wing tips.
\tcblower
\textbf{Solution.} The wings, the velocity and the vertical field component are mutually perpendicular, so
\[
e = Blv = 4.0\times 10^{-5} \times 30 \times 250 = \SI{0.30}{V}.
\]
A small but measurable e.m.f. of $0.30\ \mathrm{V}$ appears between the wing tips. (No current flows, because the circuit is not closed.)
\end{exoresolu}

\courssub{Eddy Currents and Electromagnetic Braking}

Induction does not only occur in wires. When the flux through a \emph{block} of metal changes, currents are induced in closed loops \emph{inside the material}: these are called \cle{eddy currents}. By Lenz's law, they always oppose the change producing them.

\begin{experience}[The falling magnet in a copper tube]
Drop a small strong magnet down a vertical copper tube. Astonishingly, it falls in slow motion, far slower than a non-magnetic object of the same mass. As the magnet falls, the flux through each ring-shaped slice of the tube changes, inducing eddy currents. By Lenz's law, the eddy currents \emph{below} the magnet repel the approaching pole, while those \emph{above} the magnet attract the receding pole: the net effect is an upward force opposing the fall. The magnet quickly reaches a small terminal velocity when this force balances its weight. The gravitational potential energy is converted, via the eddy currents, into heat in the copper.
\end{experience}

\begin{savaistu}[Electromagnetic braking]
The same principle brakes some trains, lorries and roller-coasters: a strong electromagnet induces eddy currents in a metal disc or rail attached to the wheel, and the opposing force slows the vehicle \emph{without any contact or wear}. The kinetic energy is dissipated as heat. Eddy currents are a nuisance in transformer cores (they waste energy), which is why cores are made of thin insulated \emph{laminations} that break the current loops.
\end{savaistu}

\begin{exercice}[Checking your understanding]
\begin{enumerate}
\item A South pole is withdrawn from a coil. State the polarity induced on the near face and justify using Lenz's law.
\item A rectangular coil is pulled horizontally out of a uniform field directed into the page. In which sense does the induced current flow?
\item A metal rod $0.80\ \mathrm{m}$ long moves at $15\ \mathrm{m\,s^{-1}}$ at $60^\circ$ to a field of $0.50\ \mathrm{T}$, the rod being perpendicular to both. Calculate the induced e.m.f.
\end{enumerate}
\end{exercice}

\begin{corrige}[Exercise]
\begin{enumerate}
\item Withdrawing a South pole decreases the flux; the coil opposes the retreat by \emph{attracting} the magnet, so the near face becomes a \textbf{North} pole.
\item The flux directed into the page decreases, so the induced current creates its own field \emph{into the page}: by the right-hand grip rule it flows \textbf{clockwise}.
\item $e = Blv\sin\theta = 0.50 \times 0.80 \times 15 \times \sin 60^\circ = 6.0 \times 0.866 = \SI{5.2}{V}$.
\end{enumerate}
\end{corrige}

\begin{aretenir}[Key points of this section]
\begin{itemize}
\item \textbf{Lenz's law}: the induced e.m.f. opposes the change producing it; this is the meaning of the minus sign in $e = -\dfrac{d\Phi}{dt}$.
\item Approaching North $\Rightarrow$ near face of the coil becomes North (repulsion); receding North $\Rightarrow$ near face becomes South (attraction).
\item \textbf{Fleming's right-hand rule} (induction): First finger Field, seCond finger Current, thuMb Motion. Right hand = generator; left hand = motor.
\item Motional e.m.f.: $e = Blv$ (mutually perpendicular), or $e = Blv\sin\theta$ in general.
\item The induced current feels a force $F = BIl$ opposing the motion; mechanical power $Fv$ equals electrical power $eI$: energy is conserved.
\item Changing flux in bulk metal produces \textbf{eddy currents}, used in electromagnetic braking and responsible for the slow fall of a magnet in a copper tube.
\end{itemize}
\end{aretenir}

\coursec{Simple AC and DC Generators}

In the previous section we established that a conductor moving through a magnetic field experiences a motional e.m.f. $e = Blv$, and that Lenz's law fixes its direction. A natural question follows: how can we make a conductor move \emph{continuously} through a magnetic field so as to produce a \emph{continuous} supply of electrical energy? The answer is one of the most consequential inventions in human history: the \cle{generator}. Instead of pushing a straight rod along rails, we \emph{rotate a coil} in a magnetic field. Rotation is easy to sustain --- a waterfall, steam from a boiler, or a petrol engine can all spin a shaft --- and, as we shall now show, a rotating coil automatically produces an alternating e.m.f. This section builds the simple AC generator from Faraday's law, then shows how a single clever modification --- the split-ring commutator --- turns it into a DC generator.

\courssub{The Simple AC Generator: Construction}

\begin{definition}[Simple AC generator]
A \cle{simple AC generator} (or \cle{alternator}) consists of a rectangular coil of $N$ turns, each of area $A$, mounted on an axle and rotated at constant angular velocity $\omega$ in a uniform magnetic field $\vec{B}$. The ends of the coil are connected to two \cle{slip rings} which rotate with the coil; stationary carbon \cle{brushes} press against the slip rings and connect the coil to the external circuit.
\end{definition}

Each component has a precise role:
\begin{itemize}
\item The \textbf{coil} ($N$ turns, area $A$): this is where the e.m.f. is induced. Using many turns multiplies the flux linkage and hence the e.m.f.
\item The \textbf{magnet} (permanent magnet or electromagnet): provides the uniform field $\vec{B}$ across the coil.
\item The \textbf{axle}: driven by an external mechanical agent (turbine, engine, hand crank); it supplies the mechanical work that will be converted into electrical energy.
\item The \textbf{slip rings}: two complete copper rings, each permanently connected to one end of the coil. They rotate with the coil, so each end of the coil always remains connected to the same terminal of the external circuit.
\item The \textbf{brushes}: small blocks of carbon held by springs against the slip rings. They rub on the rotating rings and carry the current to the stationary external circuit without twisting the wires.
\end{itemize}

\begin{popfigure}
\begin{tikzpicture}[scale=1.0, line cap=round]
% Magnet poles
\fill[popblueL] (-3.6,-1.6) rectangle (-2.4,1.6);
\fill[popblueL] (2.4,-1.6) rectangle (3.6,1.6);
\node at (-3.0,0) {\textbf{N}};
\node at (3.0,0) {\textbf{S}};
% Field lines
\foreach \y in {-1.0,-0.5,0,0.5,1.0}{
  \draw[-latex, popblue, thin] (-2.35,\y) -- (2.35,\y);
}
\node[popblue] at (0,1.35) {$\vec{B}$};
% Coil (tilted rectangle, seen in perspective)
\draw[very thick, poporange] (-1.1,-0.75) -- (1.1,-0.45) -- (1.1,0.75) -- (-1.1,0.45) -- cycle;
\node[poporange] at (0,-1.15) {coil, $N$ turns, area $A$};
% Axle
\draw[thick, popdark] (0,0.6) -- (0,2.1);
\draw[-latex, popdark] (0.55,1.9) arc[start angle=70, end angle=200, radius=0.45];
\node[popdark] at (1.35,1.95) {$\omega$};
% Slip rings (two separate rings on the axle)
\draw[thick, popgold] (-0.30,2.05) ellipse (0.10 and 0.22);
\draw[thick, popgold] (0.30,2.05) ellipse (0.10 and 0.22);
\node[popgold, align=center, text width=2.2cm] at (-1.9,2.45) {two slip rings};
\draw[-latex, popgold] (-1.15,2.35) -- (-0.42,2.10);
% Brushes
\fill[popink] (-0.62,1.95) rectangle (-0.40,2.18);
\fill[popink] (0.40,1.95) rectangle (0.62,2.18);
\node[popink, align=center, text width=1.6cm] at (1.8,2.45) {brushes};
\draw[-latex, popink] (1.35,2.35) -- (0.64,2.10);
% External circuit
\draw[thick] (-0.51,2.18) -- (-0.51,3.3) -- (-0.51,3.3);
\draw[thick] (0.51,2.18) -- (0.51,3.3);
\draw[thick] (-0.51,3.3) -- (-0.51,3.8) to[R=$R$] (0.51,3.8) -- (0.51,3.3);
\node at (0,4.15) {external circuit};
\end{tikzpicture}
\legende{The simple AC generator: a coil rotating in a uniform field, connected to the external circuit through two slip rings and two brushes.}
\end{popfigure}

\courssub{Derivation of the Induced E.M.F.: $e = NBA\omega\sin(\omega t)$}

This derivation is \textbf{examinable} and follows directly from Faraday's law.

\textbf{Setting up.} Let the coil rotate with constant angular velocity $\omega$ about an axis perpendicular to $\vec{B}$. Choose the origin of time so that at $t = 0$ the \emph{plane of the coil is perpendicular to} $\vec{B}$, i.e. the normal $\vec{n}$ to the coil is parallel to $\vec{B}$. After a time $t$, the normal makes an angle
\[ \theta = \omega t \]
with the field direction.

\textbf{Step 1: flux linkage.} The flux through one turn is $\Phi = BA\cos\theta$, so the total flux linkage is
\[ N\Phi = NBA\cos(\omega t). \]

\textbf{Step 2: Faraday's law.} The induced e.m.f. is the negative rate of change of flux linkage:
\begin{align*}
e &= -\frac{d(N\Phi)}{dt} = -\frac{d}{dt}\big[NBA\cos(\omega t)\big] \\
  &= -NBA\,\big(-\omega\sin(\omega t)\big) = NBA\omega\sin(\omega t).
\end{align*}

\begin{propriete}[E.M.F. of a rotating coil]
A coil of $N$ turns and area $A$ rotating at angular frequency $\omega$ in a uniform field $B$ develops a sinusoidal e.m.f.
\[ e = NBA\omega\sin(\omega t) = e_0\sin(\omega t), \qquad \text{with peak e.m.f. } e_0 = NBA\omega. \]
(Here $t = 0$ is taken when the coil plane is perpendicular to $\vec{B}$.) The frequency of the output is $f = \omega/(2\pi)$ and the r.m.s. value is $E_{\mathrm{rms}} = e_0/\sqrt{2}$. The derivation from Faraday's law is examinable.
\end{propriete}

Note that $e_0 = NBA\omega$ shows exactly how to build a bigger e.m.f.: stronger field $B$, larger area $A$, more turns $N$, faster rotation $\omega$. This is the design recipe for every real alternator.

\begin{attention}[The origin of time matters]
A very common mistake is to write $e = NBA\omega\cos(\omega t)$ ``because flux is a cosine''. Whether the e.m.f. is a sine or a cosine depends entirely on \textbf{where you take $t = 0$}. With $t = 0$ at the perpendicular position (normal parallel to $\vec{B}$), $N\Phi = NBA\cos(\omega t)$ and $e = NBA\omega\sin(\omega t)$. If instead $t = 0$ is taken with the coil plane \emph{parallel} to $\vec{B}$, then $N\Phi = NBA\sin(\omega t)$ and $e = NBA\omega\cos(\omega t)$. \textbf{Always state the position of the coil at $t = 0$ before writing the equation.} The peak value $e_0 = NBA\omega$ is unaffected.
\end{attention}

\courssub{Interpreting the Graph: Where is the E.M.F. Maximum?}

\begin{popfigure}
\begin{tikzpicture}
\begin{axis}[
  width=13cm, height=6.2cm,
  axis lines=middle,
  xlabel={$\omega t$}, ylabel={$e$},
  xtick={0,1.571,3.142,4.712,6.283},
  xticklabels={$0$,$\frac{\pi}{2}$,$\pi$,$\frac{3\pi}{2}$,$2\pi$},
  ytick={1,-1}, yticklabels={$e_0$,$-e_0$},
  xmin=-0.25, xmax=6.7, ymin=-1.5, ymax=1.6,
  samples=200, domain=0:6.283, thick,
  legend style={at={(0.98,0.97)}, anchor=north east, font=\small}
]
\addplot[popblue, very thick] {sin(deg(x))};
\addlegendentry{AC output: $e = e_0\sin(\omega t)$}
\addplot[poporange, very thick, dashed] {abs(sin(deg(x)))};
\addlegendentry{DC output: $e = |e_0\sin(\omega t)|$}
\end{axis}
\end{tikzpicture}
\legende{Output e.m.f. of the simple AC generator (sinusoid) and of the simple DC generator (rectified, pulsating) on the same axes.}
\end{popfigure}

The graph raises a subtle point that examiners love. There are two remarkable positions of the coil:

\begin{itemize}
\item \textbf{Coil plane perpendicular to $\vec{B}$} ($\omega t = 0, \pi, \dots$): the flux linkage $N\Phi$ is \emph{maximum}, but at that instant it is passing through its peak, so its \emph{rate of change is zero}: $e = 0$.
\item \textbf{Coil plane parallel to $\vec{B}$} ($\omega t = \pi/2, 3\pi/2, \dots$): the flux through the coil is \emph{zero}, but it is changing \emph{fastest} (the long sides of the coil are then cutting field lines at the greatest rate): $|e| = e_0$, maximum.
\end{itemize}

\begin{methode}[Locating maxima and zeros of the e.m.f.]
\begin{enumerate}
\item Sketch or picture the coil at the instant considered and ask: ``is the flux itself large, or is the flux \emph{changing} fast?''
\item The e.m.f. depends only on the \textbf{rate of change} of flux linkage (Faraday), never on the flux itself.
\item Conclusion: $e$ is \textbf{maximum} when the coil plane is \textbf{parallel} to $\vec{B}$ (flux zero, rate of change maximum); $e$ is \textbf{zero} when the coil plane is \textbf{perpendicular} to $\vec{B}$ (flux maximum, rate of change zero).
\item Check with the motional picture: the e.m.f. is greatest when the long sides of the coil move \emph{perpendicularly} across the field lines ($e = Blv$ per side, all $N$ turns adding up).
\end{enumerate}
\end{methode}

\begin{attention}[Flux maximum does not mean e.m.f. maximum]
Students frequently write ``the e.m.f. is maximum when the flux through the coil is maximum''. This is wrong and costs marks. The induced e.m.f. is $e = -d(N\Phi)/dt$: it is proportional to the \emph{slope} of the flux--time graph, not to the flux. When the flux is at its peak, its slope is zero, so $e = 0$.
\end{attention}

\courssub{Energy Conservation in the Generator}

Why must the turbine keep pushing? When the coil delivers current to the external circuit, that current flows through the coil \emph{inside the magnetic field}, so each side of the coil experiences a force $\vec{F} = I\vec{L}\wedge\vec{B}$. By Lenz's law, the direction of the induced current is always such that these forces create a \textbf{torque opposing the rotation}. The external agent must therefore work continuously against this braking torque: the mechanical work done per second against the opposing torque is exactly the electrical power delivered (plus small losses). A generator is thus an \emph{energy converter}: mechanical energy $\rightarrow$ electrical energy. If the external circuit is opened (no current), the opposing torque vanishes and the coil spins freely --- a fact anyone who has pedalled a bicycle dynamo with the lamp switched off will recognise.

\courssub{The Simple DC Generator: the Split-Ring Commutator}

The AC generator delivers a current that reverses direction every half-turn. For some applications (charging batteries, electrolysis, older DC motors) we want a current that keeps \emph{one direction} in the external circuit, even if its magnitude varies. The solution is beautifully simple.

\begin{definition}[Simple DC generator]
A \cle{simple DC generator} is an AC generator in which the two slip rings are replaced by a \cle{split-ring commutator}: a single copper ring cut into two insulated half-rings, each connected to one end of the coil. The brushes are placed so that each brush changes from one half-ring to the other exactly at the instant when the e.m.f. passes through zero.
\end{definition}

\begin{propriete}[Action of the commutator]
Every half-turn, the connections between the coil ends and the external terminals are \textbf{swapped} by the split-ring commutator, at precisely the moment the induced e.m.f. reverses. The two reversals cancel each other: although the e.m.f. inside the coil is still alternating, the e.m.f. delivered to the external circuit is
\[ e_{\text{out}} = \left| e_0\sin(\omega t) \right|, \]
a \emph{rectified, pulsating} e.m.f. that never changes sign. The external current therefore keeps one direction.
\end{propriete}

\begin{popfigure}
\begin{tikzpicture}[scale=1.05, line cap=round]
% Commutator: two half rings
\draw[very thick, popgold] (0.45,0) arc[start angle=0, end angle=170, radius=0.45];
\draw[very thick, popgold] (-0.45,0) arc[start angle=180, end angle=350, radius=0.45];
% gaps
\fill[white] (0.38,0.10) rectangle (0.55,-0.02);
\fill[white] (-0.38,-0.10) rectangle (-0.55,0.02);
% axle
\fill[popmuted] (0,0) circle (0.12);
\node[below, popmuted] at (0,-0.16) {axle};
% coil connections
\draw[thick, poporange] (0.45,0.12) -- (1.5,0.9);
\draw[thick, poporange] (-0.45,-0.12) -- (-1.5,-0.9);
\node[poporange, right, text width=2.6cm] at (1.5,0.9) {to one end of coil};
\node[poporange, left, text width=2.6cm] at (-1.5,-0.9) {to other end of coil};
% brushes
\fill[popink] (0.42,0.55) rectangle (0.72,0.85);
\fill[popink] (-0.72,-0.55) rectangle (-0.42,-0.85);
\draw[thick, popink] (0.57,0.85) -- (0.57,1.35);
\draw[thick, popink] (-0.57,-0.85) -- (-0.57,-1.35);
\node[popink, right] at (0.62,1.2) {brush $+$};
\node[popink, left] at (-0.62,-1.2) {brush $-$};
% external
\draw[thick] (0.57,1.35) -- (2.4,1.35) -- (2.4,-1.35) -- (-0.57,-1.35);
\draw[thick] (2.4,0.25) to[R=$R$] (2.4,-0.25);
% labels
\node[popgold, align=center, text width=3.4cm] at (2.6,0.75) {two half-rings (insulated gap)};
\draw[-latex, popgold] (1.9,0.6) -- (0.42,0.28);
\node[align=center, text width=4.6cm] at (-2.6,0.6) {at reversal, each brush swaps half-ring};
\draw[-latex, popdark] (-1.9,0.45) -- (-0.5,0.1);
\end{tikzpicture}
\legende{The split-ring commutator at the moment of reversal: the brushes pass over the insulating gap and swap half-rings exactly when $e = 0$.}
\end{popfigure}

The output is shown as the dashed orange curve in the graph above: the negative half-cycles are ``folded up'', giving $|e_0\sin(\omega t)|$. The current is unidirectional but far from smooth --- real DC generators use many coils at different angles to smooth it out.

\courssub{Comparing the Two Generators; Practical Improvements}

\begin{center}
\begin{tabularx}{\linewidth}{|l|X|X|}
\hline
\rowcolor{popblueL} \textbf{Feature} & \textbf{Simple AC generator} & \textbf{Simple DC generator} \\
\hline
Rotating part & Coil of $N$ turns, area $A$ & Coil of $N$ turns, area $A$ \\
\hline
Connection to coil & Two slip rings & Split-ring commutator (two half-rings) \\
\hline
Brushes & One per slip ring & One per half-ring pair, swapping every half-turn \\
\hline
E.M.F. in the coil & $e = e_0\sin(\omega t)$ & $e = e_0\sin(\omega t)$ (same) \\
\hline
Output e.m.f. & $e_0\sin(\omega t)$ (alternating) & $|e_0\sin(\omega t)|$ (pulsating, one sign) \\
\hline
Current in external circuit & Reverses each half-cycle & Keeps one direction \\
\hline
Peak e.m.f. & $e_0 = NBA\omega$ & $e_0 = NBA\omega$ \\
\hline
\end{tabularx}
\end{center}

Real generators improve on the simple model in several standard ways:
\begin{itemize}
\item \textbf{Many turns} of insulated copper wire: $e_0 \propto N$.
\item A \textbf{soft-iron core} inside the coil: it concentrates and guides the field lines, greatly increasing the effective $B$ through the coil (and is laminated to reduce eddy-current losses).
\item \textbf{Electromagnets} instead of permanent magnets: they provide a much stronger and adjustable field $B$.
\item \textbf{Several coils} at angles to one another (and several commutator segments): the outputs add to give a smoother, nearly steady DC, or three-phase AC in power-station alternators.
\end{itemize}

\begin{savaistu}[Generators and Cameroon]
Nearly all of Cameroon's electricity comes from alternators working on exactly this principle. At the hydroelectric stations of \textbf{Songloulou} on the Sanaga river and \textbf{Edéa} further downstream on the same river, falling water spins turbines coupled to huge alternators; the rotating magnetic field induces three-phase alternating e.m.f. in stationary coils at $f = \SI{50}{Hz}$. Thermal stations (gas or heavy-fuel-oil plants near Douala and Kribi) use steam or gas turbines to the same end. On a much smaller scale, the alternator of a motorcycle or the dynamo of a bicycle lamp is a miniature rotating-coil (or rotating-magnet) generator: the faster the engine turns, the larger $\omega$, and the brighter the lamp --- a direct, everyday illustration of $e_0 = NBA\omega$.
\end{savaistu}

\courssub{Worked Examples}

\begin{exoresolu}[Peak e.m.f. of a rotating coil]
A rectangular coil of $N = 200$ turns has dimensions $\SI{10}{cm} \times \SI{8.0}{cm}$. It rotates at $3000$ revolutions per minute about an axis perpendicular to a uniform magnetic field of flux density $B = \SI{0.40}{T}$. (a) Calculate the peak e.m.f. induced in the coil. (b) Write the equation of the e.m.f., taking $t = 0$ when the plane of the coil is perpendicular to the field. (c) State the frequency and the r.m.s. value of the e.m.f.
\tcblower
\textbf{(a)} Convert all data to SI units:
\[ A = 0.10 \times 0.080 = 8.0\times 10^{-3}\ \mathrm{m^2}, \qquad \omega = \frac{3000 \times 2\pi}{60} = 100\pi\ \mathrm{rad\,s^{-1}} \approx 314\ \mathrm{rad\,s^{-1}}. \]
Peak e.m.f.:
\[ e_0 = NBA\omega = 200 \times 0.40 \times 8.0\times 10^{-3} \times 314 \approx 2.0\times 10^{2}\ \mathrm{V}. \]
\textbf{(b)} With $t = 0$ at the perpendicular position, $N\Phi = NBA\cos(\omega t)$, so
\[ e = e_0\sin(\omega t) = 201\sin(314\,t)\ \mathrm{V}. \]
\textbf{(c)} Frequency: $f = \dfrac{\omega}{2\pi} = \dfrac{3000}{60} = \SI{50}{Hz}$. R.m.s. value:
\[ E_{\mathrm{rms}} = \frac{e_0}{\sqrt{2}} = \frac{201}{\sqrt{2}} \approx \SI{142}{V}. \]
Note the pleasing result: $3000$ rpm gives exactly the mains frequency of $\SI{50}{Hz}$ used in Cameroon.
\end{exoresolu}

\begin{exoresolu}[Designing for a target e.m.f.]
A simple AC generator must produce a peak e.m.f. of $\SI{170}{V}$ at $\SI{50}{Hz}$. The available magnet provides $B = \SI{0.25}{T}$ and the coil has area $A = \SI{2.0e-2}{m^2}$. How many turns are required?
\tcblower
Angular frequency: $\omega = 2\pi f = 2\pi \times 50 = 100\pi \approx \SI{314}{rad.s^{-1}}$.
From $e_0 = NBA\omega$:
\[ N = \frac{e_0}{BA\omega} = \frac{170}{0.25 \times 2.0\times 10^{-2} \times 314} = \frac{170}{1.57} \approx 108. \]
So about $\mathbf{108}$ \textbf{turns} (in practice one would wind $110$ turns to allow for losses). Check the units: $\mathrm{T\,m^2\,s^{-1}} = \mathrm{Wb\,s^{-1}} = \mathrm{V}$, as required.
\end{exoresolu}

\begin{exercice}[From AC to DC]
A coil of $150$ turns and area $\SI{5.0e-3}{m^2}$ rotates at $\SI{1200}{rpm}$ in a field of $\SI{0.60}{T}$, first with slip rings, then with a split-ring commutator. (a) Calculate the peak e.m.f. and the frequency. (b) Sketch the output e.m.f. in both cases over one full period. (c) Explain why the average value of the DC output is $\dfrac{2e_0}{\pi}$ and compute it.
\end{exercice}

\begin{corrige}[Exercise 1]
(a) $\omega = \dfrac{1200 \times 2\pi}{60} = 40\pi \approx \SI{126}{rad.s^{-1}}$; $e_0 = NBA\omega = 150 \times 0.60 \times 5.0\times10^{-3} \times 126 \approx \SI{57}{V}$; $f = \SI{20}{Hz}$.
(b) With slip rings: sinusoid $57\sin(126\,t)$. With commutator: the same curve with negative half-cycles reflected upwards, i.e. $|57\sin(126\,t)|$.
(c) Averaging over a half-period: $\langle e \rangle = \dfrac{1}{\pi}\displaystyle\int_0^{\pi} e_0\sin\theta\, d\theta = \dfrac{e_0}{\pi}\big[-\cos\theta\big]_0^{\pi} = \dfrac{2e_0}{\pi} \approx \SI{36}{V}$.
\end{corrige}

\begin{aretenir}[Key points]
\begin{itemize}
\item A coil rotating at angular frequency $\omega$ in a field $B$ produces $e = NBA\omega\sin(\omega t)$, with peak e.m.f. $e_0 = NBA\omega$ (taking $t = 0$ with the coil plane perpendicular to $\vec{B}$).
\item The e.m.f. is \textbf{maximum when the coil plane is parallel to} $\vec{B}$ (flux zero, rate of change of flux maximum) and \textbf{zero when the coil plane is perpendicular to} $\vec{B}$.
\item Slip rings + brushes give an \textbf{AC} output; a \textbf{split-ring commutator} swaps the connections every half-turn and gives a pulsating \textbf{DC} output $|e_0\sin(\omega t)|$.
\item The induced current creates a torque opposing the rotation (Lenz's law): the mechanical work done against it is the source of the electrical energy.
\item Practical generators use many turns, a laminated soft-iron core, electromagnets and several coils; alternators of this type power Cameroon from Songloulou and Edéa.
\end{itemize}
\end{aretenir}

\coursec{Self- and Mutual Inductance; the L--R Circuit and Stored Energy}

In the previous section we saw how a generator converts mechanical work into electrical energy by forcing a flux change through a coil. But a coil does not need an external moving magnet to experience induction: \emph{its own} changing current changes the flux through itself, and two neighbouring coils can induce e.m.f.s in each other. These two ideas — \cle{self-inductance} and \cle{mutual inductance} — explain why the current in a circuit containing a coil cannot switch on instantly, why a spark jumps when you open a switch, and how much energy a current-carrying coil stores. They are also the working principle of transformers, which you will meet later.

\courssub{Self-Inductance}

\textbf{Observation.} Connect a coil in series with a lamp and a battery. When the switch is closed, the lamp brightens \emph{gradually} over a fraction of a second; with a plain resistor instead of the coil, it lights instantly. The coil is opposing the growth of its own current. Where does this opposition come from? When the current $I$ in the coil changes, the magnetic field it produces changes, so the flux through the coil itself changes. By Faraday's law, an e.m.f. is induced \emph{in the same coil}, and by Lenz's law it opposes the change: it is called a \cle{back e.m.f.}.

\begin{definition}[Self-inductance]
The \cle{self-inductance} $L$ of a coil is the constant of proportionality between the \textbf{flux linkage} through the coil and the current producing it:
\[
N\Phi = LI
\]
where $N$ is the number of turns and $\Phi$ the flux through one turn. The SI unit of $L$ is the \cle{henry} (H): a coil has an inductance of $1\ \mathrm{H}$ if a current of $1\ \mathrm{A}$ produces a flux linkage of $1\ \mathrm{Wb\,turn}$. Equivalently, $1\ \mathrm{H} = 1\ \mathrm{V\,s\,A^{-1}}$. By Faraday's and Lenz's laws, the induced (back) e.m.f. is
\[
e = -\frac{d(N\Phi)}{dt} = -L\,\frac{dI}{dt}.
\]
\end{definition}

The minus sign is Lenz's law: if $I$ is increasing, $e$ opposes the battery; if $I$ is decreasing, $e$ tries to maintain it. Note that $L$ depends only on the \emph{geometry} of the coil (number of turns, length, cross-section) and the \emph{core material} — an iron core multiplies $L$ enormously because of its large relative permeability.

\begin{propriete}[Self-inductance of a long solenoid]
For a long solenoid of length $\ell$, cross-sectional area $A$, with $N$ turns ($n = N/\ell$ turns per unit length), in air (or vacuum):
\[
L = \mu_0 n^2 A \ell = \frac{\mu_0 N^2 A}{\ell}.
\]
\end{propriete}

\textbf{Derivation (examinable).} A current $I$ in a long solenoid gives a uniform internal field $B = \mu_0 n I$. The flux through one turn is $\Phi = BA = \mu_0 n I A$, so the flux linkage is
\[
N\Phi = N\mu_0 n I A = \mu_0 n^2 A \ell\, I \quad (\text{since } N = n\ell).
\]
Comparing with $N\Phi = LI$ gives $L = \mu_0 n^2 A\ell$. This shows explicitly that $L$ grows with $N^2$, with $A$, and decreases with $\ell$.

\courssub{Mutual Inductance}

\textbf{Observation.} Wind two coils on the same soft-iron core (or place them coaxially close together). Pass a varying current $I_1$ through the first (the \emph{primary}): a galvanometer connected to the second (the \emph{secondary}) deflects, although the two coils are not electrically connected. Part of the flux produced by the primary links the secondary, and its variation induces an e.m.f. there. This is \cle{mutual induction}.

\begin{definition}[Mutual inductance]
The \cle{mutual inductance} $M$ between two coils is the constant of proportionality between the flux linkage in one coil and the current in the other:
\[
N_2\Phi_2 = MI_1 \qquad\text{and}\qquad N_1\Phi_1 = MI_2,
\]
where $\Phi_2$ is the flux through one turn of coil 2 due to the current $I_1$ in coil 1 (and vice versa). It can be shown (reciprocity theorem) that the same constant $M$ appears in both relations. The unit of $M$ is the \cle{henry} (H). The e.m.f. induced in the secondary is
\[
e_2 = -M\,\frac{dI_1}{dt}.
\]
\end{definition}

\begin{popfigure}
\begin{center}
\begin{tikzpicture}[scale=1.0]
% core
\draw[thick, popmuted] (-4.2,-0.9) rectangle (4.2,0.9);
\node[popmuted] at (0,0) {soft-iron core};
% primary coil (left)
\foreach \x in {-3.4,-3.0,-2.6,-2.2}{
  \draw[thick, popblue] (\x,0.9) arc[start angle=90, end angle=450, x radius=0.18, y radius=0.55];}
\draw[thick, popblue] (-3.4,0.9) -- (-3.4,1.7) -- (-1.2,1.7);
\draw[thick, popblue] (-2.2,0.9) .. controls (-2.0,1.3) .. (-1.6,1.7) -- (-1.2,1.7);
\node[circle, draw, popblue, inner sep=1.5pt] at (-0.9,1.7) {\scriptsize A};
\draw[thick, popblue] (-0.6,1.7) -- (0.2,1.7);
\node[popblue] at (-2.6,2.15) {primary, $N_1$ turns};
\draw[->, thick, popblue] (-3.9,1.7) -- (-3.45,1.7);
\node[popblue] at (-4.15,1.95) {$I_1$};
% secondary coil (right)
\foreach \x in {2.2,2.6,3.0,3.4}{
  \draw[thick, poporange] (\x,0.9) arc[start angle=90, end angle=450, x radius=0.18, y radius=0.55];}
\draw[thick, poporange] (2.2,0.9) -- (2.2,1.7) -- (1.4,1.7);
\draw[thick, poporange] (3.4,0.9) -- (3.4,1.7) -- (4.0,1.7);
\node[circle, draw, poporange, inner sep=1.5pt] at (1.1,1.7) {\scriptsize G};
\draw[thick, poporange] (1.4,1.7) -- (0.8,1.7);
\draw[thick, poporange] (4.0,1.7) -- (4.0,0.9);
\node[poporange] at (3.0,2.15) {secondary, $N_2$ turns};
% flux arrows
\draw[->, thick, popgreen] (-1.5,0.35) -- (0.9,0.35);
\draw[->, thick, popgreen] (-1.5,-0.35) -- (0.9,-0.35);
\node[popgreen] at (0.4,-0.68) {$\Phi_2$};
\end{tikzpicture}
\end{center}
\legende{Mutual inductance: the varying current $I_1$ in the primary creates a flux $\Phi_2$ linking the secondary; the galvanometer G detects the induced e.m.f. $e_2 = -M\,dI_1/dt$.}
\end{popfigure}

\begin{aretenir}[Meaning of $L$ and $M$]
$L$ and $M$ are \textbf{constants of proportionality between flux linkage and current}. They depend only on geometry (turns, sizes, separation) and on the core material — never on the actual value of the current. A shared iron core concentrates the flux and makes $M$ large; moving the coils apart reduces $M$.
\end{aretenir}

\courssub{The L--R Circuit: Growth of Current}

\begin{popfigure}
\begin{center}
\begin{circuitikz}[european, scale=1.0]
\draw (0,0) to[battery1, l={$E$}] (0,2.4)
      to[nos, l={$K$}] (2.2,2.4)
      to[R, l={$R$}] (4.6,2.4)
      to[L, l={$L$}] (4.6,0) -- (0,0);
\draw (2.2,2.4) -- (2.2,3.4) to[ammeter] (4.6,3.4) -- (4.6,2.4);
\draw (0,0) -- (-1.4,0) to[voltmeter] (-1.4,2.4) -- (0,2.4);
\end{circuitikz}
\end{center}
\legende{Series $L$--$R$ circuit: battery $E$, switch $K$, resistor $R$, inductor $L$. The ammeter records $I(t)$; the voltmeter monitors the supply voltage.}
\end{popfigure}

\textbf{Setting up the equation.} Close the switch $K$ at $t = 0$. At each instant, Kirchhoff's voltage law around the loop gives
\[
E = RI + L\frac{dI}{dt},
\]
since the inductor presents a back e.m.f. $L\,dI/dt$ opposing the battery. This is a first-order linear differential equation. Separating variables:
\[
\frac{dI}{E - RI} = \frac{dt}{L}
\;\Longrightarrow\;
-\frac{1}{R}\ln(E - RI) = \frac{t}{L} + C.
\]
With the initial condition $I(0) = 0$ (the current through an inductor cannot jump discontinuously), we find $C = -\frac{1}{R}\ln E$. Solving for $I$:

\begin{propriete}[Growth of current in an L--R circuit]
\[
I(t) = I_0\left(1 - e^{-t/\tau}\right), \qquad I_0 = \frac{E}{R}, \qquad \tau = \frac{L}{R}.
\]
$I_0$ is the final steady current (when $dI/dt = 0$, the inductor behaves like a plain wire) and $\tau$ is the \cle{time constant} of the circuit.
\end{propriete}

\textbf{Physical meaning of $\tau$.} After a time $t = \tau$, $I = I_0(1 - e^{-1}) \approx 0.63\,I_0$: the current has reached $63\%$ of its final value. After $t = 5\tau$, $I \approx 0.993\,I_0$ — the growth is practically complete. A large $L$ or a small $R$ makes the growth slow.

\courssub{Decay of Current}

Suppose the current has reached $I_0$ and the battery is then removed from the circuit (the switch connects the coil and resistor in a closed loop without the source).

\begin{popfigure}
\begin{center}
\begin{circuitikz}[european, scale=1.0]
\draw (0,0) to[R, l={$R$}] (3,0)
      to[L, l={$L$}] (3,-2) -- (0,-2) -- (0,0);
\draw (1.5,0) -- (1.5,0.8) to[ammeter] (3,0.8) -- (3,0);
\end{circuitikz}
\end{center}
\legende{Decay circuit: the battery is replaced by a short circuit; the inductor discharges through the resistor.}
\end{popfigure}

Kirchhoff's law now gives
\[
L\frac{dI}{dt} + RI = 0
\;\Longrightarrow\;
\frac{dI}{I} = -\frac{R}{L}\,dt
\;\Longrightarrow\;
I(t) = I_0\,e^{-t/\tau}.
\]
The inductor opposes the \emph{decrease}: it releases its stored energy to keep the current flowing, and the current dies away exponentially. After one time constant, $I = I_0 e^{-1} \approx 0.37\,I_0$.

\begin{popfigure}
\begin{center}
\begin{tikzpicture}
\begin{axis}[
  width=11cm, height=6.5cm,
  axis lines=left,
  xlabel={$t/\tau$}, ylabel={$I/I_0$},
  xmin=0, xmax=5.2, ymin=0, ymax=1.15,
  xtick={1}, ytick={0.63,1},
  yticklabels={$0.63\,I_0$,$I_0$},
  domain=0:5, samples=100, thick]
\addplot[popblue, very thick] {1-exp(-x)};
\addlegendentry{growth: $I_0(1-e^{-t/\tau})$}
\addplot[poporange, very thick] {exp(-x)};
\addlegendentry{decay: $I_0 e^{-t/\tau}$}
\addplot[popmuted, dashed] coordinates {(0,1) (5.2,1)};
\addplot[popmuted, dashed] coordinates {(1,0) (1,0.632)};
\addplot[popmuted, dashed] coordinates {(0,0.632) (1,0.632)};
\node[popmuted] at (axis cs:1,-0.09) {$\tau$};
\end{axis}
\end{tikzpicture}
\end{center}
\legende{Growth and decay of current in an $L$--$R$ circuit. At $t = \tau$ the growing current reaches $0.63\,I_0$ and the decaying current has fallen to $0.37\,I_0$.}
\end{popfigure}

\begin{attention}[The spark at switch opening]
If you simply \emph{open} a switch in a highly inductive circuit, $dI/dt$ becomes enormous and the back e.m.f. $e = -L\,dI/dt$ can reach thousands of volts: a spark jumps across the switch contacts. This is why engineers place a \textbf{protective diode} (flyback diode) in parallel with the coil of a relay: when the switch opens, the coil's current circulates through the diode and decays gently instead of sparking.
\end{attention}

\courssub{Energy Stored in an Inductor}

While the current grows, the inductor absorbs energy from the source and stores it in its magnetic field. The instantaneous power absorbed by the inductor is
\[
P = e_{\text{absorbed}}\, I = L\frac{dI}{dt}\, I.
\]
The total energy stored as the current rises from $0$ to $I$ is
\[
E = \int_0^I LI'\,dI' = \tfrac{1}{2}LI^2.
\]

\begin{propriete}[Energy stored in an inductor]
\[
E = \tfrac{1}{2}LI^2 \quad (\text{joules, with } L \text{ in H and } I \text{ in A}).
\]
This is the magnetic analogue of the energy $\tfrac{1}{2}CV^2$ stored in a capacitor's electric field. The proof by integration of power, as above, is examinable.
\end{propriete}

\begin{methode}[Solving L--R problems]
\begin{enumerate}
\item \textbf{Write Kirchhoff's equation} for the loop: $E = RI + L\,dI/dt$ (growth) or $RI + L\,dI/dt = 0$ (decay).
\item \textbf{Identify} the steady current $I_0 = E/R$ and the time constant $\tau = L/R$.
\item \textbf{Use the exponential solution}: $I = I_0(1 - e^{-t/\tau})$ for growth, $I = I_0 e^{-t/\tau}$ for decay; differentiate if the back e.m.f. is asked.
\item \textbf{Check the limits}: at $t = 0$ the growth current is $0$; as $t \to \infty$ it tends to $I_0$. If your formula violates these, it is wrong.
\end{enumerate}
\end{methode}

\begin{exoresolu}[Growth, time constant and stored energy]
A coil of inductance $L = 0.40\ \mathrm{H}$ and negligible resistance is connected in series with a resistor $R = 20\ \Omega$, a switch and a battery of e.m.f. $E = 12\ \mathrm{V}$. The switch is closed at $t = 0$.\\
(a) Find the time constant and the final steady current.\\
(b) Find the current at $t = 30\ \mathrm{ms}$ and the back e.m.f. at that instant.\\
(c) Find the energy stored in the coil when the current is steady.
\tcblower
\textbf{Solution.}\\
(a) $\tau = \dfrac{L}{R} = \dfrac{0.40}{20} = \SI{0.020}{s} = \SI{20}{ms}$;\qquad
$I_0 = \dfrac{E}{R} = \dfrac{12}{20} = \SI{0.60}{A}$.\\[4pt]
(b) $I = I_0\left(1 - e^{-t/\tau}\right) = 0.60\left(1 - e^{-30/20}\right) = 0.60(1 - e^{-1.5}) = 0.60(1 - 0.223) \approx \SI{0.47}{A}$.\\
Back e.m.f.: $e = L\,\dfrac{dI}{dt} = E - RI = 12 - 20 \times 0.47 \approx \SI{2.7}{V}$.\\[4pt]
(c) $E_{\text{stored}} = \tfrac{1}{2}LI_0^2 = \tfrac{1}{2} \times 0.40 \times (0.60)^2 = \SI{0.072}{J}$.
\end{exoresolu}

\begin{attention}[Common mistakes]
\begin{itemize}
\item The current through an inductor \textbf{cannot change instantaneously}: at $t = 0$ during growth, $I = 0$, \emph{not} $E/R$. Writing $I = E/R$ at $t = 0$ is the classic GCE error.
\item Do not confuse $\tau = L/R$ (inductor) with $\tau = RC$ (capacitor).
\item In $e = -L\,dI/dt$, the back e.m.f. is largest at $t = 0$ (equal to $E$) and falls to zero as the current stabilises — not the other way round.
\item Energy is $\tfrac{1}{2}LI^2$: forgetting the factor $\tfrac{1}{2}$ or squaring $L$ instead of $I$ are frequent slips.
\end{itemize}
\end{attention}

\begin{exercice}[Decay and graph reading]
A coil of inductance $0.10\ \mathrm{H}$ carries a steady current of $2.0\ \mathrm{A}$ in series with a $5.0\ \Omega$ resistor. The source is suddenly replaced by a short circuit, so the current decays through the resistor.\\
(a) Write the expression for $I(t)$ and give the time constant.\\
(b) Calculate the current at $t = 40\ \mathrm{ms}$.\\
(c) Calculate the energy stored initially, and state what becomes of it.
\end{exercice}

\begin{corrige}[Exercise 1]
(a) $I(t) = I_0 e^{-t/\tau}$ with $I_0 = 2.0\ \mathrm{A}$ and $\tau = L/R = 0.10/5.0 = \SI{0.020}{s}$.\\
(b) $I = 2.0\,e^{-40/20} = 2.0\,e^{-2} = 2.0 \times 0.135 \approx \SI{0.27}{A}$.\\
(c) $E = \tfrac{1}{2}LI_0^2 = \tfrac{1}{2} \times 0.10 \times (2.0)^2 = \SI{0.20}{J}$. This magnetic energy is dissipated as heat in the resistor as the current decays.
\end{corrige}

\begin{savaistu}[Inductors in everyday Cameroon]
The fluorescent lamp chokes used in many classrooms and homes are simply large inductors: they limit the current through the tube once it is lit, and they provide the high voltage spike (via $e = -L\,dI/dt$) needed to start the discharge. The same principle, mutual inductance this time, is at work in the transformers that step down ENEO's high-voltage lines to the $220\ \mathrm{V}$ supplied to your house.
\end{savaistu}

\newpage

\coursec{Integration activity}

\courssub{Aims of the activity}

This integration activity closes the chapter on electromagnetic induction by placing you in the role of a young engineer. Working in groups, you will design a simple bicycle dynamo capable of charging mobile phones in a village without grid electricity. By the end of the activity you should be able to:

\begin{itemize}
\item apply \cle{Faraday's law} and the rotating-coil formula $\mathcal{E}_0 = NBA\omega$ to a concrete design problem;
\item interpret the sinusoidal output of a simple \cle{AC generator} and relate its frequency to the rotation speed;
\item use \cle{Lenz's law} to explain the back-torque that a loaded generator exerts on the cyclist;
\item compute the \cle{time constant} $\tau = L/R$ of an $L$--$R$ circuit and the \cle{energy stored in an inductor} $W = \tfrac{1}{2}LI^2$;
\item argue, quantify, justify and communicate a technical solution on a poster, as expected at the GCE Advanced Level.
\end{itemize}

\courssub{Situation}

\begin{experience}[The Nkongsamba phone-charging project]
The village of Bakassa, near Nkongsamba in the Littoral Region, is not yet connected to the national grid. The nearest charging point is a commercial generator $8\ \mathrm{km}$ away, and villagers pay for each phone charge. The science club of the local secondary school decides to build a \textbf{bicycle-powered charging station}: a cyclist pedals, the wheel drives a small coil which spins between the poles of a permanent magnet, and the induced e.m.f. charges phones through a simple regulating circuit.

The club has salvaged the following components:
\begin{itemize}
\item a permanent magnet giving an approximately uniform field $B = 0.20\ \mathrm{T}$ between its poles;
\item copper wire, to be wound into a flat rectangular coil of area $A = 20\ \mathrm{cm^2} = 2.0\times 10^{-3}\ \mathrm{m^2}$;
\item a belt drive allowing a fit cyclist to spin the coil at up to $300$ revolutions per minute;
\item a rectifier and regulating circuit whose total resistance, together with the coil, is $R = 10\ \Omega$; the coil itself has a self-inductance $L = 50\ \mathrm{mH}$.
\end{itemize}

The phone battery requires a rectified supply whose \textbf{peak e.m.f. should be about $6\ \mathrm{V}$}. The club must decide how many turns $N$ to wind, predict the output waveform, understand why pedalling ``bites back'' when a phone is connected, and check the behaviour of the coil in the DC part of the circuit. You are the consulting physicists: produce the calculations and a poster justifying every step.
\end{experience}

\courssub{Tasks}

\begin{enumerate}
\item \textbf{Choosing the number of turns.} Convert the rotation speed into angular velocity $\omega$ in $\mathrm{rad\,s^{-1}}$. Using the peak e.m.f. of a rotating coil, $\mathcal{E}_0 = NBA\omega$, determine the number of turns $N$ needed for $\mathcal{E}_0 \approx 6\ \mathrm{V}$. Round to a realistic whole number of turns.
\item \textbf{Output waveform.} Write the instantaneous e.m.f. $\mathcal{E}(t)$, sketch $\mathcal{E}$ against time for two full cycles, and state the frequency $f$ and period $T$ of the output. Mark the instants when the coil plane is parallel to the field.
\item \textbf{Lenz's law and the cyclist.} When no phone is connected, the coil circuit is open and pedalling is easy. When a phone draws current, pedalling becomes noticeably harder. Using Lenz's law, explain the origin of the extra resistive torque, and state where the extra energy goes.
\item \textbf{The coil as an inductor.} In the DC (rectified) part of the circuit, the coil of inductance $L = 50\ \mathrm{mH}$ is in series with the total resistance $R = 10\ \Omega$.
\begin{enumerate}
\item Compute the time constant $\tau$ of this $L$--$R$ circuit and interpret its meaning for the growth of the current.
\item Compute the magnetic energy stored in the coil when the steady current is $I = 0.5\ \mathrm{A}$.
\end{enumerate}
\item \textbf{Improving the design.} Propose \emph{one} modification that would increase the peak e.m.f., and justify it using the formula $\mathcal{E}_0 = NBA\omega$. Discuss briefly one practical drawback of your proposal.
\item \textbf{Poster.} Assemble your results (calculations, graph, explanations) on a poster and present them to the class in five minutes.
\end{enumerate}

\courssub{Model answer}

\textbf{Task 1 — Number of turns.}\par
The angular velocity corresponding to $300$ revolutions per minute is
\[
\omega = \frac{300 \times 2\pi}{60} = 10\pi \approx 31.4\ \mathrm{rad\,s^{-1}}.
\]
From $\mathcal{E}_0 = NBA\omega$,
\[
N = \frac{\mathcal{E}_0}{BA\omega} = \frac{6}{0.20 \times 2.0\times 10^{-3} \times 31.4} \approx 4.8\times 10^{2}.
\]
\begin{aretenir}[Design choice]
Wind about $\cle{N \approx 480}$ turns (say $500$ turns to allow for losses). Fewer turns would leave the peak e.m.f. below the $6\ \mathrm{V}$ target.
\end{aretenir}

\textbf{Task 2 — Output waveform.}\par
Take $t = 0$ when the coil plane is perpendicular to $\vec B$, so that the flux through one turn is $\Phi = BA\cos\omega t$ and the \cle{flux linkage} is $N\Phi = NBA\cos\omega t$. Faraday's law then gives
\[
\mathcal{E}(t) = -\frac{d(N\Phi)}{dt} = NBA\omega\sin\omega t = \mathcal{E}_0\sin\omega t,
\qquad \mathcal{E}_0 \approx 6\ \mathrm{V}.
\]
The frequency is
\[
f = \frac{\omega}{2\pi} = \frac{300}{60} = 5.0\ \mathrm{Hz}, \qquad T = \frac{1}{f} = 0.20\ \mathrm{s}.
\]
The e.m.f. is zero when the coil plane is perpendicular to $\vec B$ (flux maximum, but momentarily not changing) and maximum when the coil plane is parallel to $\vec B$ (flux zero but changing fastest). On the graph below, the instants $t = 0.05\ \mathrm{s}$, $0.15\ \mathrm{s}$, $0.25\ \mathrm{s}$ and $0.35\ \mathrm{s}$ (peaks and troughs) correspond to the coil plane being parallel to $\vec B$.

\begin{popfigure}
\begin{center}
\begin{tikzpicture}
\begin{axis}[
  width=11cm, height=5.5cm,
  axis lines=middle,
  xlabel={$t$ (s)}, ylabel={$\mathcal{E}$ (V)},
  xmin=0, xmax=0.42, ymin=-7.5, ymax=7.5,
  xtick={0,0.1,0.2,0.3,0.4},
  ytick={-6,0,6},
  domain=0:0.4, samples=200,
  every axis x label/.style={at={(ticklabel* cs:1.02)}, anchor=west},
  every axis y label/.style={at={(ticklabel* cs:1.02)}, anchor=south},
]
\addplot[popblue, thick] {6*sin(deg(2*pi*5*x))};
\addplot[popmuted, dashed] coordinates {(0,6) (0.4,6)};
\addplot[popmuted, dashed] coordinates {(0,-6) (0.4,-6)};
\node[popblue, anchor=south west] at (axis cs:0.055,6.2) {$\mathcal{E}_0 = 6\ \mathrm{V}$};
\node[popmuted, anchor=north west] at (axis cs:0.16,-6.2) {$-\mathcal{E}_0$};
\addplot[poporange, mark=*, mark size=1.5pt, only marks] coordinates {(0.05,6) (0.15,-6) (0.25,6) (0.35,-6)};
\node[poporange, anchor=south, align=center, text width=3.2cm] at (axis cs:0.25,6.4) {coil plane parallel to $\vec B$};
\end{axis}
\end{tikzpicture}
\end{center}
\legende{Sinusoidal output of the dynamo: $\mathcal{E}(t) = 6\sin(10\pi t)$, frequency $5.0\ \mathrm{Hz}$, period $0.20\ \mathrm{s}$. The orange dots mark the instants when the coil plane is parallel to the field.}
\end{popfigure}

\textbf{Task 3 — Lenz's law and the back-torque.}\par
With the circuit open, no current flows: the coil experiences only friction, so pedalling is easy. When a phone is connected, the induced e.m.f. drives a current $I$ around the coil. By \cle{Lenz's law}, this induced current must oppose the change that produces it: the magnetic forces on the current-carrying sides of the coil create a \textbf{braking torque} opposing the rotation. The cyclist must therefore work against this torque. Energy is conserved: the extra muscular work done by the cyclist is transferred to the phone battery (and partly dissipated as heat in $R$). The harder the phone charges (larger $I$), the larger the braking torque — the dynamo never gives energy for free.

\textbf{Task 4 — The coil as an inductor.}\par
\textbf{(a)} The time constant of the $L$--$R$ circuit is
\[
\tau = \frac{L}{R} = \frac{50\times 10^{-3}}{10} = 5.0\times 10^{-3}\ \mathrm{s} = 5.0\ \mathrm{ms}.
\]
When the rectified supply is switched on, the current grows according to $i(t) = I_0\left(1 - e^{-t/\tau}\right)$: it reaches about $63\ \%$ of its final value after $5\ \mathrm{ms}$, and is practically steady after $5\tau \approx 25\ \mathrm{ms}$. The inductor smooths the charging current — a useful side-effect for the phone.

\textbf{(b)} At $I = 0.5\ \mathrm{A}$, the magnetic energy stored in the coil is
\[
W = \frac{1}{2}LI^2 = \frac{1}{2}\times 50\times 10^{-3}\times (0.5)^2 = 6.25\times 10^{-3}\ \mathrm{J} \approx 6.3\ \mathrm{mJ}.
\]

\textbf{Task 5 — Improving the design.}\par
Since $\mathcal{E}_0 = NBA\omega$, the peak e.m.f. can be raised by increasing any one of $N$, $B$, $A$ or $\omega$. A practical suggestion: \textbf{fit a stronger neodymium magnet} (from an old loudspeaker or hard-disk drive), raising $B$ from $0.20\ \mathrm{T}$ to about $0.40\ \mathrm{T}$, which would double $\mathcal{E}_0$ to about $12\ \mathrm{V}$ without adding turns. Drawback: stronger magnets are costlier, and the larger back-torque (Lenz's law) makes pedalling harder when the phone is connected — the energy must always come from the cyclist. (Equally acceptable: doubling $N$ to about $1000$ turns, at the cost of a heavier coil and a larger resistance.)

\begin{attention}[Common mistakes to avoid on the poster]
\begin{itemize}
\item Forgetting to convert revolutions per minute into $\mathrm{rad\,s^{-1}}$: $\omega = 2\pi f$ with $f$ in hertz, not $300\ \mathrm{rad\,s^{-1}}$!
\item Mixing up $\mathrm{cm^2}$ and $\mathrm{m^2}$: $20\ \mathrm{cm^2} = 2.0\times 10^{-3}\ \mathrm{m^2}$.
\item Writing $\tau = R/L$ instead of $\tau = L/R$ — check with units: henry per ohm gives seconds.
\item Believing the dynamo ``creates'' energy: Lenz's law guarantees that every joule delivered to the phone is paid for by the cyclist's muscles.
\end{itemize}
\end{attention}

\begin{methode}[Assessment guide for the poster]
\begin{enumerate}
\item Correct use of $\mathcal{E}_0 = NBA\omega$ with unit conversions (2 marks).
\item Correct sinusoidal sketch with labelled amplitude, period and frequency (2 marks).
\item Lenz's-law explanation of the braking torque, linked to energy conservation (2 marks).
\item Correct $\tau = L/R = 5.0\ \mathrm{ms}$ and $W = \tfrac12 LI^2 = 6.3\ \mathrm{mJ}$ with units (2 marks).
\item A justified, realistic improvement and its drawback; clarity of presentation (2 marks).
\end{enumerate}
\end{methode}

\begin{savaistu}[Did you know?]
Bicycle dynamos of this kind have powered lamps in Cameroonian towns and villages for decades, and hand-cranked or pedal chargers remain vital in off-grid areas of the Far North and East Regions. The same physics — a coil spinning in a magnetic field — is scaled up enormously at the Song Loulou and Edea hydroelectric stations on the Sanaga river, where turbines spin huge rotors to supply the national grid. From a bicycle in Bakassa to a dam on the Sanaga, it is always Faraday's law at work.
\end{savaistu}

\newpage

\coursec{Methods and worked examples}

\courssub{Methods and Worked Examples}

\begin{methode}[Calculating Magnetic Flux and Applying Faraday's Law]
\begin{enumerate}
\item \textbf{Identify the surface and its orientation.} Define the area vector $\vec{A}$ (magnitude $A$, direction normal to the plane). Determine the angle $\theta$ between $\vec{B}$ and $\vec{A}$.
\item \textbf{Write the magnetic flux.} For a uniform field: $\Phi = BA\cos\theta$. For a coil of $N$ turns: $\Phi_{\text{total}} = N\Phi = NBA\cos\theta$.
\item \textbf{Determine the rate of change.} Identify what varies with time: $B$, $A$, or $\theta$ (or a combination). Compute $\dfrac{d\Phi}{dt}$ or $\dfrac{\Delta\Phi}{\Delta t}$ for average values.
\item \textbf{Apply Faraday's Law.} The magnitude of the induced e.m.f. is $\mathcal{E} = \left|\dfrac{d\Phi_{\text{total}}}{dt}\right| = \left|N\dfrac{d\Phi}{dt}\right|$.
\item \textbf{Check units.} $1\ \mathrm{Wb} = 1\ \mathrm{T\,m^2}$, $1\ \mathrm{V} = 1\ \mathrm{Wb\,s^{-1}}$.
\end{enumerate}
\end{methode}

\begin{exoresolu}[Magnetic Flux through a Rotating Coil]
A circular coil of $200$ turns has a radius of $5.0\ \mathrm{cm}$. It is placed in a uniform magnetic field of flux density $0.80\ \mathrm{T}$. Calculate the magnetic flux linkage through the coil when: (a) the plane of the coil is perpendicular to the field; (b) the plane of the coil is at $30^\circ$ to the field; (c) the plane of the coil is parallel to the field. If the coil is rotated from position (a) to position (c) in $0.25\ \mathrm{s}$, calculate the average e.m.f. induced in the coil.
\tcblower
\textbf{Solution.}
\begin{itemize}
\item \textbf{Data:} $N = 200$, $r = 5.0\ \mathrm{cm} = 0.050\ \mathrm{m}$, $B = 0.80\ \mathrm{T}$, $\Delta t = 0.25\ \mathrm{s}$.
\item \textbf{Area:} $A = \pi r^2 = \pi (0.050)^2 = 7.854 \times 10^{-3}\ \mathrm{m^2}$.
\item \textbf{Angle definition:} $\theta$ is the angle between $\vec{B}$ and the \emph{normal} to the plane (area vector $\vec{A}$).
\end{itemize}
\textbf{(a) Plane $\perp$ field $\Rightarrow$ Normal $\parallel$ field $\Rightarrow \theta = 0^\circ$.}
\[
\Phi = BA\cos 0^\circ = (0.80)(7.854\times 10^{-3})(1) = 6.283 \times 10^{-3}\ \mathrm{Wb}
\]
Flux linkage $= N\Phi = 200 \times 6.283\times 10^{-3} = \boxed{1.26\ \mathrm{Wb}\text{ (or Vs)}}$.

\textbf{(b) Plane at $30^\circ$ to field $\Rightarrow$ Normal at $60^\circ$ to field $\Rightarrow \theta = 60^\circ$.}
\[
\Phi = BA\cos 60^\circ = (0.80)(7.854\times 10^{-3})(0.5) = 3.142 \times 10^{-3}\ \mathrm{Wb}
\]
Flux linkage $= 200 \times 3.142\times 10^{-3} = \boxed{0.628\ \mathrm{Wb}}$.

\textbf{(c) Plane $\parallel$ field $\Rightarrow$ Normal $\perp$ field $\Rightarrow \theta = 90^\circ$.}
\[
\Phi = BA\cos 90^\circ = 0 \quad \Rightarrow \quad \text{Flux linkage} = \boxed{0\ \mathrm{Wb}}.
\]

\textbf{Average e.m.f. from (a) to (c):}
\[
\mathcal{E}_{\text{avg}} = \left| N \frac{\Delta \Phi}{\Delta t} \right| = \left| 200 \times \frac{0 - 6.283\times 10^{-3}}{0.25} \right| = \frac{1.2566}{0.25} = \boxed{5.03\ \mathrm{V}}.
\]
\end{exoresolu}

\begin{methode}[Applying Lenz's Law and Motional E.M.F. ($\mathcal{E} = Blv$)]
\begin{enumerate}
\item \textbf{Draw the situation.} Show $\vec{B}$, the conductor (length $l$), velocity $\vec{v}$, and the circuit loop (real or imaginary).
\item \textbf{Determine the change in flux.} Is the area of the loop increasing or decreasing? Is the flux entering or leaving the page?
\item \textbf{Apply Lenz's Law:} \emph{The induced current flows in a direction such that its magnetic field opposes the change in flux that produced it.}
\item \textbf{Find current direction.} Use the Right-Hand Grip Rule: thumb along the induced $\vec{B}_{\text{ind}}$ (opposing the change), fingers curl in direction of $I_{\text{ind}}$.
\item \textbf{Calculate magnitude.} For a straight rod moving perpendicular to $\vec{B}$: $\mathcal{E} = Blv$. If $\vec{v}$ is at angle $\theta$ to $\vec{B}$: $\mathcal{E} = Blv\sin\theta$ (component of $v \perp B$). If rod is at angle $\phi$ to $\vec{v}$: $\mathcal{E} = Blv\sin\phi$ (effective length $\perp v$). General: $\mathcal{E} = (\vec{v} \times \vec{B}) \cdot \vec{l}$.
\item \textbf{Polarity.} The end of the rod where positive charge accumulates is at higher potential. Use $\vec{F}_q = q(\vec{v} \times \vec{B})$ on free electrons ($q=-e$) to find which end becomes negative.
\end{enumerate}
\end{methode}

\begin{exoresolu}[Motional E.M.F. in a Moving Rod on Rails]
A horizontal rod $PQ$ of length $0.40\ \mathrm{m}$ and resistance $2.0\ \Omega$ slides at a constant speed of $5.0\ \mathrm{m\,s^{-1}}$ along two parallel vertical rails separated by $0.40\ \mathrm{m}$. The rails are connected at the bottom by a resistor $R = 3.0\ \Omega$. The whole apparatus is in a uniform vertical magnetic field of flux density $0.50\ \mathrm{T}$ directed \emph{downwards}. (a) Calculate the induced e.m.f. in the rod. (b) Find the magnitude and direction of the induced current. (c) Calculate the magnetic force acting on the rod due to this current. (d) What external force is required to keep the rod moving at constant speed? (e) Calculate the power dissipated in the circuit and the rate of work done by the external force.
\tcblower
\textbf{Solution.}
\begin{itemize}
\item \textbf{Setup:} Rod moves horizontally (say to the right). $\vec{B}$ downwards. Rails vertical $\Rightarrow$ loop area increases horizontally.
\end{itemize}
\textbf{(a) Induced e.m.f.:} $\vec{v} \perp \vec{B}$, rod $\perp \vec{v}$ and $\perp \vec{B}$.
\[
\mathcal{E} = Blv = (0.50)(0.40)(5.0) = \boxed{1.0\ \mathrm{V}}.
\]

\textbf{(b) Current magnitude and direction.}
Total resistance $R_{\text{tot}} = 2.0 + 3.0 = 5.0\ \Omega$.
\[
I = \frac{\mathcal{E}}{R_{\text{tot}}} = \frac{1.0}{5.0} = \boxed{0.20\ \mathrm{A}}.
\]
\textit{Direction (Lenz's Law):} Rod moves right $\Rightarrow$ Area increases $\Rightarrow$ Flux \emph{downwards} increases. Induced $\vec{B}_{\text{ind}}$ must be \emph{upwards} to oppose. Right-hand grip rule: Thumb up $\Rightarrow$ fingers curl \emph{anticlockwise} (viewed from above). Current flows $P \to Q$ along the rod (if $P$ is left, $Q$ right), down the right rail, through $R$, up the left rail.

\textit{Polarity (Lorentz force on electrons):} $\vec{F} = -e(\vec{v} \times \vec{B})$. Take coordinates: $x$ right, $y$ up, $z$ out of page. $\vec{v}=v\hat{i}$, $\vec{B}=-B\hat{j}$, so $\vec{v}\times\vec{B} = -vB(\hat{i}\times\hat{j}) = -vB\hat{k}$ (into page). The force on electrons ($q=-e$) is $+evB\,\hat{k}$ (out of page): electrons accumulate at the end facing out of the page, which becomes negative; the opposite end is positive. Conventional current inside the rod flows from the negative end to the positive end, i.e.\ \textbf{in the direction of $\vec{v}\times\vec{B}$}.

\textbf{(c) Magnetic force on rod:} $\vec{F}_B = I \vec{l} \times \vec{B}$. $I=0.20\ \mathrm{A}$, $\vec{l}$ along the rod, $\vec{B}$ down. The cross product gives a force directed \emph{opposite to the velocity} (as Lenz's law requires).
Magnitude: $F_B = BIl = (0.50)(0.20)(0.40) = \boxed{0.040\ \mathrm{N}}$ \emph{opposing the motion}.

\textbf{(d) External force:} Constant velocity $\Rightarrow$ Net force $=0$. $\vec{F}_{\text{ext}} + \vec{F}_B = 0 \Rightarrow \vec{F}_{\text{ext}} = -\vec{F}_B$.
\[
F_{\text{ext}} = \boxed{0.040\ \mathrm{N}} \text{ in the direction of motion}.
\]

\textbf{(e) Power balance.}
Power dissipated (Joule heating): $P_{\text{diss}} = I^2 R_{\text{tot}} = (0.20)^2(5.0) = \boxed{0.20\ \mathrm{W}}$.
Rate of work by external force: $P_{\text{ext}} = F_{\text{ext}} v = (0.040)(5.0) = \boxed{0.20\ \mathrm{W}}$.
\textit{Check:} $P_{\text{ext}} = P_{\text{diss}}$. Mechanical work is converted to electrical energy then heat. Energy conserved.
\end{exoresolu}

\begin{methode}[Analysing AC and DC Generators]
\begin{enumerate}
\item \textbf{Visualise the rotation.} A coil of $N$ turns, area $A$, rotates at constant angular velocity $\omega$ in uniform $\vec{B}$.
\item \textbf{Flux linkage:} $\Phi_N = NBA\cos\theta = NBA\cos(\omega t)$ (assuming $\theta=0$ at $t=0$, flux max).
\item \textbf{Induced e.m.f. (Faraday):} $\mathcal{E} = -\dfrac{d\Phi_N}{dt} = NBA\omega\sin(\omega t) = \mathcal{E}_0\sin(\omega t)$.
\item \textbf{Peak e.m.f.:} $\mathcal{E}_0 = NBA\omega = 2\pi f NBA$.
\item \textbf{RMS values:} $\mathcal{E}_{\text{rms}} = \mathcal{E}_0/\sqrt{2}$, $I_{\text{rms}} = I_0/\sqrt{2}$ (for purely resistive load).
\item \textbf{DC Generator (Commutator):} Split-ring commutator reverses connections every half-turn. Output is $|\mathcal{E}_0\sin(\omega t)|$ (pulsating DC). Brushes contact when coil plane $\parallel \vec{B}$ (flux zero, rate of change max).
\item \textbf{AC Generator (Slip rings):} Continuous rings $\Rightarrow$ sinusoidal AC output.
\end{enumerate}
\end{methode}

\begin{exoresolu}[Design Parameters of an AC Generator]
An AC generator consists of a rectangular coil of $500$ turns, each of area $2.5 \times 10^{-2}\ \mathrm{m^2}$, rotating at $50\ \mathrm{rev\,s^{-1}}$ in a uniform magnetic field of flux density $0.40\ \mathrm{T}$. The coil has a total resistance of $10\ \Omega$ and is connected to a purely resistive load of $40\ \Omega$. (a) Calculate the peak e.m.f. and the r.m.s. e.m.f. generated. (b) Write the equation for the instantaneous e.m.f. assuming $\mathcal{E}=0$ at $t=0$ and increasing positively. (c) Calculate the r.m.s. current in the circuit. (d) Determine the average power dissipated in the load resistor. (e) If the generator runs at the same speed but the magnetic field is doubled, what happens to the peak e.m.f. and the r.m.s. current?
\tcblower
\textbf{Solution.}
\textbf{Data:} $N=500$, $A=2.5\times 10^{-2}\ \mathrm{m^2}$, $f=50\ \mathrm{Hz}$, $B=0.40\ \mathrm{T}$, $R_{\text{coil}}=10\ \Omega$, $R_{\text{load}}=40\ \Omega$.
$\omega = 2\pi f = 100\pi\ \mathrm{rad\,s^{-1}}$.

\textbf{(a) Peak and RMS e.m.f.:}
\[
\mathcal{E}_0 = NBA\omega = 500 \times 0.40 \times (2.5\times 10^{-2}) \times (100\pi) = 500\pi \approx \boxed{1570\ \mathrm{V}}.
\]
\[
\mathcal{E}_{\text{rms}} = \frac{\mathcal{E}_0}{\sqrt{2}} = \frac{1570}{\sqrt{2}} \approx \boxed{1110\ \mathrm{V}}.
\]

\textbf{(b) Instantaneous e.m.f.:} $\mathcal{E}=0$ at $t=0$, increasing $\Rightarrow$ sine function.
\[
\mathcal{E}(t) = \mathcal{E}_0 \sin(\omega t) = \boxed{1570 \sin(100\pi t)\ \mathrm{V}}.
\]

\textbf{(c) RMS current:} Total resistance $R_T = 10 + 40 = 50\ \Omega$.
\[
I_{\text{rms}} = \frac{\mathcal{E}_{\text{rms}}}{R_T} = \frac{1110}{50} = \boxed{22.2\ \mathrm{A}}.
\]

\textbf{(d) Average power in load:} $P_{\text{avg}} = I_{\text{rms}}^2 R_{\text{load}} = (22.2)^2 \times 40 = 492.8 \times 40 \approx \boxed{19.7\ \mathrm{kW}}$.

\textbf{(e) Effect of doubling $B$:} $\mathcal{E}_0 \propto B \Rightarrow \mathcal{E}_0' = 2\mathcal{E}_0 = \boxed{3140\ \mathrm{V}}$.
$I_{\text{rms}} \propto \mathcal{E}_{\text{rms}} \propto B \Rightarrow I_{\text{rms}}' = 2I_{\text{rms}} = \boxed{44.4\ \mathrm{A}}$.
\end{exoresolu}

\begin{methode}[Self-Inductance and the L-R Circuit (Growth/Decay)]
\begin{enumerate}
\item \textbf{Define Self-Inductance $L$.} $\Phi_N = LI$ (flux linkage proportional to current). Unit: Henry (H) = $\mathrm{V\,s\,A^{-1}}$. For a long solenoid: $L = \mu_0 n^2 A l = \mu_0 N^2 A / l$.
\item \textbf{Back e.m.f.:} $\mathcal{E}_L = -L \dfrac{dI}{dt}$. Opposes the change in current (Lenz's law applied to own field).
\item \textbf{L-R Circuit Equation (Growth):} Switch closed at $t=0$. $E = IR + L\dfrac{dI}{dt}$.
\item \textbf{Solution (Growth):} $I(t) = \dfrac{E}{R}\left(1 - e^{-t/\tau}\right)$ where $\tau = \dfrac{L}{R}$ (time constant).
\item \textbf{Key instants:} $t=0$: $I=0$, $V_L=E$, $V_R=0$. $t\to\infty$: $I=E/R$, $V_L=0$, $V_R=E$. $t=\tau$: $I = 0.632(E/R)$.
\item \textbf{Decay:} Switch opened (or source shorted). $0 = IR + L\dfrac{dI}{dt} \Rightarrow I(t) = I_0 e^{-t/\tau}$. $t=\tau$: $I = 0.368 I_0$.
\item \textbf{Energy Stored:} $U = \frac{1}{2}LI^2$. Derived from $P = \mathcal{E}_L I = L I \frac{dI}{dt} = \frac{d}{dt}(\frac{1}{2}LI^2)$.
\end{enumerate}
\end{methode}

\begin{exoresolu}[Transient Analysis of an L-R Circuit]
A coil of inductance $L = 2.0\ \mathrm{H}$ and resistance $R = 10\ \Omega$ is connected to a battery of e.m.f. $E = 12\ \mathrm{V}$ and negligible internal resistance. (a) Calculate the time constant of the circuit. (b) Determine the initial rate of increase of current ($\frac{dI}{dt}$ at $t=0$). (c) Find the current and the potential difference across the inductor at $t = 0.10\ \mathrm{s}$. (d) Calculate the energy stored in the magnetic field of the inductor when the current reaches its steady-state value. (e) The battery is suddenly disconnected and the coil is short-circuited (connected across a wire of zero resistance). Write the equation for the decaying current and find the time for the current to fall to $1.0\ \mathrm{A}$.
\tcblower
\textbf{Solution.}
\textbf{Data:} $L=2.0\ \mathrm{H}$, $R=10\ \Omega$, $E=12\ \mathrm{V}$.

\textbf{(a) Time constant:}
\[
\tau = \frac{L}{R} = \frac{2.0}{10} = \boxed{0.20\ \mathrm{s}}.
\]

\textbf{(b) Initial rate of increase:} At $t=0$, $I=0$, so $V_R=0$. Loop rule: $E = V_L = L\frac{dI}{dt}$.
\[
\left.\frac{dI}{dt}\right|_{t=0} = \frac{E}{L} = \frac{12}{2.0} = \boxed{6.0\ \mathrm{A\,s^{-1}}}.
\]

\textbf{(c) At $t = 0.10\ \mathrm{s}$ ($t = 0.5\tau$):}
Steady current $I_0 = E/R = 12/10 = 1.2\ \mathrm{A}$.
\[
I(t) = I_0(1 - e^{-t/\tau}) = 1.2(1 - e^{-0.5}) = 1.2(1 - 0.6065) = \boxed{0.472\ \mathrm{A}}.
\]
P.D. across inductor: $V_L = L\frac{dI}{dt} = E - IR = 12 - (0.472 \times 10) = 12 - 4.72 = \boxed{7.28\ \mathrm{V}}$.
(Check: $V_L = E e^{-t/\tau} = 12 e^{-0.5} = 7.28\ \mathrm{V}$).

\textbf{(d) Energy stored at steady state ($I=I_0=1.2\ \mathrm{A}$):}
\[
U = \frac{1}{2} L I_0^2 = \frac{1}{2} \times 2.0 \times (1.2)^2 = 1.0 \times 1.44 = \boxed{1.44\ \mathrm{J}}.
\]

\textbf{(e) Decay (short-circuited):} Initial current $I_0 = 1.2\ \mathrm{A}$. Equation: $I(t) = I_0 e^{-t/\tau}$.
Find $t$ when $I = 1.0\ \mathrm{A}$.
\[
1.0 = 1.2 e^{-t/0.20} \implies e^{-t/0.20} = \frac{1.0}{1.2} = \frac{5}{6} \approx 0.8333.
\]
\[
-\frac{t}{0.20} = \ln(0.8333) \approx -0.1823 \implies t = 0.20 \times 0.1823 = \boxed{0.0365\ \mathrm{s}} \ (36.5\ \mathrm{ms}).
\]
\end{exoresolu}

\begin{methode}[Mutual Inductance and Transformer Principles]
\begin{enumerate}
\item \textbf{Define Mutual Inductance $M$.} Coil 1 (primary) carries $I_1$, creates flux $\Phi_{21}$ linking coil 2 (secondary). $\Phi_{21} \propto I_1 \Rightarrow N_2\Phi_{21} = M I_1$. Unit: Henry (H). $M$ depends on geometry, distance, orientation, and core material ($\mu_r$).
\item \textbf{Reciprocity Theorem:} $M_{12} = M_{21} = M$. $N_1\Phi_{12} = M I_2$.
\item \textbf{Induced e.m.f. in secondary:} $\mathcal{E}_2 = -M \dfrac{dI_1}{dt}$. Magnitude: $|\mathcal{E}_2| = M \left|\dfrac{dI_1}{dt}\right|$.
\item \textbf{Coupling Coefficient $k$:} $M = k\sqrt{L_1 L_2}$, where $0 \le k \le 1$. $k=1$ for perfect coupling (ideal transformer).
\item \textbf{Ideal Transformer (No losses, $k=1$, infinite $L$):} $\dfrac{V_s}{V_p} = \dfrac{N_s}{N_p} = \dfrac{I_p}{I_s}$. Power in = Power out: $V_p I_p = V_s I_s$.
\item \textbf{Step-up/Step-down:} $N_s > N_p \Rightarrow$ Step-up (Voltage increases, Current decreases). $N_s < N_p \Rightarrow$ Step-down.
\end{enumerate}
\end{methode}

\begin{exoresolu}[Mutual Inductance and Transformer Calculations]
Two coils are wound on a common ferromagnetic core. The primary coil has $N_p = 2000$ turns and inductance $L_p = 4.0\ \mathrm{H}$. The secondary coil has $N_s = 500$ turns. The coupling coefficient is $k = 0.95$. (a) Calculate the self-inductance of the secondary coil $L_s$. (b) Calculate the mutual inductance $M$. (c) An alternating current $I_p = 2.0 \sin(100\pi t)\ \mathrm{A}$ flows in the primary. Determine the amplitude of the e.m.f. induced in the secondary. (d) If the secondary is connected to a resistive load of $10\ \Omega$, calculate the r.m.s. current in the secondary, assuming the induced e.m.f. of part (c) acts as the source for the secondary circuit. (e) Explain why the primary current would increase if the secondary load resistance is decreased.
\tcblower
\textbf{Solution.}
\textbf{(a) Self-inductance $L_s$:} For coils on the same core, $L \propto N^2$ (assuming same geometry, length, area, $\mu$).
\[
\frac{L_s}{L_p} = \left(\frac{N_s}{N_p}\right)^2 \implies L_s = L_p \left(\frac{N_s}{N_p}\right)^2 = 4.0 \times \left(\frac{500}{2000}\right)^2 = 4.0 \times \frac{1}{16} = \boxed{0.25\ \mathrm{H}}.
\]

\textbf{(b) Mutual Inductance $M$:}
\[
M = k\sqrt{L_p L_s} = 0.95 \times \sqrt{4.0 \times 0.25} = 0.95 \times \sqrt{1.0} = \boxed{0.95\ \mathrm{H}}.
\]
(Note: For perfect coupling $k=1$, $M = \sqrt{L_p L_s} = 1.0\ \mathrm{H}$. Also $M = k \frac{N_s}{N_p} L_p = 0.95 \times \frac{1}{4} \times 4.0 = 0.95\ \mathrm{H}$. Consistent.)

\textbf{(c) Induced e.m.f. amplitude in secondary:}
$I_p = I_0 \sin(\omega t)$ with $I_0 = 2.0\ \mathrm{A}$, $\omega = 100\pi\ \mathrm{rad\,s^{-1}}$.
$\frac{dI_p}{dt} = I_0 \omega \cos(\omega t)$. Max magnitude $= I_0 \omega$.
\[
\mathcal{E}_{s0} = M \left|\frac{dI_p}{dt}\right|_{\text{max}} = M I_0 \omega = 0.95 \times 2.0 \times 100\pi = 190\pi \approx \boxed{597\ \mathrm{V}}.
\]

\textbf{(d) Secondary RMS current:}
The induced e.m.f. of part (c) acts as the AC source for the secondary loop (standard GCE simplification: leakage and coil resistance neglected).
\[
V_{s,\text{rms}} = \mathcal{E}_{s,\text{rms}} = \frac{\mathcal{E}_{s0}}{\sqrt{2}} = \frac{597}{\sqrt{2}} \approx 422\ \mathrm{V}.
\]
Load $R_L = 10\ \Omega$.
\[
I_{s,\text{rms}} = \frac{V_{s,\text{rms}}}{R_L} = \frac{422}{10} = \boxed{42.2\ \mathrm{A}}.
\]

\textbf{(e) Primary current increase with load:} When secondary load $R_L$ decreases, $I_s$ increases. The secondary current creates a magnetic flux opposing the primary flux (Lenz's Law). This reduces the net flux in the core, which reduces the back e.m.f. in the primary. With reduced back e.m.f., the primary current $I_p$ increases (drawn from the source) to restore the flux balance. The primary current has two components: a magnetizing component (lagging $90^\circ$) and a load component (in phase with $V_p$, reflected from the secondary).
\end{exoresolu}

\begin{methode}[Energy Stored in an Inductor and Magnetic Energy Density]
\begin{enumerate}
\item \textbf{Work done against back e.m.f.:} To establish current $I$ in an inductor $L$, the source does work against $\mathcal{E}_L = -L dI/dt$. Power $P = \mathcal{E}_L I = L I \frac{dI}{dt}$.
\item \textbf{Total Energy:} $U = \int_0^I P dt = \int_0^I L I dI = \frac{1}{2} L I^2$.
\item \textbf{Magnetic Energy Density:} For a solenoid, $U = \frac{1}{2} (\mu_0 n^2 A l) I^2$. Volume $V = Al$. Magnetic field $B = \mu_0 n I$.
\[
u = \frac{U}{V} = \frac{\frac{1}{2} \mu_0 n^2 A l I^2}{Al} = \frac{1}{2} \mu_0 n^2 I^2 = \frac{1}{2} \frac{B^2}{\mu_0} \quad (\text{in vacuum/air}).
\]
In a material with permeability $\mu = \mu_0 \mu_r$: $u = \frac{1}{2} \frac{B^2}{\mu} = \frac{1}{2} B H$.
\item \textbf{Application:} Energy stored in the magnetic field of a transformer core, or in the stray inductance of cables (surge protection).
\end{enumerate}
\end{methode}

\begin{exoresolu}[Energy in a Solenoid and Magnetic Energy Density]
A long air-cored solenoid has length $l = 0.50\ \mathrm{m}$, cross-sectional area $A = 4.0 \times 10^{-4}\ \mathrm{m^2}$, and $N = 1000$ turns. It carries a current $I = 2.0\ \mathrm{A}$. (a) Calculate the self-inductance $L$ of the solenoid. (b) Calculate the magnetic flux density $B$ inside the solenoid. (c) Determine the total energy stored in the magnetic field using both $U = \frac{1}{2}LI^2$ and the energy density formula $u = \frac{B^2}{2\mu_0}$. (d) If the current is reduced to zero uniformly in $1.0\ \mathrm{ms}$, calculate the average magnitude of the induced e.m.f. across the solenoid.
\tcblower
\textbf{Solution.}
\textbf{Constants:} $\mu_0 = 4\pi \times 10^{-7}\ \mathrm{H\,m^{-1}}$.
\textbf{Data:} $l=0.50\ \mathrm{m}$, $A=4.0\times 10^{-4}\ \mathrm{m^2}$, $N=1000$, $I=2.0\ \mathrm{A}$.

\textbf{(a) Self-inductance $L$:}
\[
L = \frac{\mu_0 N^2 A}{l} = \frac{(4\pi \times 10^{-7}) \times (1000)^2 \times (4.0 \times 10^{-4})}{0.50}
\]
\[
L = \frac{4\pi \times 10^{-7} \times 10^6 \times 4.0 \times 10^{-4}}{0.5} = 32\pi \times 10^{-5} = \boxed{1.01 \times 10^{-3}\ \mathrm{H}} \ (1.01\ \mathrm{mH}).
\]

\textbf{(b) Magnetic Flux Density $B$:}
$n = N/l = 1000 / 0.50 = 2000\ \mathrm{turns/m}$.
\[
B = \mu_0 n I = (4\pi \times 10^{-7}) \times 2000 \times 2.0 = 16\pi \times 10^{-4} = \boxed{5.03 \times 10^{-3}\ \mathrm{T}} \ (5.03\ \mathrm{mT}).
\]

\textbf{(c) Energy Stored:}
\textit{Method 1:} $U = \frac{1}{2} L I^2 = 0.5 \times (1.005 \times 10^{-3}) \times (2.0)^2 = \boxed{2.01 \times 10^{-3}\ \mathrm{J}} \ (2.01\ \mathrm{mJ})$.
\textit{Method 2:} Energy density $u = \frac{B^2}{2\mu_0} = \frac{(5.027 \times 10^{-3})^2}{2 \times 4\pi \times 10^{-7}} = \frac{25.27 \times 10^{-6}}{8\pi \times 10^{-7}} \approx 10.05\ \mathrm{J\,m^{-3}}$.
Volume $V = A l = (4.0 \times 10^{-4}) \times 0.50 = 2.0 \times 10^{-4}\ \mathrm{m^3}$.
$U = u V = 10.05 \times 2.0 \times 10^{-4} = \boxed{2.01 \times 10^{-3}\ \mathrm{J}}$. \textit{Results agree.}

\textbf{(d) Average induced e.m.f. during switch-off:}
$\Delta I = 0 - 2.0 = -2.0\ \mathrm{A}$, $\Delta t = 1.0\ \mathrm{ms} = 1.0 \times 10^{-3}\ \mathrm{s}$.
\[
|\mathcal{E}_{\text{avg}}| = L \left|\frac{\Delta I}{\Delta t}\right| = (1.005 \times 10^{-3}) \times \frac{2.0}{1.0 \times 10^{-3}} = \boxed{2.01\ \mathrm{V}}.
\]
\textit{Note:} This voltage spike is why inductive loads need flyback diodes/snubbers.
\end{exoresolu}

\begin{methode}[Solving Complex Induction Problems: Combined Motion and Field Change]
\begin{enumerate}
\item \textbf{General Faraday's Law:} $\mathcal{E} = -\dfrac{d\Phi}{dt} = -\dfrac{d}{dt}\int \vec{B} \cdot d\vec{A}$.
\item \textbf{Two contributions:} a \emph{transformer e.m.f.} from the time-varying $\vec{B}$ through a fixed loop, and a \emph{motional e.m.f.} from conductors moving in $\vec{B}$: $\mathcal{E}_{\text{mot}} = \oint (\vec{v} \times \vec{B}) \cdot d\vec{l}$.
\item \textbf{Strategy:}
    \begin{enumerate}
    \item Define the loop (closed circuit) clearly at instant $t$.
    \item Compute total flux $\Phi(t)$ through that loop.
    \item Differentiate $\Phi(t)$ w.r.t. time. This automatically accounts for both changing $B$ and changing area/orientation.
    \item Alternatively, compute motional emf on moving parts ($Blv$) and transformer emf on stationary parts separately, then sum them (Kirchhoff's loop rule).
    \end{enumerate}
\item \textbf{Sign convention:} Use Lenz's law / Right-hand rule to assign polarity to the calculated magnitude.
\end{enumerate}
\end{methode}

\begin{exoresolu}[Rod Moving in a Time-Varying Magnetic Field]
A conducting rod of length $l = 0.20\ \mathrm{m}$ slides on two parallel horizontal rails separated by $l$. The rails are connected at one end by a resistor $R = 5.0\ \Omega$. The rod moves with constant velocity $v = 3.0\ \mathrm{m\,s^{-1}}$ away from the resistor, increasing the loop area. The magnetic field is vertical, directed downwards, and its magnitude varies with time as $B(t) = (0.10 + 0.050 t)\ \mathrm{T}$, where $t$ is in seconds. At $t = 0$, the rod is at a distance $x_0 = 0.10\ \mathrm{m}$ from the resistor. (a) Derive an expression for the magnetic flux through the loop as a function of time. (b) Calculate the induced e.m.f. at $t = 2.0\ \mathrm{s}$. (c) Determine the magnitude and direction of the induced current at $t = 2.0\ \mathrm{s}$. (d) Calculate the magnetic force on the rod at this instant and the mechanical power required to maintain constant velocity.
\tcblower
\textbf{Solution.}
\textbf{Setup:} Loop area $A(t) = l \times x(t)$ with $x(t) = x_0 + vt = 0.10 + 3.0t$.
$\vec{B}(t) = (0.10 + 0.050t)\ \mathrm{T}$ downwards. Define positive flux downwards, so $\Phi(t) = B(t)A(t)$; Lenz's law will give the direction.

\textbf{(a) Flux $\Phi(t)$:}
\[
\Phi(t) = B(t) \cdot A(t) = (0.10 + 0.050t) \times [0.20 \times (0.10 + 3.0t)]
\]
Expand the product:
\[
(0.10 + 0.05t)(0.10 + 3t) = 0.01 + 0.30t + 0.005t + 0.15t^2 = 0.01 + 0.305t + 0.15t^2
\]
\[
\Phi(t) = 0.20(0.01 + 0.305t + 0.15t^2) = \boxed{0.002 + 0.061t + 0.030t^2\ \mathrm{Wb}}.
\]

\textbf{(b) Induced e.m.f. at $t=2.0\ \mathrm{s}$:}
\[
\mathcal{E} = -\frac{d\Phi}{dt} = -(0.061 + 0.060t)
\]
At $t=2.0\ \mathrm{s}$:
\[
\mathcal{E} = -(0.061 + 0.120) = -0.181\ \mathrm{V}.
\]
Magnitude $\boxed{0.181\ \mathrm{V}}$. The negative sign means the induced current produces flux opposing the increase. Since $\Phi$ (downwards) is increasing, the induced flux is upwards, so the current is \emph{anticlockwise} (viewed from above).

\textbf{(c) Induced Current:}
\[
I = \frac{|\mathcal{E}|}{R} = \frac{0.181}{5.0} = \boxed{0.0362\ \mathrm{A}} \ (36.2\ \mathrm{mA}).
\]
Direction: anticlockwise viewed from above. With the resistor at the near end and the rod moving away, the current in the rod flows \emph{towards the resistor side of the loop} (from the far rail to the near rail along the rod).

\textbf{(d) Magnetic Force and Power:}
Force on rod: $\vec{F}_B = I \vec{l} \times \vec{B}$. With $\vec{l}$ along the current in the rod and $\vec{B}$ downwards, the cross product gives a horizontal force \textbf{opposing the motion} (as Lenz's law requires).
Magnitude: $F_B = B I l$.
At $t=2.0\ \mathrm{s}$: $B = 0.10 + 0.050(2) = 0.20\ \mathrm{T}$, $I = 0.0362\ \mathrm{A}$, $l = 0.20\ \mathrm{m}$.
\[
F_B = 0.20 \times 0.0362 \times 0.20 = \boxed{1.45 \times 10^{-3}\ \mathrm{N}} \text{ (opposing motion)}.
\]
External force: $F_{\text{ext}} = \boxed{1.45 \times 10^{-3}\ \mathrm{N}}$ in the direction of motion.
Mechanical Power: $P_{\text{mech}} = F_{\text{ext}} v = (1.448 \times 10^{-3}) \times 3.0 = \boxed{4.34 \times 10^{-3}\ \mathrm{W}} \ (4.34\ \mathrm{mW})$.

\textit{Energy check.} Electrical power dissipated: $P_{\text{elec}} = I^2 R = (0.0362)^2 \times 5.0 = 6.55 \times 10^{-3}\ \mathrm{W}$, which is \emph{larger} than $P_{\text{mech}}$. There is no contradiction: the magnetic field is also changing with time, so the external agent driving $B(t)$ supplies the difference. Splitting the e.m.f.:
\[
\frac{d\Phi}{dt} = B\frac{dA}{dt} + A\frac{dB}{dt}, \qquad \frac{dA}{dt} = lv = 0.60\ \mathrm{m^2\,s^{-1}}.
\]
At $t=2.0\ \mathrm{s}$: $A = lx = 0.20 \times (0.10 + 6.0) = 1.22\ \mathrm{m^2}$, $dB/dt = 0.050\ \mathrm{T\,s^{-1}}$.
Motional part: $\mathcal{E}_{\text{mot}} = B\,\dfrac{dA}{dt} = 0.20 \times 0.60 = 0.12\ \mathrm{V}$.
Transformer part: $\mathcal{E}_{\text{trans}} = A\,\dfrac{dB}{dt} = 1.22 \times 0.050 = 0.061\ \mathrm{V}$.
Total: $0.12 + 0.061 = 0.181\ \mathrm{V}$, matching (b).
Power balance: $P_{\text{mech}} = \mathcal{E}_{\text{mot}} I = 0.12 \times 0.0362 = 4.34\ \mathrm{mW}$; field source supplies $\mathcal{E}_{\text{trans}} I = 0.061 \times 0.0362 = 2.21\ \mathrm{mW}$; sum $= 6.55\ \mathrm{mW} = P_{\text{elec}}$. Energy is conserved.
\end{exoresolu}

\begin{methode}[Exam Technique: Graphs and Quantitative Analysis]
\begin{enumerate}
\item \textbf{Flux vs Time Graph:} Slope $= d\Phi/dt$. Induced e.m.f. $\mathcal{E} = - \text{slope}$. Steepest slope $\Rightarrow$ max $\mathcal{E}$. Zero slope $\Rightarrow \mathcal{E}=0$.
\item \textbf{Current vs Time (L-R Growth):} Exponential rise $I = I_0(1-e^{-t/\tau})$. Tangent at $t=0$ has slope $I_0/\tau = E/L$. Intercept of tangent at $I=I_0$ gives $\tau$.
\item \textbf{Current vs Time (L-R Decay):} Exponential fall $I = I_0 e^{-t/\tau}$. Tangent at $t=0$ slope $-I_0/\tau$. Time to halve $t_{1/2} = \tau \ln 2 \approx 0.693\tau$.
\item \textbf{AC Generator Output:} Sinusoidal $\mathcal{E} = \mathcal{E}_0 \sin(\omega t)$. Peak $\mathcal{E}_0 = NBA\omega$. Period $T = 2\pi/\omega$. RMS $\mathcal{E}_{\text{rms}} = \mathcal{E}_0/\sqrt{2}$.
\item \textbf{DC Generator Output (Commutator):} Rectified sine $|\mathcal{E}_0 \sin(\omega t)|$. Average (DC) value $\mathcal{E}_{\text{avg}} = \frac{2}{\pi}\mathcal{E}_0 \approx 0.637\mathcal{E}_0$.
\item \textbf{Units Check:} Always write units in final answer. $\mathrm{Wb}, \mathrm{T}, \mathrm{H}, \mathrm{V}, \mathrm{A}, \mathrm{s}, \mathrm{J}$.
\item \textbf{Direction Language:} ``Clockwise viewed from...'', ``From X to Y'', ``End A is positive relative to End B''.
\end{enumerate}
\end{methode}

\begin{exemplebox}[Reading an L-R Growth Graph]
A graph of current $I$ against time $t$ for an L-R circuit rises from $0$ and levels off at $I_0 = 1.2\ \mathrm{A}$. The tangent at the origin reaches the line $I = I_0$ at $t = 0.20\ \mathrm{s}$. We read directly: time constant $\tau = 0.20\ \mathrm{s}$, initial slope $= I_0/\tau = 6.0\ \mathrm{A\,s^{-1}}$, and since $E = 12\ \mathrm{V}$, $R = E/I_0 = 10\ \Omega$ and $L = E/(\text{initial slope}) = 2.0\ \mathrm{H}$. A single graph therefore yields all the circuit parameters.
\end{exemplebox}

\newpage

\coursec{Exercises}

\courssub{I check my knowledge}

\begin{exercice}[Multiple choice questions]
For each question, choose the single correct answer and justify your choice briefly.
\begin{enumerate}
\item The magnetic flux through a plane coil of $N$ turns, area $A$, whose normal makes an angle $\theta$ with a uniform field $B$ is:
\begin{enumerate}
\item $\Phi = NBA\sin\theta$
\item $\Phi = NBA\cos\theta$
\item $\Phi = BA\cos\theta$
\item $\Phi = NBA\tan\theta$
\end{enumerate}
\item The magnitude of the e.m.f. induced in a circuit equals:
\begin{enumerate}
\item the magnetic flux through the circuit;
\item the rate of change of flux linkage with time;
\item the flux linkage divided by the resistance;
\item the product of the flux and the time.
\end{enumerate}
\item The current in an $L$--$R$ circuit connected to a DC supply grows according to:
\begin{enumerate}
\item $I = I_0\,e^{-t/\tau}$
\item $I = I_0\left(1 - e^{-t/\tau}\right)$
\item $I = I_0\,t/\tau$
\item $I = I_0\left(1 + e^{-t/\tau}\right)$
\end{enumerate}
\item The energy stored in an inductor carrying a current $I$ is:
\begin{enumerate}
\item $W = \tfrac{1}{2}LI^2$
\item $W = LI^2$
\item $W = \tfrac{1}{2}L^2I$
\item $W = \tfrac{1}{2}LI$
\end{enumerate}
\end{enumerate}
\end{exercice}

\begin{exercice}[True or false?]
State whether each statement is true or false, and justify your answer in one or two sentences.
\begin{enumerate}
\item An e.m.f. is induced in a coil whenever a magnetic field exists near it, even if the field is constant.
\item According to Lenz's law, the induced current always flows so as to oppose the change of flux that produces it.
\item A split-ring commutator is fitted to a simple AC generator.
\item The self-inductance $L$ of a coil is defined by the flux-linkage equation $N\Phi = LI$.
\item In the steady state, an ideal inductor connected to a DC supply behaves like a simple connecting wire.
\end{enumerate}
\end{exercice}

\begin{exercice}[Definitions and statements]
\begin{enumerate}
\item Define the \cle{weber} in terms of the tesla and state the SI unit of flux linkage.
\item State Faraday's law of electromagnetic induction and write its mathematical form, explaining every symbol.
\item State Lenz's law and explain how it follows from the principle of conservation of energy.
\item Define the \cle{mutual inductance} $M$ of two coils from the equation $N_2\Phi_2 = MI_1$, and state its SI unit.
\end{enumerate}
\end{exercice}

\begin{exercice}[Quick calculations]
\begin{enumerate}
\item A flat coil of area $25\ \mathrm{cm^2}$ lies perpendicular to a uniform field of $0.30\ \mathrm{T}$. Calculate the flux through the coil.
\item A conductor of length $0.50\ \mathrm{m}$ moves at $8.0\ \mathrm{m\,s^{-1}}$ perpendicular to a field of $0.12\ \mathrm{T}$. Calculate the e.m.f. induced between its ends.
\item An inductor of $L = 0.50\ \mathrm{H}$ carries a steady current of $2.0\ \mathrm{A}$. Calculate the energy stored in its magnetic field.
\item The current in a coil of self-inductance $80\ \mathrm{mH}$ changes at the rate of $150\ \mathrm{A\,s^{-1}}$. Calculate the magnitude of the self-induced e.m.f.
\end{enumerate}
\end{exercice}

\courssub{I practise}

\begin{exercice}[E.M.F. and induced charge in a coil]
A coil of $200$ turns and area $15\ \mathrm{cm^2}$ is placed with its plane perpendicular to a uniform magnetic field which falls uniformly from $0.40\ \mathrm{T}$ to zero in $0.050\ \mathrm{s}$.
\begin{enumerate}
\item Calculate the change of flux linkage of the coil.
\item Calculate the magnitude of the e.m.f. induced in the coil.
\item The coil forms part of a closed circuit of total resistance $5.0\ \Omega$. Calculate the induced current and the total charge that flows round the circuit while the field is falling.
\end{enumerate}
\end{exercice}

\begin{exercice}[Motional e.m.f. in the Earth's field]
A metal rod of length $0.80\ \mathrm{m}$ is moved horizontally at $12\ \mathrm{m\,s^{-1}}$ in a direction perpendicular both to its length and to the horizontal component of the Earth's magnetic field, of magnitude $2.0\times10^{-5}\ \mathrm{T}$.
\begin{enumerate}
\item Calculate the e.m.f. induced between the ends of the rod.
\item Using Fleming's right-hand rule (or the force on positive charges, $\vec{F}=q\vec{v}\times\vec{B}$), explain how you decide which end of the rod becomes positive.
\item Explain why no current flows unless the rod forms part of a closed circuit, although the e.m.f. still exists.
\end{enumerate}
\end{exercice}

\begin{exercice}[Growth of current in an L--R circuit]
An inductor of inductance $L = 0.20\ \mathrm{H}$ and negligible resistance is connected in series with a resistor $R = 40\ \Omega$ across a $12\ \mathrm{V}$ DC supply.
\begin{enumerate}
\item Write down the equation giving the growth of the current $I(t)$ after the switch is closed, defining all symbols.
\item Calculate the time constant $\tau$ of the circuit.
\item Calculate the steady-state current $I_0$.
\item Calculate the current $2.0\ \mathrm{ms}$ after the switch is closed.
\item Calculate the energy stored in the inductor in the steady state.
\end{enumerate}
\end{exercice}

\begin{exercice}[Mutual induction between two coils]
Two neighbouring coils have mutual inductance $M = 15\ \mathrm{mH}$. The current in the primary coil changes uniformly from $3.0\ \mathrm{A}$ to zero in $10\ \mathrm{ms}$. The secondary coil has $800$ turns.
\begin{enumerate}
\item Define the mutual inductance $M$ from the flux-linkage equation $N_2\Phi_2 = MI_1$.
\item Calculate the average e.m.f. induced in the secondary coil.
\item Calculate the average change of flux through each turn of the secondary coil during the $10\ \mathrm{ms}$ interval.
\item State two practical ways of increasing the mutual inductance of the pair of coils.
\end{enumerate}
\end{exercice}

\courssub{I prepare for the GCE}

\begin{exercice}[Structured question: the rotating-coil generator]
A rectangular coil of $N = 500$ turns and area $A = 100\ \mathrm{cm^2}$ rotates at a steady frequency of $50\ \mathrm{Hz}$ about an axis perpendicular to a uniform magnetic field of flux density $B = 0.15\ \mathrm{T}$.
\begin{enumerate}
\item[(a)] Write down the expression for the flux linkage of the coil at time $t$, taking $t = 0$ when the plane of the coil is perpendicular to the field. \hfill (2 marks)
\item[(b)] Hence show that the induced e.m.f. is $e(t) = E_0\sin(\omega t)$ and calculate the peak e.m.f. $E_0$. \hfill (4 marks)
\item[(c)] Calculate the r.m.s. value of the e.m.f. \hfill (2 marks)
\item[(d)] Sketch, on the same time axis, two complete cycles of the output e.m.f. when the coil is connected to an external circuit through:
\begin{enumerate}
\item[(i)] two slip rings;
\item[(ii)] a split-ring commutator.
\end{enumerate}
Label the axes with numerical values. \hfill (4 marks)
\item[(e)] State and explain the effect on the peak e.m.f. and on the frequency of doubling the speed of rotation. \hfill (3 marks)
\end{enumerate}
\hfill (Total: 15 marks)
\end{exercice}

\begin{exercice}[Essay question: the magnet falling through a copper tube]
A small but strong bar magnet is dropped through a long vertical copper tube. It is observed to fall much more slowly than an identical unmagnetised object falling through the same tube.
\begin{enumerate}
\item[(a)] Explain, with the help of Lenz's law and a clear description of the induced currents in the tube, why the magnet experiences an upward force opposing its fall. \hfill (5 marks)
\item[(b)] Show, using the principle of conservation of energy, that the gravitational potential energy lost by the magnet reappears as thermal energy in the tube. \hfill (3 marks)
\item[(c)] Explain why the magnet eventually reaches a constant terminal velocity. \hfill (3 marks)
\item[(d)] State and explain what changes, and why, if the tube is replaced by an identical tube having a narrow vertical slit along its whole length. \hfill (3 marks)
\item[(e)] Suggest one practical application of this phenomenon in engineering. \hfill (1 mark)
\end{enumerate}
\hfill (Total: 15 marks)
\end{exercice}

\begin{exercice}[Data analysis: decay of current in an L--R circuit]
The steady current through a coil of inductance $L$ in series with a resistor of resistance $R = 20\ \Omega$ is suddenly allowed to decay by short-circuiting the supply. The current $I$ is measured at successive times $t$:

\begin{center}
\begin{tabular}{|l|c|c|c|c|c|c|}
\hline
\rowcolor{popblueL} $t$ / ms & 0 & 2.0 & 4.0 & 6.0 & 8.0 & 10.0 \\
\hline
$I$ / A & 0.60 & 0.40 & 0.27 & 0.18 & 0.12 & 0.080 \\
\hline
\end{tabular}
\end{center}

\begin{enumerate}
\item[(a)] Starting from Kirchhoff's voltage law applied to the $L$--$R$ loop, show that the decaying current obeys $I = I_0\,e^{-Rt/L}$. \hfill (4 marks)
\item[(b)] Copy and complete the table with values of $\ln(I/\mathrm{A})$. \hfill (2 marks)
\item[(c)] Plot a graph of $\ln(I/\mathrm{A})$ against $t$, choosing suitable scales. \hfill (3 marks)
\item[(d)] Explain the shape of the graph and determine its gradient. \hfill (3 marks)
\item[(e)] Hence determine the time constant $\tau$ of the circuit and deduce the inductance $L$ of the coil. \hfill (3 marks)
\end{enumerate}
\hfill (Total: 15 marks)
\end{exercice}



\newpage

\coursec{Solutions to the exercises}

\courssub{I check my knowledge}

\begin{corrige}[Exercise 1 --- Multiple choice questions]
\begin{enumerate}
\item \textbf{(b)} $\Phi = NBA\cos\theta$. Only the component of $\vec{B}$ along the normal to the coil contributes to the flux; multiplying by the $N$ turns gives the flux linkage $N\Phi = NBA\cos\theta$. Option (c) forgets the number of turns, and (a) would be correct only if $\theta$ were measured from the plane of the coil.
\item \textbf{(b)}. Faraday's law: $e = -\dfrac{d(N\Phi)}{dt}$; the magnitude of the induced e.m.f. equals the rate of change of flux linkage with time. Flux itself (a) induces nothing if it is constant.
\item \textbf{(b)} $I = I_0\left(1 - e^{-t/\tau}\right)$ with $\tau = L/R$: the current starts at zero (an inductor opposes any sudden change of current) and rises asymptotically towards $I_0 = E/R$. Option (a) describes the \emph{decay} of current.
\item \textbf{(a)} $W = \tfrac{1}{2}LI^2$, obtained by integrating the power $p = LI\,\dfrac{dI}{dt}$ supplied to the inductor; it is the magnetic analogue of $\tfrac{1}{2}CV^2$ for a capacitor.
\end{enumerate}
\end{corrige}

\begin{corrige}[Exercise 2 --- True or false?]
\begin{enumerate}
\item \textbf{False.} An e.m.f. is induced only when the flux \emph{linkage changes}; a constant field produces a constant flux and hence no induced e.m.f., however strong the field may be.
\item \textbf{True.} This is exactly the statement of Lenz's law; it is a direct consequence of the conservation of energy (the induced current can never aid the change that creates it).
\item \textbf{False.} A simple AC generator uses two \cle{slip rings}, one connected to each end of the coil; the \cle{split-ring commutator} is fitted to a simple DC generator to reverse the connections every half-turn and rectify the output.
\item \textbf{True.} $L$ is the constant of proportionality in $N\Phi = LI$, i.e. the flux linkage per unit current; it is measured in henry (H), with $1\ \mathrm{H} = 1\ \mathrm{Wb\,A^{-1}}$.
\item \textbf{True.} In the steady state $\dfrac{dI}{dt} = 0$, so the back e.m.f. $L\dfrac{dI}{dt} = 0$; an ideal (resistance-free) inductor then offers no opposition at all and behaves like a plain connecting wire.
\end{enumerate}
\end{corrige}

\begin{corrige}[Exercise 3 --- Definitions and statements]
\begin{enumerate}
\item One weber is the magnetic flux through an area of $1\ \mathrm{m^2}$ placed perpendicular to a uniform magnetic field of flux density $1\ \mathrm{T}$:
\[ 1\ \mathrm{Wb} = 1\ \mathrm{T\,m^2}. \]
The flux linkage $N\Phi$ is measured in weber-turns (Wb-turns), usually written simply Wb.
\item \textbf{Faraday's law:} whenever the magnetic flux linkage of a circuit changes, an e.m.f. is induced in the circuit whose magnitude is equal to the rate of change of flux linkage:
\[ e = -\frac{d(N\Phi)}{dt}, \]
where $e$ is the induced e.m.f. (V), $N$ the number of turns of the coil, $\Phi$ the magnetic flux through one turn (Wb) and $t$ the time (s). The minus sign expresses Lenz's law: the induced e.m.f. opposes the change producing it.
\item \textbf{Lenz's law:} the direction of the induced current is always such that its magnetic effect opposes the change of flux producing it. Suppose the opposite were true: the induced current would reinforce the change of flux, which would induce a still larger current, and so on --- energy would be created from nothing, contradicting the conservation of energy. In reality, external work must be done against the opposing (Lenz) force, and this work is exactly the electrical energy that appears in the circuit.
\item $M$ is the constant of proportionality between the flux linkage $N_2\Phi_2$ of the secondary coil and the current $I_1$ in the primary coil that produces it:
\[ M = \frac{N_2\Phi_2}{I_1} \qquad \text{(equivalently } N_1\Phi_1 = MI_2\text{)}. \]
Its SI unit is the henry (H).
\end{enumerate}
\end{corrige}

\begin{corrige}[Exercise 4 --- Quick calculations]
\begin{enumerate}
\item Convert the area: $A = 25\ \mathrm{cm^2} = 2.5\times10^{-3}\ \mathrm{m^2}$. The field is normal to the plane of the coil, so
\[ \Phi = BA = 0.30 \times 2.5\times10^{-3} = 7.5\times10^{-4}\ \mathrm{Wb}. \]
\item Motional e.m.f.:
\[ e = Blv = 0.12 \times 0.50 \times 8.0 = 0.48\ \mathrm{V}. \]
\item Energy stored:
\[ W = \tfrac{1}{2}LI^2 = \tfrac{1}{2}\times 0.50 \times (2.0)^2 = 1.0\ \mathrm{J}. \]
\item Self-induced e.m.f.:
\[ |e| = L\frac{dI}{dt} = 0.080 \times 150 = 12\ \mathrm{V}. \]
\end{enumerate}
\end{corrige}

\begin{attention}[Units first!]
Always convert $\mathrm{cm^2}$ to $\mathrm{m^2}$ ($1\ \mathrm{cm^2} = 10^{-4}\ \mathrm{m^2}$) and $\mathrm{mH}$ to $\mathrm{H}$ before substituting; most lost marks in these questions come from unit errors, not from physics.
\end{attention}

\courssub{I practise}

\begin{corrige}[Exercise 5 --- E.M.F. and induced charge in a coil]
\begin{enumerate}
\item The field is normal to the plane of the coil, so the flux through one turn is $\Phi = BA$. The change of flux linkage is
\[ \Delta(N\Phi) = NA\,\Delta B = 200 \times 15\times10^{-4} \times (0.40 - 0) = 0.12\ \mathrm{Wb\text{-}turns}. \]
\item Faraday's law, using the uniform (average) rate of change:
\[ |e| = \frac{\Delta(N\Phi)}{\Delta t} = \frac{0.12}{0.050} = 2.4\ \mathrm{V}. \]
\item Induced current:
\[ I = \frac{|e|}{R} = \frac{2.4}{5.0} = 0.48\ \mathrm{A}. \]
Total charge circulated:
\[ Q = I\,\Delta t = 0.48 \times 0.050 = 2.4\times10^{-2}\ \mathrm{C}. \]
Check with the general result $Q = \dfrac{\Delta(N\Phi)}{R} = \dfrac{0.12}{5.0} = 0.024\ \mathrm{C}$: the charge depends only on the total change of flux linkage and the resistance, not on how quickly the change occurs.
\end{enumerate}
\end{corrige}

\begin{corrige}[Exercise 6 --- Motional e.m.f. in the Earth's field]
\begin{enumerate}
\item The rod, its velocity and the field are mutually perpendicular, so
\[ e = Blv = 2.0\times10^{-5} \times 0.80 \times 12 = 1.9\times10^{-4}\ \mathrm{V} \approx 0.19\ \mathrm{mV}. \]
\item The free charge carriers in the rod are carried along with velocity $\vec{v}$ and each experiences the magnetic force $\vec{F} = q\vec{v}\times\vec{B}$, directed along the rod. Fleming's right-hand rule (First finger = Field, thuMb = Motion, seCond finger = induced Current) gives the same direction. Positive charges are pushed towards one end of the rod, which becomes positive, while the other end, depleted of positive charge, becomes negative. Charge accumulates until the electric field it creates inside the rod exactly balances the magnetic force ($qE = qvB$, giving $e = El = Blv$).
\item The separation of charge \emph{is} the e.m.f.: it exists whenever the rod moves through the field, whether or not the circuit is closed. But a current requires a complete conducting path; with the rod open-ended, charges merely pile up at the ends and no continuous flow is possible. Closing the circuit allows the e.m.f. to drive a current $I = e/R$.
\end{enumerate}
\end{corrige}

\begin{corrige}[Exercise 7 --- Growth of current in an L--R circuit]
\begin{enumerate}
\item Kirchhoff's voltage law around the loop gives $E = RI + L\dfrac{dI}{dt}$. Solving this differential equation with the initial condition $I(0) = 0$:
\[ I(t) = I_0\left(1 - e^{-t/\tau}\right), \qquad I_0 = \frac{E}{R},\quad \tau = \frac{L}{R}, \]
where $E$ is the supply e.m.f., $I_0$ the steady-state (final) current and $\tau$ the time constant of the circuit.
\item Time constant:
\[ \tau = \frac{L}{R} = \frac{0.20}{40} = 5.0\times10^{-3}\ \mathrm{s} = 5.0\ \mathrm{ms}. \]
\item Steady-state current (the inductor then behaves like a wire):
\[ I_0 = \frac{E}{R} = \frac{12}{40} = 0.30\ \mathrm{A}. \]
\item At $t = 2.0\ \mathrm{ms}$:
\[ I = 0.30\left(1 - e^{-2.0/5.0}\right) = 0.30\left(1 - e^{-0.40}\right) = 0.30 \times 0.330 = 9.9\times10^{-2}\ \mathrm{A}. \]
\item Energy stored in the steady state:
\[ W = \tfrac{1}{2}LI_0^2 = \tfrac{1}{2}\times 0.20 \times (0.30)^2 = 9.0\times10^{-3}\ \mathrm{J}. \]
\end{enumerate}
\end{corrige}

\begin{corrige}[Exercise 8 --- Mutual induction between two coils]
\begin{enumerate}
\item $M$ is the flux linkage of the secondary coil per unit current in the primary coil:
\[ M = \frac{N_2\Phi_2}{I_1}, \]
where $\Phi_2$ is the flux through one turn of the secondary produced by the primary current $I_1$; $M$ is measured in henry.
\item The current changes uniformly, so $\left|\dfrac{dI_1}{dt}\right| = \dfrac{3.0}{10\times10^{-3}} = 300\ \mathrm{A\,s^{-1}}$, and
\[ |e_2| = M\left|\frac{dI_1}{dt}\right| = 15\times10^{-3} \times 300 = 4.5\ \mathrm{V}. \]
\item From $N_2\,\Delta\Phi_2 = M\,\Delta I_1$:
\[ \Delta\Phi_2 = \frac{M\,\Delta I_1}{N_2} = \frac{15\times10^{-3}\times 3.0}{800} = 5.6\times10^{-5}\ \mathrm{Wb}. \]
\item Any two of: wind both coils on a common soft-iron core, which concentrates and guides the flux from one coil to the other; bring the coils closer together or wind one directly on top of the other; increase the number of turns of either coil; orient the coils coaxially so that maximum flux links them.
\end{enumerate}
\end{corrige}

\courssub{I prepare for the GCE}

\begin{corrige}[Exercise 9 --- Structured question: the rotating-coil generator]
\begin{enumerate}
\item[(a)] At $t = 0$ the plane of the coil is perpendicular to $\vec{B}$, so the normal to the coil is parallel to $\vec{B}$; after time $t$ the normal makes the angle $\theta = \omega t$ with $\vec{B}$. Hence
\[ N\Phi = NBA\cos(\omega t), \qquad \omega = 2\pi f = 2\pi \times 50 = 314\ \mathrm{rad\,s^{-1}}. \]
\textbf{Mark scheme (2):} 1 mark for $\theta = \omega t$ (or $\cos$ dependence consistent with the stated initial condition); 1 mark for the complete expression including $N$.
\item[(b)] Faraday's law:
\[ e = -\frac{d(N\Phi)}{dt} = -\frac{d}{dt}\left[NBA\cos(\omega t)\right] = NBA\omega\sin(\omega t) = E_0\sin(\omega t), \]
with peak value
\[ E_0 = NBA\omega = 500 \times 0.15 \times 100\times10^{-4} \times 314 = 2.4\times10^{2}\ \mathrm{V}\ (\approx 236\ \mathrm{V}). \]
\textbf{Mark scheme (4):} 1 mark for quoting Faraday's law; 1 mark for correct differentiation ($\frac{d}{dt}\cos\omega t = -\omega\sin\omega t$); 1 mark for identifying $E_0 = NBA\omega$; 1 mark for the numerical value with unit (accept $235$--$236\ \mathrm{V}$). \emph{Examiner's expectation:} the area must be converted to $\mathrm{m^2}$ ($100\ \mathrm{cm^2} = 1.0\times10^{-2}\ \mathrm{m^2}$).
\item[(c)] For a sinusoidal e.m.f.:
\[ E_{\mathrm{rms}} = \frac{E_0}{\sqrt{2}} = \frac{236}{\sqrt{2}} = 1.7\times10^{2}\ \mathrm{V}\ (\approx 167\ \mathrm{V}). \]
\textbf{Mark scheme (2):} 1 mark for the relation $E_{\mathrm{rms}} = E_0/\sqrt{2}$; 1 mark for the value with unit.
\item[(d)] Period $T = \dfrac{1}{f} = 20\ \mathrm{ms}$.
\begin{popfigure}
\begin{center}
\begin{tikzpicture}[scale=1.0]
\draw[->,thick] (0,0) -- (9.2,0) node[below left,align=center,text width=1.6cm]{$t$ / ms};
\draw[->,thick] (0,-2.4) -- (0,2.6) node[above]{$e$ / V};
\foreach \x/\lab in {2/10, 4/20, 6/30, 8/40}{\draw (\x,0.06) -- (\x,-0.06) node[below=2pt]{\lab};}
\draw (0.06,2) -- (-0.06,2) node[left=2pt]{$236$};
\draw (0.06,-2) -- (-0.06,-2) node[left=2pt]{$-236$};
\draw[very thick,popblue,domain=0:8,samples=120] plot (\x,{2*sin(\x*90)});
\node[popblue] at (7.3,2.35) {slip rings (AC)};
\draw[very thick,poporange,domain=0:8,samples=120] plot (\x,{abs(2*sin(\x*90))});
\node[poporange] at (2.6,-2.3) {split ring (rectified)};
\end{tikzpicture}
\end{center}
\legende{(i) Slip rings: sinusoidal AC output, $e = 236\sin(314\,t)$. (ii) Split-ring commutator: the negative half-cycles are reversed, giving a pulsating unidirectional output of the same peak value and the same period.}
\end{popfigure}
\textbf{Mark scheme (4):} 1 mark for a correct sine wave through the origin (slip rings); 1 mark for the rectified shape with all half-cycles on the same side and the same peak (split ring); 1 mark for labelled axes with units; 1 mark for correct numerical scales ($T = 20\ \mathrm{ms}$, $E_0 = 236\ \mathrm{V}$).
\item[(e)] Since $E_0 = NBA\omega = 2\pi f NBA$, both the frequency and the peak e.m.f. are directly proportional to the speed of rotation. Doubling the speed doubles the frequency ($50\ \mathrm{Hz} \to 100\ \mathrm{Hz}$) and doubles the peak e.m.f. ($236\ \mathrm{V} \to 471\ \mathrm{V}$).
\textbf{Mark scheme (3):} 1 mark for $f$ doubles; 1 mark for $E_0$ doubles; 1 mark for justification via $E_0 = 2\pi f NBA$.
\end{enumerate}
\end{corrige}

\begin{corrige}[Exercise 10 --- Essay question: the magnet falling through a copper tube]
\begin{enumerate}
\item[(a)] As the magnet falls, the magnetic flux through each section of the tube changes: it \emph{increases} in the region below the magnet (which the magnet approaches) and \emph{decreases} in the region above it (which the magnet leaves). Copper is a good conductor, so these flux changes drive induced (\cle{eddy}) currents in closed loops around the tube wall. By Lenz's law, the currents below the magnet flow in such a direction as to \emph{repel} the approaching pole, while the currents above flow so as to \emph{attract} the receding pole. Both effects produce an upward force on the magnet opposing its fall, so it descends far more slowly than a non-magnetic object.\\
\textbf{Mark scheme (5):} 1 mark for identifying the changing flux; 1 mark for induced currents in the tube wall (copper a conductor); 1 mark for correct use of Lenz's law below the magnet (repulsion); 1 mark for attraction above; 1 mark for the conclusion (net upward braking force).
\item[(b)] The induced currents flow through the resistance of the copper and dissipate electrical energy as heat (Joule heating, $P = I^2R$). The only source of this energy is the work done by gravity. As the magnet descends through a height $h$, it loses gravitational potential energy $mgh$. Once the magnet moves at constant (terminal) velocity its kinetic energy no longer changes, so the \emph{whole} of $mgh$ is converted into thermal energy in the tube: energy is conserved, merely transferred from gravitational potential energy to internal energy of the copper.\\
\textbf{Mark scheme (3):} 1 mark for Joule heating of the tube; 1 mark for identifying gravity as the energy source; 1 mark for the balance $mgh \to$ heat at constant speed.
\item[(c)] The faster the magnet falls, the greater the rate of change of flux $\dfrac{d\Phi}{dt}$, hence the larger the induced e.m.f.s, the induced currents and therefore the braking force. The upward electromagnetic force thus grows with speed until it exactly balances the weight $mg$ (air resistance being negligible). The resultant force is then zero and, by Newton's first law, the magnet continues at a constant \cle{terminal velocity}.\\
\textbf{Mark scheme (3):} 1 mark for braking force increasing with speed (via rate of change of flux); 1 mark for force balance $F = mg$; 1 mark for zero acceleration $\Rightarrow$ constant velocity.
\item[(d)] The slit cuts the closed conducting loops around the circumference of the tube, so the eddy currents can no longer circulate around it. The induced currents, and therefore the braking force, become very small; the magnet then falls almost freely, nearly like the unmagnetised object.\\
\textbf{Mark scheme (3):} 1 mark for recognising that the current loops are interrupted; 1 mark for the consequent near-absence of eddy currents/braking; 1 mark for the conclusion (near-free fall).
\item[(e)] Electromagnetic (eddy-current) braking, e.g. on trains or roller-coasters; electromagnetic damping of moving-coil meters; induction furnaces; electromagnetic sorting of metals. (Any one, 1 mark.)
\end{enumerate}
\end{corrige}

\begin{corrige}[Exercise 11 --- Data analysis: decay of current in an L--R circuit]
\begin{enumerate}
\item[(a)] With the supply short-circuited, the loop contains only $L$ and $R$. Kirchhoff's voltage law gives
\[ L\frac{dI}{dt} + RI = 0 \quad\Longrightarrow\quad \frac{dI}{I} = -\frac{R}{L}\,dt. \]
Integrating between $(0, I_0)$ and $(t, I)$:
\[ \int_{I_0}^{I}\frac{dI'}{I'} = -\frac{R}{L}\int_0^t dt' \quad\Longrightarrow\quad \ln I - \ln I_0 = -\frac{R}{L}t, \]
so
\[ I = I_0\,e^{-Rt/L} = I_0\,e^{-t/\tau}, \qquad \tau = \frac{L}{R}. \]
\textbf{Mark scheme (4):} 1 mark for the Kirchhoff equation; 1 mark for correct separation of variables; 1 mark for correct integration with limits; 1 mark for the final exponential form.
\item[(b)] Completed table:
\begin{center}
\begin{tabular}{|l|c|c|c|c|c|c|}
\hline
\rowcolor{popblueL} $t$ / ms & 0 & 2.0 & 4.0 & 6.0 & 8.0 & 10.0 \\
\hline
$I$ / A & 0.60 & 0.40 & 0.27 & 0.18 & 0.12 & 0.080 \\
\hline
$\ln(I/\mathrm{A})$ & $-0.51$ & $-0.92$ & $-1.31$ & $-1.71$ & $-2.12$ & $-2.53$ \\
\hline
\end{tabular}
\end{center}
\textbf{Mark scheme (2):} all six values correct to 2 d.p. (2 marks); one error (1 mark).
\item[(c)]
\begin{popfigure}
\begin{center}
\begin{tikzpicture}[scale=1.0]
\draw[->,thick] (0,0) -- (10.8,0) node[below left,align=center,text width=1.6cm]{$t$ / ms};
\draw[->,thick] (0,-5.6) -- (0,0.6) node[above]{$\ln(I/\mathrm{A})$};
\foreach \x in {0,2,4,6,8,10}{\draw (\x,0.05) -- (\x,-0.05) node[below=2pt]{\x};}
\foreach \y/\lab in {-1/-0.5, -2/-1.0, -3/-1.5, -4/-2.0, -5/-2.5}{\draw (0.05,\y) -- (-0.05,\y) node[left=2pt]{\lab};}
\draw[popmuted,dashed] (0,-1) -- (10,-1);
\draw[popmuted,dashed] (0,-2) -- (10,-2);
\draw[popmuted,dashed] (0,-3) -- (10,-3);
\draw[popmuted,dashed] (0,-4) -- (10,-4);
\draw[popmuted,dashed] (0,-5) -- (10,-5);
\draw[very thick,popblue] (0,-1.02) -- (10,-5.05);
\foreach \p in {(0,-1.02),(2,-1.83),(4,-2.62),(6,-3.43),(8,-4.24),(10,-5.05)}{\fill[poporange] \p circle (2.2pt);}
\end{tikzpicture}
\end{center}
\legende{Graph of $\ln(I/\mathrm{A})$ against $t$: a straight line of negative gradient passing through the experimental points.}
\end{popfigure}
\textbf{Mark scheme (3):} 1 mark for sensible scales using more than half of each axis; 1 mark for all points plotted correctly; 1 mark for labelled axes with units and a best-fit straight line.
\item[(d)] Taking logarithms of $I = I_0 e^{-Rt/L}$:
\[ \ln I = \ln I_0 - \frac{R}{L}\,t, \]
which is the equation of a straight line of gradient $-\dfrac{R}{L}$ and intercept $\ln I_0$ --- exactly what the graph shows. Using two well-separated points on the best-fit line:
\[ \text{gradient} = \frac{-2.53 - (-0.51)}{(10.0 - 0)\times10^{-3}} = \frac{-2.02}{1.0\times10^{-2}} = -2.0\times10^{2}\ \mathrm{s^{-1}}. \]
\textbf{Mark scheme (3):} 1 mark for deriving $\ln I = \ln I_0 - Rt/L$; 1 mark for identifying the gradient with $-R/L$; 1 mark for a gradient read from a large triangle, value in the range $(1.9$--$2.1)\times10^{2}\ \mathrm{s^{-1}}$.
\item[(e)] Time constant:
\[ \tau = -\frac{1}{\text{gradient}} = \frac{1}{2.0\times10^{2}} = 5.0\times10^{-3}\ \mathrm{s} = 5.0\ \mathrm{ms}. \]
Inductance, from $\tau = L/R$:
\[ L = R\tau = 20 \times 5.0\times10^{-3} = 0.10\ \mathrm{H}. \]
\textbf{Mark scheme (3):} 1 mark for $\tau = -1/\text{gradient}$; 1 mark for $\tau = 5.0\ \mathrm{ms}$; 1 mark for $L = R\tau = 0.10\ \mathrm{H}$ with unit.
\end{enumerate}
\end{corrige}

\begin{attention}[Examiner's advice]
In graphical questions, always draw the gradient triangle using more than half of the drawn line, quote the gradient with its unit (here $\mathrm{s^{-1}}$), and carry the unit through to the final answer. A value of $L$ without ``H'' loses the last mark.
\end{attention}

\newpage

\coursec{Summary sheet}

\begin{popfigure}
\begin{tikzpicture}[
    box/.style={draw=popline, rounded corners=2pt, fill=popblueL, text width=3.6cm, align=center, font=\small, inner sep=4pt},
    head/.style={draw=popblue, rounded corners=2pt, fill=popblue, text=white, text width=3.6cm, align=center, font=\small\bfseries, inner sep=4pt},
    link/.style={draw=popmuted, thick}
]
\node[head] (root) at (0,0) {Electromagnetic Induction};
\node[box, align=center] (flux) at (-5.2,-2.2){\textbf{Magnetic flux}\\ $\Phi = BA\cos\theta$\\ unit: weber (Wb)\\ flux linkage: $N\Phi$};
\node[box, align=center] (faraday) at (0,-2.2){\textbf{Faraday's law}\\ $\mathcal{E} = -N\dfrac{d\Phi}{dt}$};
\node[box, align=center] (lenz) at (5.2,-2.2){\textbf{Lenz's law}\\ induced current opposes the \emph{change} of flux};
\node[box, align=center] (mot) at (-5.2,-4.8){\textbf{Motional e.m.f.}\\ $\mathcal{E} = Blv$\\ ($\vec{v}\perp\vec{B}\perp\vec{l}$)};
\node[box, align=center] (gen) at (0,-4.8){\textbf{Generators}\\ AC: slip rings, $\mathcal{E}=NBA\omega\sin\omega t$\\ DC: split-ring commutator};
\node[box, align=center] (ind) at (5.2,-4.8){\textbf{Inductance}\\ $N\Phi = LI$,\; $N_2\Phi_2 = MI_1$\\ $U = \tfrac{1}{2}LI^2$,\; $\tau = L/R$};
\draw[link] (root) -- (flux);
\draw[link] (root) -- (faraday);
\draw[link] (root) -- (lenz);
\draw[link] (flux) -- (mot);
\draw[link] (faraday) -- (gen);
\draw[link] (lenz) -- (ind);
\end{tikzpicture}
\legende{Chapter PHY08 concept map: from magnetic flux to inductance and the L--R circuit.}
\end{popfigure}

\courssub{Key Definitions}
\begin{itemize}
\item \cle{Magnetic flux} ($\Phi$): $\Phi = \vec{B}\cdot\vec{A} = BA\cos\theta$, where $\theta$ is the angle between $\vec{B}$ and the \emph{normal} to the surface. Scalar quantity; SI unit the weber ($1\ \mathrm{Wb} = 1\ \mathrm{T\,m^2}$).
\item \cle{Flux linkage}: for a coil of $N$ turns, flux linkage $= N\Phi$ (unit: Wb-turns).
\item \cle{Induced e.m.f.} ($\mathcal{E}$): the e.m.f. produced by a changing flux linkage; work done per unit charge; unit volt (V).
\item \cle{Self-inductance} ($L$): the constant in $N\Phi = LI$, i.e.\ the flux linkage per unit current in the \emph{same} coil: $L = N\Phi/I$; unit the henry (H).
\item \cle{Mutual inductance} ($M$): the constant in $N_2\Phi_2 = MI_1$ (flux linkage of coil 2 per unit current in coil 1); symmetrically $N_1\Phi_1 = MI_2$; unit the henry (H).
\item \cle{Time constant} ($\tau$) of an L--R circuit: $\tau = L/R$; the time for the current to reach about $63\%$ of its final value during growth, or to fall to about $37\%$ of its initial value during decay.
\end{itemize}

\courssub{Fundamental Laws and Formulas}
\begin{itemize}
\item \textbf{Faraday's law}: $\mathcal{E} = -\dfrac{d(N\Phi)}{dt}$. The magnitude of the induced e.m.f.\ equals the rate of change of flux linkage: $|\mathcal{E}| = N\left|\dfrac{d\Phi}{dt}\right|$.
\item \textbf{Lenz's law}: the induced current flows in such a direction that its own magnetic effect opposes the \emph{change} of flux producing it. This is the meaning of the minus sign in Faraday's law and is a direct consequence of the conservation of energy.
\item \textbf{Motional e.m.f.}: for a straight conductor of length $l$ moving with velocity $v$ perpendicular to a uniform field $\vec{B}$ (with $\vec{v}$, $\vec{B}$, $\vec{l}$ mutually perpendicular):
\[ \mathcal{E} = Blv. \]
If $\vec{v}$ makes an angle $\theta$ with $\vec{B}$, only the perpendicular component counts: $\mathcal{E} = Blv\sin\theta$.
\item \textbf{AC generator} (coil of $N$ turns, area $A$, rotating at angular speed $\omega$ in uniform $\vec{B}$):
\[ \mathcal{E} = \mathcal{E}_0\sin(\omega t), \qquad \mathcal{E}_0 = NBA\omega. \]
\item \textbf{Inductor (back) e.m.f.}: $\mathcal{E}_L = -L\dfrac{dI}{dt}$.
\item \textbf{L--R circuit} (d.c.\ source of e.m.f.\ $\mathcal{E}_0$, total resistance $R$):
\begin{align*}
\text{Growth: } & I(t) = \frac{\mathcal{E}_0}{R}\left(1 - e^{-t/\tau}\right), \qquad \tau = \frac{L}{R},\\[2pt]
\text{Decay: } & I(t) = I_0\, e^{-t/\tau}, \qquad I_0 = \frac{\mathcal{E}_0}{R}.
\end{align*}
\item \textbf{Energy stored in an inductor}: $U = \dfrac{1}{2}LI^2$. During growth, at time $t$:
\[ U(t) = \frac{1}{2}L\,I(t)^2 = \frac{1}{2}L\left(\frac{\mathcal{E}_0}{R}\right)^2\left(1 - e^{-t/\tau}\right)^2. \]
\item \textbf{Mutual induction}: $\mathcal{E}_2 = -M\dfrac{dI_1}{dt}$ and $\mathcal{E}_1 = -M\dfrac{dI_2}{dt}$.
\end{itemize}

\courssub{Standard Methods}
\begin{enumerate}
\item \textbf{Flux calculation}: identify $\vec{B}$, the area $A$, and the angle $\theta$ between $\vec{B}$ and the normal to the surface. Then $\Phi = BA\cos\theta$; multiply by $N$ for flux linkage.
\item \textbf{Direction of the induced current (Lenz's law)}:
      \begin{enumerate}
      \item State the direction of the external field $\vec{B}$ through the circuit.
      \item Decide whether the flux is increasing or decreasing.
      \item The induced field $\vec{B}_{\text{ind}}$ opposes the \emph{change}: if the flux increases, $\vec{B}_{\text{ind}}$ is opposite to $\vec{B}$; if it decreases, $\vec{B}_{\text{ind}}$ is in the same direction as $\vec{B}$.
      \item Apply the right-hand grip rule to deduce the current direction from $\vec{B}_{\text{ind}}$.
      \end{enumerate}
\item \textbf{Polarity of a motional e.m.f.}: apply the magnetic force $\vec{F} = q(\vec{v}\times\vec{B})$ to the free electrons (negative charges); the end where electrons accumulate becomes the negative terminal.
\item \textbf{L--R circuit analysis}: write Kirchhoff's loop rule, $\mathcal{E}_0 - IR - L\dfrac{dI}{dt} = 0$; separate variables and integrate with the initial condition $I(0)=0$ (growth) or set $\mathcal{E}_0 = 0$ with $I(0)=I_0$ (decay).
\item \textbf{Energy balance in an L--R circuit}: power supplied by the source $P = \mathcal{E}_0 I$; power dissipated in the resistor $P_R = I^2R$; rate of storage in the inductor $P_L = LI\dfrac{dI}{dt}$. Energy conservation gives $\mathcal{E}_0 I = I^2R + LI\dfrac{dI}{dt}$, which is the loop equation multiplied by $I$.
\end{enumerate}

\courssub{Common Mistakes and Traps}
\begin{itemize}
\item \textbf{Sign errors}: forgetting the minus sign in Faraday's law when only the magnitude is required; confusing Lenz's law with opposing the \emph{field} itself instead of the \emph{change} of flux.
\item \textbf{The angle $\theta$}: using the angle between $\vec{B}$ and the \emph{plane} of the coil instead of the angle between $\vec{B}$ and the \emph{normal} to the plane in $\Phi = BA\cos\theta$.
\item \textbf{Motional e.m.f.}: applying $\mathcal{E} = Blv$ when $\vec{v}$, $\vec{B}$ and $\vec{l}$ are not mutually perpendicular; the correct form is then $\mathcal{E} = Blv\sin\theta$.
\item \textbf{Inductance formulas}: confusing the definition $N\Phi = LI$ with the induced e.m.f.\ $\mathcal{E} = -L\,dI/dt$. $L$ is a geometric property of the coil and does not depend on $I$.
\item \textbf{Time constant}: writing $\tau = R/L$ instead of $\tau = L/R$. Check physically: a larger $L$ opposes changes more strongly, so the current changes more slowly and $\tau$ must be larger.
\item \textbf{Energy}: believing the energy is stored in the resistor; it is stored in the \emph{magnetic field} of the inductor. The correct expression is $U = \tfrac{1}{2}LI^2$, not $LI^2$.
\item \textbf{Generators}: confusing slip rings (AC output) with the split-ring commutator (DC output). The DC output is pulsating (rectified), not steady; the peak e.m.f.\ $NBA\omega$ is the same in both cases.
\end{itemize}

\courssub{GCE Examination Tips}
\begin{itemize}
\item \textbf{Derivations}: be ready to derive $\mathcal{E} = Blv$ from the magnetic force on the free electrons of the moving rod, and to derive the growth and decay equations of the L--R circuit. Show the steps clearly: loop rule $\rightarrow$ differential equation $\rightarrow$ separation of variables $\rightarrow$ integration $\rightarrow$ application of the initial condition.
\item \textbf{Graph sketching}: be able to sketch $I$ against $t$ for growth and decay (exponential curves, with $\tau$ marked), $\mathcal{E}$ against $t$ for an AC generator (sinusoid, peak $\mathcal{E}_0 = NBA\omega$), and $\mathcal{E}$ against $t$ for a DC generator (rectified half-waves). Always label axes and peak values.
\item \textbf{Units}: always quote SI units: $\Phi$ in Wb, $L$ and $M$ in H, $\mathcal{E}$ in V, $\tau$ in s. Check dimensional consistency: $[L] = [\Phi]/[I] = \mathrm{Wb\,A^{-1}} = \mathrm{H}$, and $[L/R] = \mathrm{s}$.
\item \textbf{Quantitative problems}: typical values are $B \sim 0.1$--$1\ \mathrm{T}$, $A \sim 10^{-3}$--$10^{-2}\ \mathrm{m^2}$, $N \sim 100$--$500$ turns, $\omega = 2\pi f$ with $f = 50\ \mathrm{Hz}$, $L \sim \mathrm{mH}$--$\mathrm{H}$, $R \sim \Omega$--$\mathrm{k}\Omega$. Give final answers to 2--3 significant figures with units.
\item \textbf{Mutual inductance}: frequently examined with two coaxial coils; remember that the same constant $M$ appears in both $N_2\Phi_2 = MI_1$ and $N_1\Phi_1 = MI_2$, and that $\mathcal{E}_2 = -M\,dI_1/dt$.
\item \textbf{Language of definitions}: ``State Faraday's law'' $\rightarrow$ give the equation \emph{and} a sentence in words. ``Explain Lenz's law'' $\rightarrow$ mention opposition to the change of flux \emph{and} conservation of energy. ``Define self-inductance'' $\rightarrow$ write $N\Phi = LI$ and state ``the flux linkage per unit current in the coil''.
\end{itemize}

\findecours

\end{document}
